% Mechanically extracted after \mainmatter from the frozen standalone master.

\markboth{Exposé XVII. Unipotent Algebraic Groups}{Exposé XVII. Unipotent Algebraic Groups}
% BEGIN inputs/XVII/tex/components/00_expose_XVII_title_and_section0.tex
% Authority: Expo17.pdf, local p.1 / combined p.963 / printed pp.531--532.

\chapter*{EXPOSÉ XVII\\[0.4em]
UNIPOTENT ALGEBRAIC GROUPS. EXTENSIONS BETWEEN UNIPOTENT GROUPS AND GROUPS OF MULTIPLICATIVE TYPE}
\addcontentsline{toc}{chapter}{Exposé XVII. Unipotent Algebraic Groups. Extensions}

\begin{center}
\textit{by M. Raynaud}%
\begingroup
\renewcommand{\thefootnote}{*}%
\footnote{Cf. the note on page 1 of Exposé XV.}%
\endgroup
\end{center}

\section*{0.\ Some notation%
\begingroup
\renewcommand{\thefootnote}{0}%
\footnote{Version x.y of 18 November 2008.}%
\endgroup}
\addcontentsline{toc}{section}{0.\ Some notation}
\SGARefTarget{sga3:xvii:section:0}

In this chapter we shall be concerned chiefly with algebraic groups defined over a
field \(k\).  The number \(p\geq 0\) will always denote the characteristic of \(k\);
\(\mathbf F_p\) will denote the prime field with \(p\) elements when \(p>0\);
\(\bar{k}\) will denote an algebraically closed extension of \(k\); and \(q\) will
denote a prime number different from \(p\).

For every \(S\)-prescheme, \((\mathbb G_a)_S\) (respectively
\((\mathbb G_m)_S\)) denotes the additive group (respectively the multiplicative
group) over \(S\) (cf.\ \SGARefLink{sga3:i:subsection:4-3}{Exposé I, 4.3}).  For every integer \(n>0\),
\((\boldsymbol\mu_n)_S\) (respectively \((\mathbb Z/n\mathbb Z)_S\)) denotes the
group of \(n\)-th roots of unity (\SGARefLink{sga3:i:paragraph:4-4-4}{Exposé I, 4.4.4}) (respectively the constant group
over \(S\) associated with the abstract group \(\mathbb Z/n\mathbb Z\)
(\SGARefLink{sga3:i:subsection:4-1}{Exposé I, 4.1})).  The group \((\mathbb Z/n\mathbb Z)_S\) is finite and étale
over \(S\); the group \((\boldsymbol\mu_n)_S\) is flat and finite over \(S\), and
is étale over \(S\) if and only if \(n\) is invertible on \(S\)
(\SGARefLink{sga3:viii:proposition:2:1}{Exposé VIII, 2.1}).

Suppose that \(S\) is a prescheme of characteristic \(p>0\).  For every integer
\(n>0\) and every \(S\)-group prescheme \(G\), we denote by \(F^n(G)\) the
radicial \(S\)-subgroup prescheme of \(G\) equal to the kernel of the \(n\)-th
iterate of the relative Frobenius morphism of \(G\) (Exposé VII\textsubscript{A}).
In particular, if \(G=(\mathbb G_a)_S\), we set
\[
  F^n(G)=(\boldsymbol\alpha_{p^n})_S .
\]
This is a radicial \(S\)-group, flat and finite over \(S\), representing the
following functor: for every \(S\)-prescheme \(S'\),
\((\boldsymbol\alpha_{p^n})_S(S')\) is the set of
\(x'\in\Gamma(S',\mathcal O_{S'})\) such that \(x'^{p^n}=0\).

The already defined group \((\boldsymbol\mu_{p^n})_S\) is canonically isomorphic
to \(F^n((\mathbb G_m)_S)\).

When there is no ambiguity about the base scheme \(S\), we shall simply write
\(\mathbb G_a\), \(\mathbb G_m\), \(\boldsymbol\alpha_{p^n}\), and so on, in
place of \((\mathbb G_a)_S\), \((\mathbb G_m)_S\),
\((\boldsymbol\alpha_{p^n})_S\), and so on.

If \(G\) is a commutative \(S\)-group prescheme, then for every integer \(n>0\),
\({}_nG\) denotes the subgroup prescheme of \(G\) equal to the kernel of
raising to the \(n\)-th power in \(G\).

For the reader's convenience, we have collected in an appendix some properties
of algebraic groups proved in Exposés VI and VII, together with elementary
properties of Hochschild cohomology that will be useful in this chapter.

% END inputs/XVII/tex/components/00_expose_XVII_title_and_section0.tex

% BEGIN inputs/XVII/tex/components/01_expose_XVII_section1_opening.tex
% Authority: Expo17.pdf, local p.2 / combined p.964 / printed pp.533--534.

\SGASection{1}{Definition of unipotent algebraic groups}
\SGARefTarget{sga3:xvii:section:1}

\SGAStatement{Definition 1.1}{%
An algebraic group \(G\) defined over an algebraically closed field \(k\) is
called \emph{unipotent} if \(G\) admits a composition series whose successive
quotients are isomorphic to algebraic subgroups of \((\mathbb G_a)_k\).}
\SGARefTarget{sga3:xvii:definition:1:1}

\SGAStatement{Proposition 1.2}{%
Let \(k\) be an algebraically closed field, let \(K\) be an algebraically closed
extension of \(k\), and let \(G\) be an algebraic group defined over \(k\).
Then \(G\) is unipotent if and only if \(G_K\) is unipotent.}
\SGARefTarget{sga3:xvii:proposition:1:2}

The necessity is clear, since a composition series remains a composition series
after extension of the base.

The proof of sufficiency is standard.  The field \(K\) is the filtered direct
limit of its finitely generated \(k\)-subalgebras.  By Exposé
VI\textsubscript{B}, \S10, one can find a finitely generated \(k\)-subalgebra
\(A\) of \(K\), a composition series \(G_i\) of \(G_S\), where
\(S=\operatorname{Spec}A\), and immersions
\[
  u_i:H_i=G_i/G_{i+1}\longrightarrow(\mathbb G_a)_S .
\]
To prove that \(G\) is unipotent, it therefore suffices to make a \(k\)-valued
base change \(A\to k\).  This is possible because
\(\operatorname{Hom}_{k\text{-alg}}(A,k)\) is nonempty: \(A\) is a nonzero
finitely generated algebra over the algebraically closed field \(k\).

\SGAStatement{Definition 1.3}{%
Let \(G\) be an algebraic group defined over a field \(k\).  We say that \(G\)
is \emph{unipotent} if there exists an algebraically closed extension
\(\bar{k}\) of \(k\) such that \(G_{\bar{k}}\) is unipotent in the sense of
\SGARefLink{sga3:xvii:definition:1:1}{Definition 1.1}.}
\SGARefTarget{sga3:xvii:definition:1:3}

By \SGARefLink{sga3:xvii:proposition:1:2}{Proposition 1.2}, this property is independent of the chosen algebraically
closed extension \(\bar{k}\).

\SGAStatement{Definition 1.4}{%
Let \(G\) be an algebraic group defined over a field \(k\), and let \(H\) be an
algebraic group defined over an extension \(k'\) of \(k\).  We say that \(H\)
is a \emph{form} of \(G\) over \(k'\) if the algebraic groups
\(G_{\bar{k}'}\) and \(H_{\bar{k}'}\) become isomorphic over an algebraic
closure \(\bar{k}'\) of \(k'\).  As above, this property is independent of
the choice of \(\bar{k}'\).  We shall also say that \(H\) is a
\emph{twisted form} of \(G\).}
\SGARefTarget{sga3:xvii:definition:1:4}

We can now describe the algebraic subgroups of \(\mathbb G_a\).

\SGAStatement{Proposition 1.5}{%
Let \(k\) be a field of characteristic \(p\geq0\).  An algebraic subgroup
\(H\) of \((\mathbb G_a)_k\) is of one of the following types:
\begin{enumerate}
\item[\rm(i)] \(H=0\);
\item[\rm(ii)] \(H=\mathbb G_a\);
\item[\rm(iii)]\SGAInlineRefTarget{sga3:xvii:proposition:1:5:item:iii} if \(p>0\), \(H\) is an extension of a twisted constant group
  \((\mathbb Z/p\mathbb Z)^r\) by a radicial group
  \(\boldsymbol\alpha_{p^n}\), where \(r,n\geq0\) and \(r+n>0\).
  If, moreover, \(k\) is perfect, this extension is necessarily trivial.
\end{enumerate}}
\SGARefTarget{sga3:xvii:proposition:1:5}

% END inputs/XVII/tex/components/01_expose_XVII_section1_opening.tex

% BEGIN inputs/XVII/tex/components/02_expose_XVII_section1_remainder_and_section2_opening.tex
% Authority: Expo17.pdf, local pp.3--4 / combined pp.965--966 /
% printed pp.535--539.

\SGAProof{Proof}
If \(\dim H=1\), then clearly \(H=\mathbb G_a\).  Otherwise \(H\) has
dimension zero and is therefore an extension\footnote{Editor's note:
reference to Exposé VII\textsubscript{A}?} of an étale group \(H''\) by its
identity component \(H'\), which is radicial.  To describe \(H''_{\bar{k}}\),
it is enough to know the abstract group \(H''(\bar{k})\), which is isomorphic
to \(H(\bar{k})\).

The latter is a finite subgroup of \(\mathbb G_a(\bar{k})\).  It is therefore
zero when \(p=0\), and otherwise is of the form
\((\mathbb Z/p\mathbb Z)^r\), since it is annihilated by \(p\).  The group
\(H'\) is closed in \(\mathbb G_a\) and is defined by a single equation
(the coordinate ring \(\bar{k}[T]\) of \(\mathbb G_a\) is principal).
Since this equation has \(0\) as its only root over \(\bar{k}\), it is of the
form \(T^n=0\).  Compatibility with the group law implies that
\((T+T')^n\) belongs to the ideal generated by \(T^n\) and \(T'^n\) in
\(\bar{k}[T,T']\).  Thus, if \(p=0\), then \(n=1\) and \(H'\) is the trivial
group; while if \(p>0\), then \(n=p^m\) and
\(H'=\boldsymbol\alpha_{p^m}\).

The last assertion of \SGARefLink{sga3:xvii:proposition:1:5}{Proposition 1.5} follows more generally from the
following lemma.

\SGAStatement{Lemma 1.6}{%
If \(k\) is a perfect field, every extension \(H\) of an étale algebraic group
\(H''\) by a radicial group \(H'\) is trivial.  Moreover, \(H''\) has a unique
lifting in \(H\), namely \(H_{\mathrm{red}}\).}
\SGARefTarget{sga3:xvii:lemma:1:6}

Indeed, since \(k\) is perfect, the reduced \(k\)-scheme
\(H_{\mathrm{red}}\) is an algebraic subgroup of \(H\)
(Exposé VI\textsubscript{A}, \SGARefLink{sga3:VIA:0.2}{0.2}), is geometrically reduced and hence smooth
over \(k\) (Exposé VI\textsubscript{A}, \SGARefLink{prop:sga3-VIA-1.3.1}{1.3.1}), and is therefore étale because
\(H\) has dimension zero.  To see that the canonical projection
\(H_{\mathrm{red}}\to H''\) is an isomorphism, it is enough to check this after
the base-field extension \(k\to\bar{k}\); there it is enough to check that it
induces an isomorphism on \(\bar{k}\)-valued points, which is clear.  The final
assertion follows because every lifting of \(H''\) in \(H\), being étale over
\(k\), is reduced and hence necessarily contained in \(H_{\mathrm{red}}\).

Notice that \(\boldsymbol\alpha_{p^n}\) is an iterated extension of groups
isomorphic to \(\boldsymbol\alpha_p\).  \SGARefLink{sga3:xvii:proposition:1:5}{Proposition 1.5} therefore gives the
following corollary.

\SGAStatement{Corollary 1.7}{%
Let \(G\) be an algebraic group defined over an algebraically closed field
\(k\).  Then \(G\) is unipotent if and only if it has a composition series
whose successive quotients are isomorphic to \(\mathbb G_a\) when \(p=0\),
and to one of the groups
\(\mathbb G_a,\mathbb Z/p\mathbb Z,\boldsymbol\alpha_p\) when \(p>0\).
We shall call these the \emph{elementary unipotent groups}.}
\SGARefTarget{sga3:xvii:corollary:1:7}

\SGASection{2}{First properties of unipotent groups}
\SGARefTarget{sga3:xvii:section:2}

\SGAStatement{Proposition 2.1}{%
A unipotent algebraic group defined over a field \(k\) is affine over \(k\).}
\SGARefTarget{sga3:xvii:proposition:2:1}

By fpqc descent for affine morphisms, it suffices to prove the assertion when
\(k\) is algebraically closed.  In that case, \(G\) is by definition an
iterated extension of affine algebraic groups, and is therefore affine by
Exposé VI\textsubscript{B}, \SGARefLink{concl:sga3-VIB-9.2-viii}{9.2 (viii)}, applied to affine
morphisms.\footnote{Editor's note: verify this reference.}

\SGAStatement{Proposition 2.2}{%
\begin{enumerate}
\item[\rm\SGAInlineRefTarget{sga3:xvii:proposition:2:1:item:i}(i)] The property that an algebraic group is unipotent is invariant
  under extension of the base field.
\item[\rm\SGAInlineRefTarget{sga3:xvii:proposition:2:1:item:ii}(ii)] Every algebraic subgroup of a unipotent group is unipotent.
\item[\rm\SGAInlineRefTarget{sga3:xvii:proposition:2:1:item:iii}(iii)] Every algebraic quotient group of a unipotent group is
  unipotent.
\item[\rm\SGAInlineRefTarget{sga3:xvii:proposition:2:1:item:iv}(iv)] Every extension of a unipotent algebraic group by a unipotent
  algebraic group is likewise unipotent.
\end{enumerate}}
\SGARefTarget{sga3:xvii:proposition:2:2}

\SGAProof{Proof}
Assertion \SGARefLink{sga3:xvii:proposition:2:1:item:i}{(i)} follows immediately from \SGARefLink{sga3:xvii:definition:1:3}{Definitions 1.3} and \SGARefLink{sga3:xvii:proposition:1:2}{Proposition 1.2}.
To prove the remaining assertions, we may suppose that \(k\) is algebraically
closed.  Assertion \SGARefLink{sga3:xvii:proposition:2:1:item:iv}{(iv)} is then immediate from the definition.

Let \(G\) be a unipotent algebraic group, \(G'\) an algebraic subgroup of
\(G\), \(G''\) an algebraic quotient group of \(G\), and let
\(G_i\), \(1\leq i\leq n\), be a composition series of \(G\) such that
\[
  H_i=G_i/G_{i+1}
\]
is an elementary unipotent group (\SGARefLink{sga3:xvii:corollary:1:7}{Corollary 1.7}).

For \SGARefLink{sga3:xvii:proposition:2:1:item:ii}{(ii)}, consider the composition series of \(G'\) induced by that of \(G\):
\(G'_i=G_i\cap G'\).  The group \(G'_i/G'_{i+1}\) identifies with an
algebraic subgroup of \(H_i\), hence is isomorphic to a subgroup of
\(\mathbb G_a\).  It follows that \(G'\) is unipotent.

For \SGARefLink{sga3:xvii:proposition:2:1:item:iii}{(iii)}, consider the composition series of \(G''\) obtained as the image
of that of \(G\): let \(G''_i\) be the image of \(G_i\) in \(G''\).  Then
\(G''_i/G''_{i+1}\) is a quotient of \(H_i\), so it remains only to prove the
following lemma.

\SGAStatement{Lemma 2.3}{%
Let \(H\) be an elementary unipotent group (\SGARefLink{sga3:xvii:corollary:1:7}{Corollary 1.7}) defined over a
field \(k\).  Every algebraic quotient group \(H''\) of \(H\) is either zero
or isomorphic to \(H\) over \(k\).}
\SGARefTarget{sga3:xvii:lemma:2:3}

\SGAProof{Proof}
If \(H=\mathbb G_a\) in characteristic zero, or if
\(H=\boldsymbol\alpha_p\) or \(\mathbb Z/p\mathbb Z\) in characteristic
\(p>0\), it is enough to observe from \SGARefLink{sga3:xvii:proposition:1:5}{Proposition 1.5} that \(H\) has no
algebraic subgroups other than \(0\) and \(H\).  Indeed, a nonzero algebraic
subgroup of \(\boldsymbol\alpha_p\) is defined by a \(k\)-algebra of rank at
least \(p\) over \(k\), by \SGARefLink{sga3:xvii:proposition:1:5:item:iii}{Proposition 1.5(iii)}, and is therefore
\(\boldsymbol\alpha_p\).

Now let \(H=\mathbb G_a\) in characteristic \(p>0\), so that there is an exact
sequence
\[
  0\longrightarrow N\longrightarrow\mathbb G_a
    \longrightarrow H''\longrightarrow0.
\]
If \(N=\mathbb G_a\), then \(H''=0\).  Otherwise, arguing by induction on the
length of a composition series of \(N\), we may suppose that
\(N=\boldsymbol\alpha_p\), or that \(N\) is a form of
\((\mathbb Z/p\mathbb Z)^r\).

\begin{enumerate}
\item[\rm(a)] If \(N\simeq\boldsymbol\alpha_p\), the proof of
  Proposition 1.5(iii) shows that \(N\) is necessarily the kernel of the
  Frobenius morphism \(F\) of \(\mathbb G_a\), and the assertion follows from
  the exact sequence
  \[
    0\longrightarrow\boldsymbol\alpha_p\longrightarrow\mathbb G_a
      \xrightarrow{F}\mathbb G_a\longrightarrow0.
  \]
\item[\rm(b)] If \(N\simeq\mathbb Z/p\mathbb Z\), there is plainly an
  \(a\in k^\times\) such that \(N\) is the closed subscheme of
  \(\mathbb G_a=\operatorname{Spec}k[X]\) defined by
  \(X^p-aX=0\).  The assertion then follows from the exact sequence
  \[
    0\longrightarrow N\longrightarrow\mathbb G_a
      \xrightarrow{P}\mathbb G_a\longrightarrow0,
  \]
  where \(P\) is the Artin--Schreier morphism \(x\mapsto x^p-ax\).
\item[\rm(c)] If \(N\) is a form of \((\mathbb Z/p\mathbb Z)^r\), there is a
  finite Galois extension \(k'/k\) that trivializes \(N\).  By \SGARefLink{sga3:xvii:section:b}{(b)} and an
  evident induction on \(r\), the quotient \(\mathbb G_a/N\) is a form of
  \(\mathbb G_a\) trivialized by \(k'\).  It remains to apply the next lemma.
\end{enumerate}

\SGAStatement{Lemma 2.3 bis}{%
Let \(k\) be a field, and let \(G\) be a \(k\)-algebraic group that is a form
of \(\mathbb G_a\), trivialized by a finite separable extension \(k'/k\).
Then \(G\) is isomorphic to \(\mathbb G_a\).}
\SGARefTarget{sga3:xvii:lemma:2:3:bis}

% END inputs/XVII/tex/components/02_expose_XVII_section1_remainder_and_section2_opening.tex

% BEGIN inputs/XVII/tex/components/03_expose_XVII_section2_continued_page005.tex
% Authority: Expo17.pdf, local p.5 / combined p.967 / printed pp.539--541.

\SGAProof{Proof of \SGARefLink{sga3:xvii:lemma:2:3:bis}{Lemma 2.3 bis}}
Indeed, the group of \(k'\)-automorphisms of the algebraic group
\((\mathbb G_a)_{k'}\) is the group of nonzero homotheties
\((k')^\times\) (\textit{Bible}, Exposé 9, \SGARefLink{sga3:xvii:lemma:3:1}{Lemma 1}), and the Galois
cohomology group \(H^1(k'/k,\mathbb G_m)\) is zero (we may suppose that
\(k'\) is a Galois extension of \(k\)), by Hilbert's Theorem 90.
\SGARefLink{sga3:xvii:lemma:2:3:bis}{Lemma 2.3 bis} then follows from the classification of the \(k\)-forms
of an algebraic group (cf.\ J.-P. Serre, \textit{Cohomologie
galoisienne}, Chapter III, \SGARefLink{sga3:iii:subsection:1-3}{1.3}).

This completes the proof of \SGARefLink{sga3:xvii:proposition:2:2}{2.2}.

\SGAStatement{Proposition 2.4}{Let \(k\) be a field, let \(M\) be a
\(k\)-group of multiplicative type (Exposé VIII) and of finite type,
and let \(U\) be a unipotent \(k\)-algebraic group.  Then:
\begin{enumerate}
\item[\SGAInlineRefTarget{sga3:xvii:proposition:2-4:item:i}(i)]\SGAInlineRefTarget{sga3:xvii:proposition:2:4:item:i} \(\underline{\operatorname{Hom}}_{k\text{-gr}}(M,U)
  =\underline e\)
  (a fortiori \(\operatorname{Hom}_{k\text{-gr}}(M,U)=e\));
\item[(ii)] \(\operatorname{Hom}_{k\text{-gr}}(U,M)=e\).
\end{enumerate}}
\SGARefTarget{sga3:xvii:proposition:2:4}

To prove (i), we must establish that, for every prescheme \(S\) over
\(k\),
\[
  \operatorname{Hom}_{S\text{-gr}}(M_S,U_S)=e.
\]
This follows from the next lemma.

\SGAStatement{Lemma 2.5}{Let \(S\) be a prescheme, let \(M\) be an
\(S\)-group of multiplicative type and of finite type over \(S\), and
let \(U\) be an \(S\)-group prescheme, of finite presentation over
\(S\), whose fibers are unipotent.  Then
\(\operatorname{Hom}_{S\text{-gr}}(M,U)=e\).}
\SGARefTarget{sga3:xvii:lemma:2:5}

\SGAProof{Proof}
Under these hypotheses, for an \(S\)-group morphism \(u:M\to U\) to
be the unit homomorphism, it suffices that the restriction of \(u\) to
the fibers of \(M\) over the points of \(S\) be the zero morphism
(\SGARefLink{sga3:ix:corollary:5:2}{Exposé IX, 5.2}).  We are therefore reduced to the case in which \(S\)
is the spectrum of a field, which we may moreover suppose to be
algebraically closed.  In view of \SGARefLink{sga3:xvii:definition:1:1}{Definition 1.1}, it is enough to take
\(U=\mathbb G_a\), in which case the assertion has already been proved
(\SGARefLink{sga3:xii:4.4.1}{Exposé XII, 4.4.1}).

We now prove Proposition 2.4(ii).  Let \(u:U\to M\) be a \(k\)-group
morphism.  The image \(u(U)\) is represented by an algebraic subgroup
\(U''\) of \(M\) (Exposé VI\textsubscript{B}, \SGARefLink{cor:sga3-VIB-5.4}{5.4}).\footnote{Editor's
note: verify this reference.}  The group \(U''\) is unipotent as a
quotient of a unipotent group (2.2(iii)), and is of multiplicative type
as a subgroup of a group of multiplicative type (cf.\ \textit{Bible},
Exposé 4, Theorem 2, \SGARefLink{sga3:xvii:corollary:C:5:1}{Corollary 1}, or \SGARefLink{sga3:ix:theorem:6:8}{Exposé IX, 6.8}).  Thus \(U''\) is
the unit group by \SGARefLink{sga3:xvii:proposition:2-4:item:i}{Proposition 2.4(i)}.

\SGAStatement{\SGAInlineRefTarget{sga3:xvii:remark:2-6}Remark 2.6}{Keeping the notation of \SGARefLink{sga3:xvii:proposition:2:4}{Proposition 2.4}, it
is no longer true in general that the functor
\(\underline{\operatorname{Hom}}_{k\text{-gr}}(U,M)\) is equal to
\(\underline e\).  Indeed, let \(S\) be a prescheme such that
\(\Gamma(S,\mathcal O_S)\) contains a nonzero element \(\varepsilon\)
with \(\varepsilon^2=0\) (for example, the spectrum of the algebra of
dual numbers over a ring \(A\)).  For every \(S'\) over \(S\), the map
\[
  u\longmapsto 1+\varepsilon_{S'}u
\]
defines, functorially in \(S'\), a homomorphism from the additive group
\(\Gamma(S',\mathcal O_{S'})\) to the multiplicative group
\(\Gamma(S',\mathcal O_{S'}^\times)\).  It therefore defines an
\(S\)-group morphism
\[
  (\mathbb G_a)_S\longrightarrow(\mathbb G_m)_S,
\]
and, since \(\varepsilon\ne0\), this morphism is not zero.}
\SGARefTarget{sga3:xvii:remark:2:6}

Recall (Exposé VII\textsubscript{A}, \(\S3\)) that when \(G\) is a
commutative \(S\)-group prescheme, finite and flat over \(S\), Cartier
duality is available and \(G\) is reflexive in the sense of Exposé
VIII, \(\S1\).  More precisely, the functor
\[
  G\longmapsto
  \underline{\operatorname{Hom}}_{S\text{-gr}}(G,\mathbb G_m)
\]
is represented

% END inputs/XVII/tex/components/03_expose_XVII_section2_continued_page005.tex

% BEGIN inputs/XVII/tex/components/04_expose_XVII_cartier_duality_and_section3_opening_page006.tex
% Authority: Expo17.pdf, local p.6 / combined p.968 / printed pp.541--543.

by a commutative \(S\)-group prescheme \(D(G)\), finite and flat over
\(S\), and the canonical morphism
\[
  G\longrightarrow D(D(G))
\]
is an isomorphism.  In particular:
\begin{enumerate}
\item[(i)] \(D(\boldsymbol\mu_n)\simeq\mathbb Z/n\mathbb Z\), and
  consequently \(D(\mathbb Z/n\mathbb Z)\simeq\boldsymbol\mu_n\);
\item[(ii)] if \(S\) has characteristic \(p>0\), then
  \(D(\boldsymbol\alpha_p)\simeq\boldsymbol\alpha_p\).
\end{enumerate}

\SGASection{3}{Unipotent groups acting on a vector space}
\SGARefTarget{sga3:xvii:section:3}

Recall (\SGARefLink{sga3:i:definition:4-6-1}{Exposé I, 4.6.1}) that if \(S\) is a prescheme and \(M\) is a
sheaf of \(\mathcal O_S\)-modules, \(\mathbf W(M)\) denotes the
\(S\)-functor
\[
  (\operatorname{Sch}/S)^\circ\longrightarrow(\operatorname{Ens})
\]
defined, for every \(S\)-prescheme \(S'\), by
\[
  \mathbf W(M)(S')
  =\Gamma\!\left(S',M\otimes_{\mathcal O_S}\mathcal O_{S'}\right).
\]

Recall also that if an \(S\)-group \(G\) acts on an \(S\)-functor
\(V\), the \(S\)-functor \(V^G\) of invariants of \(V\) under \(G\) is
the subfunctor of \(V\) whose \(S'\)-valued points, for \(S'\) over
\(S\), are those \(x\in V(S')\) such that \(x_{S''}\) is fixed under
\(G(S'')\) for every \(S''\) over \(S'\).

We then have the following lemma.

\SGAStatement{\SGAInlineRefTarget{sga3:xvii:lemma:3-1}Lemma 3.1}{Let \(S\) be a prescheme, and let \(G\) be an
\(S\)-group prescheme, affine over \(S\), defined by the
quasi-coherent \(\mathcal O_S\)-algebra \(A\).  Suppose that \(G\) acts
on a quasi-coherent sheaf \(M\) of \(\mathcal O_S\)-modules.  Let
\[
  \mu:M\longrightarrow A\otimes_{\mathcal O_S}M
\]
be the comorphism defining the action of \(G\) on \(M\) (\SGARefLink{sga3:i:proposition:4-7-2}{Exposé I,
4.7.2}), and let
\[
  \nu:M\longrightarrow A\otimes_{\mathcal O_S}M
\]
be the morphism \(x\mapsto\mu(x)-1\otimes x\).  Then:
\begin{enumerate}
\item[\SGAInlineRefTarget{sga3:xvii:lemma:3-1:item:i}(i)] \(\underline M^{\,G}(S)=\Gamma\operatorname{Ker}(\nu)\);
\item[\SGAInlineRefTarget{sga3:xvii:lemma:3-1:item:ii}(ii)] if \(S\) is the spectrum of a field \(k\), then
  \(\underline M^{\,G}\) is of the form \(\mathbf W(N)\), where \(N\)
  is a vector subspace of \(M\);
\item[\SGAInlineRefTarget{sga3:xvii:lemma:3-1:item:iii}(iii)] if \(S\) is the spectrum of a field \(k\), every element
  \(x\) of \(M\) is contained in a finite-dimensional \(k\)-vector
  subspace of \(M\) stable under the action of \(G\).
\end{enumerate}}
\SGARefTarget{sga3:xvii:lemma:3:1}

\SGAProof{Proof}
Assertion \SGARefLink{sga3:xvii:lemma:3-1:item:iii}{(iii)} is recalled here for reference and was already proved
in Exposé VI\textsubscript{B}, \SGARefLink{cor:sga3-VIB-11.2}{11.2}.

For \SGARefLink{sga3:xvii:lemma:3-1:item:i}{(i)}, it is clear that \(\underline M^{\,G}(S)\) contains
\(\Gamma\operatorname{Ker}(\nu)\).  To prove the reverse inclusion, we
may suppose that \(S\) is affine, with ring \(B\).  Let
\(m\in\underline M^{\,G}(S)\).  For every \(B\)-algebra \(B'\) and
every
\[
  u\in\operatorname{Hom}_{B\text{-alg}}(A,B')
\]
(where \(u\) corresponds to an element of
\(G(\operatorname{Spec}B')\)), we have
\[
  (u\otimes1_M)\nu(m)=0
  \qquad\text{in }B'\otimes_BM.
\]
Taking \(B'=A\) and \(u\) to be the identity of \(A\), we obtain
\(\nu(m)=0\).

For (ii), let \(N\) be the kernel of \(\nu\), equal to
\(\underline M^{\,G}(k)\) by \SGARefLink{sga3:xvii:lemma:3-1:item:i}{(i)}.  Every \(k\)-prescheme \(S\) is flat
over \(k\); hence
\[
  \underline M^{\,G}(S)
  =\Gamma\operatorname{Ker}(\nu\otimes_k S)
  =\Gamma(N\otimes_k S).
\]
Thus \(\underline M^{\,G}\) is isomorphic to the functor
\(\mathbf W(N)\).

\SGAStatement{Proposition 3.2}{Let \(G\) be a unipotent algebraic
group, defined over a field \(k\), acting on a \(k\)-vector space
\(V\).  If \(V\ne0\), then \(V^G\ne0\).}
\SGARefTarget{sga3:xvii:proposition:3:2}

% END inputs/XVII/tex/components/04_expose_XVII_cartier_duality_and_section3_opening_page006.tex

% BEGIN inputs/XVII/tex/components/05_expose_XVII_proposition3_2_proof_and_lemma3_3_page007.tex
% Authority: Expo17.pdf, local p.7 / combined p.969 / printed pp.544--545.

In view of \SGARefLink{sga3:xvii:lemma:3-1:item:ii}{3.1(ii)} and (iii), we may suppose that \(k\) is
algebraically closed and that \(V\) is finite-dimensional over \(k\).

Let
\[
  0\longrightarrow G'\longrightarrow G\longrightarrow G''
    \longrightarrow0
\]
be an exact sequence of algebraic groups, and suppose that \(G\) acts
on a sheaf \(V\) (for the fpqc topology).  Of course
\(V^G\subset V^{G'}\), and the quotient presheaf \(G/G'\) acts
naturally on \(V^{G'}\).  But \(V^{G'}\) is a sheaf, being the kernel
of the familiar pair of morphisms
\[
  V\rightrightarrows\underline{\operatorname{Hom}}(G',V).
\]
Consequently the sheaf associated with \(G/G'\), namely \(G''\),
acts on \(V^{G'}\), and it is immediate that
\[
  \left(V^{G'}\right)^{G''}
  =V^G
  =\left(V^{G'}\right)^{G/G'}.
\]

This remark reduces the proof of \SGARefLink{sga3:xvii:proposition:3:2}{3.2} to the case in which \(G\) is an
elementary unipotent group (\SGARefLink{sga3:xvii:corollary:1:7}{1.7}).

\SGARefTarget{sga3:xvii:proposition:3:2:case:a}
\textit{\SGARefLink{sga3:xvii:appendix:A}{(a)} \(G=\mathbb G_a\), \(p=0\).}
It follows from (BIBLE 4, Proposition 4) that a morphism from
\(\mathbb G_a\) to the linear group \(\operatorname{GL}(V)\) is given
by an exponential map
\[
  T\longmapsto\sum_{q=0}^{\infty}T^q\,\frac{n^q}{q!},
\]
where \(n\) is a nilpotent endomorphism of \(V\).  But then
\(V\ne0\) implies \(\operatorname{Ker}n\ne0\), and it is clear that
every vector of \(V\) annihilated by \(n\) is fixed by \(G\).

Suppose now that \(p>0\).

\SGARefTarget{sga3:xvii:proposition:3:2:case:b}
\textit{\SGARefLink{sga3:xvii:section:b}{(b)} \(G=\boldsymbol\alpha_p\).}
Since \(\boldsymbol\alpha_p\) is a radicial group of height \(1\)
(Exposé VII\textsubscript{A}, \S7), giving a representation of
\(\boldsymbol\alpha_p\) in \(V\) amounts to giving a representation
of the \(p\)-Lie algebra \(\boldsymbol\alpha_p\) in
\(\mathfrak{gl}(V)\) (Appendix II, \SGARefLink{sga3:ii:subsection:2-2}{2.2}), that is, in this case, an
element \(X\in\operatorname{End}(V)\) such that \(X^p=0\)
(Appendix II, \SGARefLink{sga3:xvii:proposition:2:1}{2.1}).  But then
\[
  V\ne0\quad\Longrightarrow\quad
  W=\operatorname{Ker}(X)\ne0,
\]
and, again by Appendix II, \SGARefLink{sga3:ii:subsection:2-2}{2.2}, one has
\(W=V^{\boldsymbol\alpha_p}\).

\SGARefTarget{sga3:xvii:proposition:3:2:case:c}
\textit{\SGARefLink{sga3:xvii:appendix:C}{(c)} \(G=\mathbb Z/p\mathbb Z\).}
A representation of \(G\) in \(V\) is equivalent to the choice of an
element \(x\in\operatorname{Aut}(V)\) such that \(x^p=1\), that is,
\((1-x)^p=0\).  Thus \(x\) is of the form \(1+n\), with \(n\)
nilpotent, and \(W=\operatorname{Ker}n\) is fixed by \(x\).

\SGARefTarget{sga3:xvii:proposition:3:2:case:d}
\textit{(d) \(G=\mathbb G_a\).}
Let \(G_i\), \(i\in I\), be the increasing filtered family of étale
algebraic subgroups of \(\mathbb G_a\), hence groups isomorphic to
\((\mathbb Z/p\mathbb Z)^{r_i}\) (\SGARefLink{sga3:xvii:proposition:1:5}{Proposition 1.5}), and put
\(V_i=V^{G_i}\).  Since \(V\) has finite nonzero dimension and each
\(V_i\) is nonzero by \SGARefLink{sga3:xvii:appendix:C}{(c)}, the decreasing filtered family of the
\(V_i\) is stationary, and
\[
  W=\bigcap_{i\in I}V_i\ne0.
\]
We shall use the following lemma.

\SGAStatement{Lemma 3.3}{The family of étale subgroups of
\((\mathbb G_a)_S\), where \(S\) is a prescheme of characteristic
\(p>0\), is schematically dense in \(G\) (\SGARefLink{sga3:ix:definition:4:1}{Exposé IX, 4.1}).}
\SGARefTarget{sga3:xvii:lemma:3:3}

By \SGARefLink{sga3:ix:lemma:4:4}{Exposé IX, 4.4}, it suffices to prove the lemma when \(S\) is the
spectrum of the prime field \(\mathbb F_p\).  In that case it is
enough to consider the family of étale subgroups \(G_n\), \(n\geq1\),
defined by
\[
  X^{p^n}-X=0.
\]
This family is schematically dense in
\((\mathbb G_a)_{\mathbb F_p}\), since it contains every closed point.

With this established, let us return to the proof of \SGARefLink{sga3:xvii:proposition:3:2:case:d}{3.2(d)}.  If
\(w\in W\), its stabilizer in \(G\) is an algebraic subgroup of
\(\mathbb G_a\) containing \(G_i\) for every \(i\in I\); it is
therefore equal to \(\mathbb G_a\) by \SGARefLink{sga3:xvii:lemma:3:3}{3.3}.  Consequently
\(W=V^G\).

\SGARefLink{sga3:xvii:proposition:3:2}{Proposition 3.2} immediately gives the following.

% END inputs/XVII/tex/components/05_expose_XVII_proposition3_2_proof_and_lemma3_3_page007.tex

% BEGIN inputs/XVII/tex/components/06_expose_XVII_corollary3_4_theorem3_5_definition3_6_page008.tex
% Authority: Expo17.pdf, local p.8 / combined p.970 / printed pp.546--547.

\SGAStatement{Corollary 3.4}{Let \(k\) be a field, and let \(G\) be a
unipotent \(k\)-algebraic group acting on a finite-dimensional
\(k\)-vector space \(V\).  Then \(V\) has a sequence of vector
subspaces \(V_i\), defined over \(k\) and stable under \(G\),
\[
  0=V_0\subset V_1\subset\cdots\subset V_n=V,
\]
such that \(G\) acts trivially on \(V_{i+1}/V_i\).  Moreover, the
\(V_i\) may be chosen so that \(V_{i+1}/V_i\) has dimension \(1\).}
\SGARefTarget{sga3:xvii:corollary:3:4}

We shall now summarize and supplement the properties of unipotent
groups already proved in the following theorem.

\SGAStatement{\SGAInlineRefTarget{sga3:xvii:theorem:3-5}Theorem 3.5}{Let \(G\) be an algebraic group defined
over a field \(k\).  The following properties are equivalent:
\begin{enumerate}
\item[(i)] \SGARefTarget{sga3:xvii:theorem:3:5:item:i}
  \(G\) is unipotent.
\item[(ii)] \SGARefTarget{sga3:xvii:theorem:3:5:item:ii}
  \(G\) has a composition series defined over \(k\), whose successive
  quotients are isomorphic to \(\mathbb G_a\) if \(p=0\), and,
  respectively, to \(\boldsymbol\alpha_p\), \(\mathbb G_a\), or a
  twisted form of \((\mathbb Z/p\mathbb Z)^r\) (\SGARefLink{sga3:xvii:definition:1:4}{1.4}) if \(p>0\).
\item[(iii)] \SGARefTarget{sga3:xvii:theorem:3:5:item:iii}
  As in (ii), with the additional requirement that the composition
  series be central.
\item[(iv)] \SGARefTarget{sga3:xvii:theorem:3:5:item:iv}
  \(G\) has a characteristic composition series
  (Exposé VI\textsubscript{B}) defined over \(k\), whose successive
  quotients are isomorphic to \((\mathbb G_a)^r\) if \(p=0\), and,
  respectively, to \((\boldsymbol\alpha_p)^r\), a twisted form of
  \((\mathbb G_a)^s\), or a twisted form of
  \((\mathbb Z/p\mathbb Z)^t\), taken in this order, if \(p>0\).
\item[(v)] \SGARefTarget{sga3:xvii:theorem:3:5:item:v}
  \(G\) is isomorphic to an algebraic subgroup of the group
  \(\operatorname{Trigstr}(n)_k\) of strictly upper-triangular
  matrices in the linear group \(\operatorname{GL}(n)_k\), for a
  suitable integer \(n\geq0\).
\item[(vi)] \SGARefTarget{sga3:xvii:theorem:3:5:item:vi}
  \(G\) is affine, and for every linear representation of \(G\) on a
  nonzero finite-dimensional \(k\)-vector space \(V\), one has
  \(V^G\ne0\).
\end{enumerate}}
\SGARefTarget{sga3:xvii:theorem:3:5}

\SGAProof{Proof}
The implication \SGARefLink{sga3:xvii:theorem:3:5:item:i}{(i)}\(\Rightarrow\)\SGARefLink{sga3:xvii:theorem:3:5:item:vi}{(vi)} follows from \SGARefLink{sga3:xvii:proposition:2:1}{2.1} and \SGARefLink{sga3:xvii:proposition:3:2}{3.2}.

For \SGARefLink{sga3:xvii:theorem:3:5:item:vi}{(vi)}\(\Rightarrow\)\SGARefLink{sga3:xvii:theorem:3:5:item:v}{(v)}, since the algebraic group \(G\) is
affine, it is an algebraic subgroup of a suitable linear group
\(\operatorname{GL}(V)\) (Exposé VI\textsubscript{B}, \SGARefLink{def:sga3-VIB-11.3}{11.3}).
Applying \SGARefLink{sga3:xvii:corollary:3:4}{3.4} to the representation of \(G\) on \(V\) defined by this
embedding gives \SGARefLink{sga3:xvii:theorem:3:5:item:v}{(v)}.

For \SGARefLink{sga3:xvii:theorem:3:5:item:v}{(v)}\(\Rightarrow\)\SGARefLink{sga3:xvii:theorem:3:5:item:iii}{(iii)}, the algebraic group
\(\operatorname{Trigstr}(n)\) is known to have a central composition
series whose successive quotients are isomorphic to \(\mathbb G_a\).
The composition series induced on \(G\) gives property \SGARefLink{sga3:xvii:theorem:3:5:item:iii}{(iii)}, in view
of \SGARefLink{sga3:xvii:proposition:1:5}{1.5}.

The implications
\[
  \text{(iii)}\Rightarrow\text{(ii)}\Rightarrow\text{(i)}
  \qquad\text{and}\qquad
  \text{(iv)}\Rightarrow\text{(i)}
\]
are clear.

We shall prove \SGARefLink{sga3:xvii:theorem:3:5:item:i}{(i)}\(\Rightarrow\)\SGARefLink{sga3:xvii:theorem:3:5:item:iv}{(iv)} shortly, but first note some
consequences of what has already been proved.

\SGAStatement{Definition 3.6}{A \(p\)-Lie algebra \(\mathfrak g\),
\(p>0\) (cf. Exposé VII\textsubscript{A}, \S5), is called
\emph{unipotent} if the map
\[
  x\longmapsto x^{(p)}
\]
is nilpotent; that is, if for every \(x\in\mathfrak g\) there is an
integer \(n>0\) such that \(x^{(p^n)}=0\).}
\SGARefTarget{sga3:xvii:definition:3:6}

% END inputs/XVII/tex/components/06_expose_XVII_corollary3_4_theorem3_5_definition3_6_page008.tex

% BEGIN inputs/XVII/tex/components/07_expose_XVII_corollaries3_7_3_8_and_lemma3_9_page009.tex
% Authority: Expo17.pdf, local pp.8--9 / combined pp.970--971 /
% printed pp.547--549.

\SGAStatement{Corollary 3.7}{A unipotent algebraic group \(G\) is
nilpotent (Exposé VI\textsubscript{B}, \S8); its Lie algebra
\(\mathfrak g\) is nilpotent (Bourbaki, \emph{Groupes et algèbres de
Lie}, Chap.~I, \S4) and is isomorphic to a Lie algebra of nilpotent
endomorphisms of a finite-dimensional vector space.  In
characteristic \(p>0\), \(\mathfrak g\) is a unipotent \(p\)-Lie
algebra (\SGARefLink{sga3:xvii:definition:3:6}{3.6}).}
\SGARefTarget{sga3:xvii:corollary:3:7}

Since \SGARefLink{sga3:xvii:theorem:3:5:item:i}{(i)} implies (v), it is enough to prove \SGARefLink{sga3:xvii:corollary:3:7}{3.7} when
\(G=\operatorname{Trigstr}(n)\).  We have already used the fact that
\(\operatorname{Trigstr}(n)\) is a nilpotent algebraic group.
Moreover, the Lie algebra \(\mathfrak h\) of
\(\operatorname{Trigstr}(n)\) consists of the upper-triangular
endomorphisms of \(V\) that have zeros on the main diagonal.  These
endomorphisms are therefore nilpotent, and consequently
\(\mathfrak h\) is nilpotent (Bourbaki, \emph{loc.\ cit.},
Chap.~I, \S4, Cor.~3).  If \(p>0\), the \(p\)-th power in the
\(p\)-Lie algebra
\(\mathfrak{gl}(V)=\operatorname{End}(V)\) coincides with the
\(p\)-th power of endomorphisms of \(V\)
(Exposé VII\textsubscript{A}, \SGARefLink{sga3:VIIA:6.4.4}{6.4.4}); hence \(\mathfrak h\) is
unipotent.

\SGAStatement{Corollary 3.8}{Let \(k\) be an algebraically closed
field, and let \(G\) be a smooth affine algebraic group over \(k\).
Then the following properties are equivalent:
\begin{enumerate}
\item[(i)] \SGARefTarget{sga3:xvii:corollary:3:8:item:i}
  \(G\) is unipotent.
\item[(ii)] \SGARefTarget{sga3:xvii:corollary:3:8:item:ii}
  \(G(k)\) consists of unipotent elements
  (BIBLE~4, Prop.~4, Cor.~1); that is, \(G\) is unipotent in the
  sense of BIBLE.
\end{enumerate}}
\SGARefTarget{sga3:xvii:corollary:3:8}

The implication \SGARefLink{sga3:xvii:corollary:3:8:item:i}{(i)}\(\Rightarrow\)\SGARefLink{sga3:xvii:corollary:3:8:item:ii}{(ii)} follows because, by
\SGARefLink{sga3:xvii:proposition:3:2}{3.2}, \SGARefLink{sga3:xvii:corollary:3:8:item:i}{(i)}\(\Rightarrow\)(v), \(G\) is isomorphic to an algebraic
subgroup of a group \(\operatorname{Trigstr}(n)\).

For \SGARefLink{sga3:xvii:corollary:3:8:item:ii}{(ii)}\(\Rightarrow\)\SGARefLink{sga3:xvii:corollary:3:8:item:i}{(i)}, suppose that \(G\) is a unipotent
algebraic group in the sense of BIBLE, and write \(G^0\) for its
identity component.  The maximal tori of \(G^0\), being made up of
unipotent elements, are trivial; thus \(G^0\) is equal to each of its
Cartan subgroups.  It follows that \(G^0\) is solvable
(BIBLE~6, Th.~6), and hence triangularizable
(BIBLE~6, Th.~1).  In short, \(G^0\) is an algebraic subgroup of a
group \(\operatorname{Trigstr}(n)\), and is therefore unipotent in
the sense of this Exposé.

The group \((G/G^0)(k)\) is a finite group consisting of unipotent
elements.  It is therefore trivial if \(p=0\), and is a finite
\(p\)-group if \(p>0\) (BIBLE~4, Prop.~4).  But then \(G/G^0\) is
unipotent in the sense of this Exposé, since it is an iterated
extension of groups isomorphic to \(\mathbb Z/p\mathbb Z\).  This
proves that \(G\) is unipotent.

\SGAProof{Completion of the proof of \SGARefLink{sga3:xvii:theorem:3-5}{3.5}}
It remains to prove \SGARefLink{sga3:xvii:theorem:3:5:item:i}{(i)}\(\Rightarrow\)\SGARefLink{sga3:xvii:theorem:3:5:item:iv}{(iv)}.

\smallskip
\noindent\textit{\SGARefLink{sga3:xvii:appendix:A}{(a)} \(p>0\).}\quad
Consider the increasing sequence of algebraic subgroups of \(G\)
\[
  \{e\}\subset F(G)\subset F^2(G)\subset\cdots\subset
  F^n(G)\subset G^0\subset G.
\]
This gives a characteristic composition series of \(G\)
(Appendix~II, 1), and for \(n\) sufficiently large,
\(G/F^n(G)\) is smooth (Appendix~II, \SGARefLink{sga3:xvii:lemma:3-1}{3.1}).  Thus its successive
quotients are, in order:
\begin{enumerate}
\item[(1)] radicial groups of height \(1\);
\item[(2)] a smooth connected group;
\item[(3)] an étale group.
\end{enumerate}
To prove \SGARefLink{sga3:xvii:lemma:3-1:item:i}{(i)}\(\Rightarrow\)(iv), it is therefore enough to prove the
following lemma.

\SGAStatement{\SGAInlineRefTarget{sga3:xvii:lemma:3-9}Lemma 3.9}{Let \(G\) be a unipotent algebraic group
defined over a field \(k\) of characteristic \(p>0\).  Then \(G\)
has a characteristic composition series, defined over \(k\), whose
successive quotients are isomorphic to:
\begin{enumerate}
\item[(i)] \SGARefTarget{sga3:xvii:lemma:3:9:item:i}
  \((\boldsymbol\alpha_p)^r\), if \(G\) is radicial;
\item[(ii)] \SGARefTarget{sga3:xvii:lemma:3:9:item:ii}
  a twisted form of \((\mathbb G_a)^r\), if \(G\) is smooth and
  connected;
\item[(iii)] \SGARefTarget{sga3:xvii:lemma:3:9:item:iii}
  a twisted form of \((\mathbb Z/p\mathbb Z)^r\), if \(G\) is
  étale.
\end{enumerate}}
\SGARefTarget{sga3:xvii:lemma:3:9}

% END inputs/XVII/tex/components/07_expose_XVII_corollaries3_7_3_8_and_lemma3_9_page009.tex

% BEGIN inputs/XVII/tex/components/08_expose_XVII_lemma3_9_proof_page010.tex
% Authority: Expo17.pdf, local p.10 / combined p.972 /
% printed pp.550--551.

\SGAProof{Proof of \SGARefLink{sga3:xvii:lemma:3:9:item:i}{3.9(i)}}
The group \(G\) is radicial.  Filtering \(G\) by the \(F^n(G)\), we
reduce to the case in which \(G\) is radicial of height \(1\).
Since \(G\) is nilpotent (\SGARefLink{sga3:xvii:theorem:3-5}{3.5}, \SGARefLink{sga3:xvii:theorem:3:5:item:i}{(i)}\(\Rightarrow\)\SGARefLink{sga3:xvii:theorem:3:5:item:iii}{(iii)}), and since
the center of an algebraic group is representable
(\SGARefLink{sga3:viii:examples:6:5:item:e:01}{Exposé VIII, 6.5(e)}), we may consider the ascending central series
of \(G\).  This series is plainly characteristic in \(G\), and
reduces us to the case in which \(G\) is also commutative.

Put \(\mathfrak g=\operatorname{Lie}(G)\).  The \(p\)-th-power
morphism \(\pi\) is then additive on \(\mathfrak g\)
(Exposé VII\textsubscript{A}); we shall reduce to the case in which
it is zero.  For every prescheme \(S\) over \(k\), set
\[
  \mathfrak g_S=\mathfrak g\otimes_k\mathcal O_S,
\]
and let \(\mathfrak h_S\) be the subsheaf of abelian groups of
\(\mathfrak g_S\) that is the image of \(\mathfrak g_S\) under
\(\pi_S\).  Finally, let \(\overline{\mathfrak h}_S\) be the
subsheaf of \(\mathcal O_S\)-modules of \(\mathfrak g_S\) generated
by \(\mathfrak h_S\).  It is clear that
\[
  \overline{\mathfrak h}_S
  =\overline{\mathfrak h}_k\otimes_k\mathcal O_S
\]
and that \(\overline{\mathfrak h}_S\)\footnote[4]{Editor's note:
\(\mathfrak h\) has been changed to \(\overline{\mathfrak h}\).}
is a characteristic \(p\)-Lie subalgebra of \(\mathfrak g_S\);
that is, it is stable under the \(S\)-functor
\(\operatorname{Aut}_{p\text{-}\mathrm{Lie}}(\mathfrak g)\).
It follows from Appendix~II, \SGARefLink{sga3:ii:subsection:2-2}{2.2} that
\(\overline{\mathfrak h}_k\) is the Lie algebra of an algebraic
subgroup \(H\) of \(G\) that is characteristic in \(G\).

Moreover, in view of \SGARefLink{sga3:xvii:corollary:3:7}{3.7} and \SGARefLink{sga3:xvii:lemma:3:9:bis}{Lemma 3.9 bis} below, if
\(G\ne\{e\}\), then \(H\ne G\), since
\(\overline{\mathfrak h}_k\ne\mathfrak g\).  On the other hand, if
\(G/H=G''\), then
\[
  \operatorname{Lie}G''
  =\mathfrak g/\overline{\mathfrak h}_k
  \qquad\text{(Appendix~II, 2.2),}
\]
and consequently the \(p\)-th power is zero in
\(\operatorname{Lie}G''\).  Induction on
\(\dim\operatorname{Lie}G\) therefore reduces us to the case in
which \(\operatorname{Lie}G\) is a \(p\)-Lie algebra whose
\(p\)-th power is zero.  But then \(\operatorname{Lie}G''\) is
isomorphic to \(\operatorname{Lie}((\boldsymbol\alpha_p)^r)\) for
some suitable integer \(r\geq0\) (Appendix~II, \SGARefLink{sga3:xvii:proposition:2:1}{2.1}), and hence
\(G''\) is isomorphic to \((\boldsymbol\alpha_p)^r\)
(Appendix~II, \SGARefLink{sga3:ii:subsection:2-2}{2.2}).  It remains to prove the following lemma.

\SGAStatement{Lemma 3.9 bis}{Let \(k\) be a field of characteristic
\(p>0\), let \(\mathfrak g\) be a commutative unipotent
\(p\)-Lie algebra (\SGARefLink{sga3:xvii:definition:3:6}{3.6}) of finite dimension over \(k\), and let
\(\mathfrak h\) be the \(p\)-Lie subalgebra of \(\mathfrak g\)
generated by the image of the \(p\)-th power in \(\mathfrak g\).
If \(\mathfrak g\ne0\), then \(\mathfrak h\ne\mathfrak g\).}
\SGARefTarget{sga3:xvii:lemma:3:9:bis}

Indeed, since \(\mathfrak g\) is commutative, \(\mathfrak h\) is
simply the \(k\)-vector subspace of \(\mathfrak g\) generated by the
\(X^{(p)}\), \(X\in\mathfrak g\).  If \(\mathfrak g\ne0\) and is
unipotent, there exists an \(X\in\mathfrak g\), \(X\ne0\), such
that \(X^{(p)}=0\).  Let \(X_1,\ldots,X_n\) be a basis of a
complement in \(\mathfrak g\) to the line \(kX\).  The Lie algebra
\(\mathfrak h\) is then the \(k\)-vector subspace of
\(\mathfrak g\) generated by
\[
  X_1^{(p)},\ldots,X_n^{(p)};
\]
it therefore has dimension at most
\(n=\dim\mathfrak g-1\).

\SGAProof{Proof of \SGARefLink{sga3:xvii:lemma:3:9:item:ii}{3.9(ii)}}
The group \(G\) is smooth and connected.  In this case the
descending central series of \(G\) is represented by smooth,
connected, characteristic algebraic subgroups \(G_i\)
(Exposé VI\textsubscript{B}, \SGARefLink{prop:sga3-VIB-8.3}{8.3} and \SGARefLink{prop:sga3-VIB-7.4}{7.4}), and \(G_i=0\) for
sufficiently large \(i\), since \(G\) is nilpotent
(\SGARefLink{sga3:xvii:theorem:3-5}{3.5}, \SGARefLink{sga3:xvii:theorem:3:5:item:i}{(i)}\(\Rightarrow\)\SGARefLink{sga3:xvii:theorem:3:5:item:iii}{(iii)}).  It is enough to prove \SGARefLink{sga3:xvii:lemma:3-9}{3.9} for the
groups \(G_i/G_{i+1}\), which reduces us to the case in which \(G\)
is also commutative.

For every integer \(n>0\), let \(G_n\) be the algebraic subgroup of
\(G\) that is the image of \(G\) under the \(p^n\)-th-power
morphism.  The group \(G_n\) is therefore smooth, connected, and
characteristic; and it follows from \SGARefLink{sga3:xvii:definition:1:1}{Definition 1.1} of unipotent
groups that \(G_n=0\) for sufficiently large \(n\).  Replacing \(G\)
by \(G_n/G_{n+1}\), we may further suppose that \(G\) is annihilated
by raising to the \(p\)-th power.  But then, by J.-P. Serre,
\emph{Groupes algébriques et corps de classes}, Chap.~VII,
Prop.~11, \(G\) is a form of \((\mathbb G_a)^r\) for a suitable
integer \(r\).

% END inputs/XVII/tex/components/08_expose_XVII_lemma3_9_proof_page010.tex

% BEGIN inputs/XVII/tex/components/09_expose_XVII_lemma3_9_iii_and_3_9_ter_page011.tex
% Authority: Expo17.pdf, local p.11 / combined p.973 /
% printed pp.552--553, stopping before Proposition 3.10.

\SGAProof{Proof of \SGARefLink{sga3:xvii:lemma:3:9:item:iii}{3.9(iii)}}
The group \(G\) is étale.  Proceeding as in (ii), we reduce first to
the case in which \(G\) is commutative, and then to the case in which
\(G\) is annihilated by \(p\).  But then
\[
  G_{\overline{k}}\simeq(\mathbb Z/p\mathbb Z)^r.
\]

\SGAProof{Proof of \SGARefLink{sga3:xvii:theorem:3-5}{3.5}, \SGARefLink{sga3:xvii:theorem:3:5:item:i}{(i)}\(\Rightarrow\)\SGARefLink{sga3:xvii:theorem:3:5:item:iv}{(iv)}, in case \SGARefLink{sga3:xvii:section:b}{(b)},
\(p=0\)}
The group \(G\) is then smooth and connected; proceeding as in
\SGARefLink{sga3:xvii:lemma:3:9:item:ii}{3.9(ii)}, we reduce to the case in which \(G\) is also commutative.
We then have the following more precise result.

\SGAStatement{Lemma 3.9 ter}{Let \(k\) be a field of characteristic
\(0\), let \(G\) be a commutative unipotent \(k\)-algebraic group,
and put \(\mathfrak g=\operatorname{Lie}G\).  There is a canonical
isomorphism
\[
  \exp:\mathbf W(\mathfrak g)\xrightarrow{\ \sim\ }G.
\]
The morphism \(\exp\) is the unique homomorphism
\(\mathbf W(\mathfrak g)\to G\) that induces the identity on Lie
algebras.}
\SGARefTarget{sga3:xvii:lemma:3:9:ter}

Since \(G\) is unipotent, it can be realized as an algebraic
subgroup of \(\operatorname{Trigstr}(n)\), for a suitable integer
\(n\) (\SGARefLink{sga3:xvii:theorem:3-5}{3.5}, \SGARefLink{sga3:xvii:theorem:3:5:item:i}{(i)}\(\Rightarrow\)\SGARefLink{sga3:xvii:theorem:3:5:item:v}{(v)}).  Choosing such an embedding
identifies \(\mathfrak g\) with a Lie subalgebra of
\(\mathfrak{gl}(n)\) consisting of nilpotent endomorphisms.  We
thereby obtain a \(k\)-morphism
\[
  \begin{aligned}
  \exp:\mathbf W(\mathfrak g)&\longrightarrow
    \operatorname{GL}(n),\\
  T&\longmapsto\sum_{i\geq0}\frac{T^i}{i!}.
  \end{aligned}
\]

Since \(G\) is commutative, the morphism \(\exp\) is a
homomorphism.  Let \(G'\) be the algebraic group that is the image
of \(\mathbf W(\mathfrak g)\) under \(\exp\).  Under the canonical
identification
\(\operatorname{Lie}\mathbf W(\mathfrak g)=\mathfrak g\), the
tangent linear map of \(\exp\) is simply the inclusion
\(\mathfrak g\hookrightarrow\mathfrak{gl}(n)\).  Hence
\[
  \operatorname{Lie}(G\cap G')
  =\mathfrak g\cap\mathfrak g=\mathfrak g.
\]
Since \(G\cap G'\) is smooth (Exposé VI\textsubscript{B}, \SGARefLink{cor:sga3-VIB-1.6.1}{1.6.1})
and \(G\) is connected---being an iterated extension of groups
\(\mathbb G_a\) (\SGARefLink{sga3:xvii:theorem:3-5}{3.5}, \SGARefLink{sga3:xvii:theorem:3:5:item:i}{(i)}\(\Rightarrow\)\SGARefLink{sga3:xvii:theorem:3:5:item:ii}{(ii)})---one necessarily has
\(G=G'\).  The kernel \(\ker(\exp)\) is an étale unipotent group,
and hence is the trivial group.  Consequently \(\exp\) is an
isomorphism from \(\mathbf W(\mathfrak g)\) onto \(G\).

If \(h:\mathbf W(\mathfrak g)\to G\) is another homomorphism such
that \(\operatorname{Lie}(h)\) is the identity map of
\(\mathfrak g\), then \(h-\exp\) is a homomorphism, since \(G\) is
commutative, and its tangent linear map is zero.  Because \(k\) has
characteristic \(0\) and \(\mathbf W(\mathfrak g)\) is connected,
Cartier's theorem again implies that \(h=\exp\).

% END inputs/XVII/tex/components/09_expose_XVII_lemma3_9_iii_and_3_9_ter_page011.tex

% BEGIN inputs/XVII/tex/components/10_expose_XVII_proposition3_10_through_section4_opening_page012.tex
% Authority: Expo17.pdf, local pp.11--12 / combined pp.973--974 /
% printed pp.553--556, stopping before subsection 4.1.

\SGAStatement{\SGAInlineRefTarget{sga3:xvii:proposition:3-10}Proposition 3.10}{Let \(G\) be an algebraic group
defined over an algebraically closed field \(k\).  The following
properties are equivalent:
\begin{enumerate}
\item[(i)] \SGARefTarget{sga3:xvii:proposition:3:10:item:i}
  Every homomorphism from a group \(M\) of multiplicative type into
  \(G\) is trivial.
\item[(i bis)] \SGARefTarget{sga3:xvii:proposition:3:10:item:i:bis}
  \(G\) has no nontrivial subgroup of multiplicative type.
\item[(ii)] \SGARefTarget{sga3:xvii:proposition:3:10:item:ii}
  \begin{enumerate}
  \item[(a)] \SGARefTarget{sga3:xvii:proposition:3:10:item:ii:a}
    If \(p=0\), then \(G(k)\) has no point of finite order other
    than \(e\).
  \item[(b)] \SGARefTarget{sga3:xvii:proposition:3:10:item:ii:b}
    If \(p\ne0\), then:
    \begin{itemize}
    \item for every prime number \(q\ne p\), if \(x\in G(k)\) and
      \(x^q=e\), then \(x=e\);
    \item \SGARefTarget{sga3:xvii:proposition:3:10:item:ii:b:lie}
      for every \(X\in\mathfrak g=\operatorname{Lie}G\) such that
      \(X^{(p)}=X\), one has \(X=0\).
    \end{itemize}
  \end{enumerate}
\item[(ii bis)] \SGARefTarget{sga3:xvii:proposition:3:10:item:ii:bis}
  \begin{enumerate}
  \item[(a)] \SGARefTarget{sga3:xvii:proposition:3:10:item:ii:bis:a}
    If \(p=0\), the condition is the same as in \SGARefLink{sga3:xvii:proposition:3:10:item:ii}{(ii)}\SGARefLink{sga3:xvii:appendix:A}{(a)}.
  \item[(b)] \SGARefTarget{sga3:xvii:proposition:3:10:item:ii:bis:b}
    If \(p\ne0\), then:
    \begin{itemize}
    \item every finite subgroup of \(G(k)\) is a \(p\)-group;
    \item \SGARefTarget{sga3:xvii:proposition:3:10:item:ii:bis:b:lie}
      for every \(X\in\mathfrak g=\operatorname{Lie}G\) such that
      \(X^{(p)}=X\), one has \(X=0\).
    \end{itemize}
  \end{enumerate}
\end{enumerate}}
\SGARefTarget{sga3:xvii:proposition:3:10}

\SGAProof{Proof}
The equivalence \SGARefLink{sga3:xvii:proposition:3:10:item:i}{(i)}\(\Leftrightarrow\)(i bis) follows because the
image of a group of multiplicative type is of multiplicative type
(\SGARefLink{sga3:ix:proposition:2:7}{Exposé IX, 2.7}).

The equivalence \SGARefLink{sga3:ix:proposition:2:7:item:ii:01}{(ii)}\(\Leftrightarrow\)(ii bis) follows from the
following observation.  If an ordinary finite group \(H\) has order
that is not a power of \(p\), then there exist a prime number
\(q\ne p\) and an element \(x\in H\), distinct from \(e\), such
that \(x^q=e\) (Sylow's theorem; cf. J.-P. Serre,
\emph{Corps locaux}, Chapter IX, \S2).

The implication \SGARefLink{sga3:xvii:proposition:3:10:item:i}{(i)}\(\Rightarrow\)\SGARefLink{sga3:xvii:proposition:3:10:item:ii}{(ii)} follows from the next lemma.

\SGAStatement{\SGAInlineRefTarget{sga3:xvii:lemma:3-11}Lemma 3.11}{Let \(G\) be an algebraic group defined
over a field \(k\).
\begin{enumerate}
\item[(a)] \SGARefTarget{sga3:xvii:lemma:3:11:item:a}
  If \(k\) contains the \(n\)-th roots of unity, where \(n\) is
  prime to \(p\), then
  \[
    \operatorname{Hom}_{k\text{-}\mathrm{gr}}(\mu_n,G)
      \simeq
    \operatorname{Hom}_{k\text{-}\mathrm{gr}}
      (\mathbb Z/n\mathbb Z,G)
      \simeq {}_nG(k),
  \]
  where \({}_nG(k)\) denotes the points of order \(n\) in \(G(k)\).
\item[(b)] \SGARefTarget{sga3:xvii:lemma:3:11:item:b}
  If \(p>0\), then
  \[
    \operatorname{Hom}_{k\text{-}\mathrm{gr}}(\mu_p,G)
      \simeq
    \{X\in\mathfrak g=\operatorname{Lie}G\mid X^{(p)}=X\}.
  \]
\end{enumerate}}
\SGARefTarget{sga3:xvii:lemma:3:11}

Indeed, for \SGARefLink{sga3:xvii:appendix:A}{(a)}, note that \(\mu_n\) is then isomorphic over \(k\)
to \(\mathbb Z/n\mathbb Z\); assertion \SGARefLink{sga3:xvii:section:b}{(b)} follows from Appendix II,
\SGARefLink{sga3:xvii:proposition:2:1}{2.1}.

For \SGARefLink{sga3:xvii:proposition:2:1:item:ii}{(ii)}\(\Rightarrow\)(i bis), \SGARefLink{sga3:xvii:lemma:3-11}{Lemma 3.11} shows that \SGARefLink{sga3:xvii:proposition:2:1:item:ii}{(ii)} is
equivalent to the assertion that, for every prime number \(r\), the
group \(G\) contains no subgroup \(\mu_r\).  The implication now
follows from:

\SGAStatement{Lemma 3.12}{Let \(G\) be a diagonalizable
\(S\)-group, of finite type over \(S\), and distinct from the
trivial group.  There exist a prime number \(r\) and a subgroup of
\(G\) isomorphic to \((\mu_r)_S\).}
\SGARefTarget{sga3:xvii:lemma:3:12}

Write \(G=D_S(M)\), where \(M\) is a finitely generated abelian
group, hence an extension of a free group \(M''\) by a finite group
\(M'\).  If \(M''\ne0\), it is clear that \(M\) has quotients
isomorphic to \(\mathbb Z/r\mathbb Z\) for every integer \(r\).  If
\(M''=0\), then \(M'\) has a quotient isomorphic to
\(\mathbb Z/r\mathbb Z\) for every prime number \(r\) dividing the
order of \(M\).  The lemma follows by duality.

We saw in \SGARefLink{sga3:xvii:proposition:2:4}{Proposition 2.4} that a unipotent group satisfies the
equivalent conditions of \SGARefLink{sga3:xvii:proposition:3-10}{3.10}.  The purpose of the next section is
to prove the converse.

\SGASection{4}{A Characterization of Unipotent Groups}
\SGARefTarget{sga3:xvii:section:4}

As announced, we shall prove that an algebraic group \(G\) defined
over an algebraically closed field \(k\), and containing no
nontrivial subgroup of multiplicative type, is unipotent.  In fact,
it is enough that \(G\) contain none of certain particular
``elementary'' subgroups of multiplicative type, which depend on the
hypotheses imposed on \(G\).  Before stating the general theorem,
let us study some special cases in detail.

% END inputs/XVII/tex/components/10_expose_XVII_proposition3_10_through_section4_opening_page012.tex

% BEGIN inputs/XVII/tex/components/11_expose_XVII_proposition4_1_1_and_lemma4_1_2_page013.tex
% Authority: Expo17.pdf, local p.13 and local p.14 top /
% combined pp.975--976 / printed pp.557--558.

\SGASubsection{4.1}{Smooth, Connected, Affine Algebraic Groups}
\SGARefTarget{sga3:xvii:subsection:4:1}

\SGAStatement{\SGAInlineRefTarget{sga3:xvii:proposition:4-1-1}Proposition 4.1.1}{Let \(k\) be a field, let \(G\)
be a smooth, connected, affine \(k\)-algebraic group, and put
\(\mathfrak g=\operatorname{Lie}G\).  The following properties are
equivalent:
\begin{enumerate}
\item[(i)] \SGARefTarget{sga3:xvii:proposition:4:1:1:item:i}
  \(G\) is unipotent.
\item[(ii)] \SGARefTarget{sga3:xvii:proposition:4:1:1:item:ii}
  \(G\) has a central composition series whose successive quotients
  are forms of \(\mathbb G_a\).
\item[(iii)] \SGARefTarget{sga3:xvii:proposition:4:1:1:item:iii}
  \(G\) has a central, characteristic composition series whose
  successive quotients are forms of \((\mathbb G_a)^r\).
\item[(iv)] \SGARefTarget{sga3:xvii:proposition:4:1:1:item:iv}
  There is an integer \(n>1\) such that \(G_{\overline{k}}\)
  contains no subgroup isomorphic to \(\mu_n\).
\item[(v)] \SGARefTarget{sga3:xvii:proposition:4:1:1:item:v}
  Every maximal torus of \(G\) is the trivial group.
\end{enumerate}
Suppose, moreover, that \(G\) is an algebraic subgroup of a linear
group \(\operatorname{GL}(n)\).  The preceding conditions are then
also equivalent to:
\begin{enumerate}
\item[(vi)] \SGARefTarget{sga3:xvii:proposition:4:1:1:item:vi}
  \(\mathfrak g\subset\mathfrak{gl}(n)\) consists of nilpotent
  endomorphisms.
\item[(vii)] \SGARefTarget{sga3:xvii:proposition:4:1:1:item:vii}
  \(\mathfrak g\) is nilpotent, and its center contains no nonzero
  semisimple endomorphism.
\end{enumerate}}
\SGARefTarget{sga3:xvii:proposition:4:1:1}

\SGAProof{Proof}
The implication \SGARefLink{sga3:xvii:proposition:4:1:1:item:ii}{(ii)}\(\Rightarrow\)\SGARefLink{sga3:xvii:proposition:4:1:1:item:i}{(i)} is clear, and
\SGARefLink{sga3:xvii:proposition:4:1:1:item:i}{(i)}\(\Rightarrow\)\SGARefLink{sga3:xvii:proposition:4:1:1:item:iii}{(iii)} was proved in \SGARefLink{sga3:xvii:lemma:3-9}{3.9}.  The implication
\SGARefLink{sga3:xvii:lemma:3:9:item:iii}{(iii)}\(\Rightarrow\)\SGARefLink{sga3:xvii:lemma:3:9:item:ii}{(ii)} will follow from the next lemma.

\SGAStatement{Lemma 4.1.2}{Let \(k\) be a field, and let \(G\) be
a \(k\)-algebraic group that is a form of \((\mathbb G_a)^r\).
Then:
\begin{enumerate}
\item[(a)] \SGARefTarget{sga3:xvii:lemma:4:1:2:item:a}
  \(G\) can be realized as an algebraic subgroup of
  \((\mathbb G_a)^n\), for a suitable integer \(n\).
\item[(b)] \SGARefTarget{sga3:xvii:lemma:4:1:2:item:b}
  \(G\) has a composition series whose successive quotients are
  forms of \(\mathbb G_a\).
\end{enumerate}}
\SGARefTarget{sga3:xvii:lemma:4:1:2}

Indeed, by hypothesis there is an extension \(k'\) of \(k\) such
that \(G_{k'}\) is isomorphic to \((\mathbb G_{a,k'})^r\).  By the
principle of finite extension (EGA IV, 9.1.1), we may assume that
\(k'/k\) is finite.  For (a), it then suffices to consider the
canonical closed immersion (EGA V\footnote[5]{Editor's note:
provide another reference here\ldots})
\[
  G\longrightarrow
  \prod_{k'/k}(G_{k'})_{/k'}
  \xrightarrow{\ \sim\ }(\mathbb G_{a,k})^n,
  \qquad n=r\deg(k'/k).
\]

To prove \SGARefLink{sga3:xvii:lemma:4:1:2:item:b}{(b)}, in view of \SGARefLink{sga3:xvii:lemma:4:1:2:item:a}{(a)}, we may assume that \(G\) is a closed
subgroup of \(G'=(\mathbb G_a)^n\).  If \(G\ne0\), there is a
hyperplane \(\mathfrak h\) in
\(\mathfrak g'=\operatorname{Lie}G'\) that does not contain
\(\mathfrak g\).  Let \(H\) be the vector subgroup
\(\mathbf W(\mathfrak h)\) of
\(\mathbf W(\mathfrak g')=G'\).  Since \(H\) is defined by one
equation in \(G'\), the intersection \(G\cap H\) is defined by one
equation in \(G\), and
\[
\begin{aligned}
  \dim G-1
  &\leq \dim(G\cap H)\\
  &\leq \dim\operatorname{Lie}(G\cap H)
   =\dim_k(\mathfrak g\cap\mathfrak h)\\
  &=\dim_k\mathfrak g-1
   =\dim G-1.
\end{aligned}
\]
Thus
\(\dim(G\cap H)=\dim\operatorname{Lie}(G\cap H)\), and consequently
\(G\cap H\) is smooth.  The group
\[
  G_1=(G\cap H)^0
\]
is a smooth, connected algebraic subgroup of \(G\) such that
\(G/G_1\) is smooth, connected, and one-dimensional, hence is a
form of \(\mathbb G_a\) (\SGARefLink{sga3:xvii:subsection:4:1}{4.1}, \SGARefLink{sga3:xvii:proposition:4:1:1:item:i}{(i)}\(\Rightarrow\)\SGARefLink{sga3:xvii:proposition:4:1:1:item:iii}{(iii)}).  Induction
on the dimension of \(G\) completes the proof.

% END inputs/XVII/tex/components/11_expose_XVII_proposition4_1_1_and_lemma4_1_2_page013.tex

% BEGIN inputs/XVII/tex/components/12_expose_XVII_corollary4_1_3_through_lemma4_1_5_page014.tex
% Authority: Expo17.pdf, local p.14 / combined p.976 /
% printed pp.558--560, stopping before the proof of Lemma 4.1.5.

Before continuing the proof of \SGARefLink{sga3:xvii:subsection:4:1}{4.1}, note that the equivalence
(i)\(\Leftrightarrow\)(ii), together with \SGARefLink{sga3:xvii:lemma:2:3:bis}{2.3 bis}, gives the
following corollary.

\SGAStatement{Corollary 4.1.3}{If \(k\) is a perfect field, a
smooth, connected \(k\)-algebraic group is unipotent if and only if
it has a composition series whose successive quotients are
isomorphic to \(\mathbb G_a\).}
\SGARefTarget{sga3:xvii:corollary:4:1:3}

\SGAProof{Continuation of the proof of \SGARefLink{sga3:xvii:subsection:4:1}{4.1}}
The implication (i)\(\Rightarrow\)(iv) follows from \SGARefLink{sga3:xiv:2.4:item:i}{2.4(i)}.

For (iv)\(\Rightarrow\)(v), by \SGARefLink{sga3:xiv:4.1}{Exposé XIV, 4.1}, the group \(G\)
has a maximal torus \(T\) defined over \(k\).  If \(r=\dim T\),
then
\[
  ({}_nT)_{\overline{k}}\simeq(\mu_n)^r.
\]
It follows that \(r=0\).

The implication (v)\(\Rightarrow\)(i) follows from the observation
made in the proof of \SGARefLink{sga3:xvii:corollary:3:8}{3.8}.

For (i)\(\Rightarrow\)(vi), by \SGARefLink{sga3:xvii:corollary:3:4}{3.4} the group \(G\) is in fact
contained in an algebraic subgroup of \(\operatorname{GL}(n)\)
isomorphic to \(\operatorname{Trigstr}(n)\).  Hence
\(\mathfrak g\) consists of nilpotent endomorphisms.

For (vi)\(\Rightarrow\)(vii), the Lie algebra \(\mathfrak g\) is
nilpotent by Engel's theorem (Bourbaki, \emph{Groupes et algèbres
de Lie}, Chapter I, \S4, Corollary 3).

For (vii)\(\Rightarrow\)(v), let \(T\) be a maximal subtorus of
\(G\) (\SGARefLink{sga3:xiv:1.1}{Exposé XIV, 1.1}), and let \(\mathfrak t\) be its Lie
algebra.  The embedding of \(G\) in \(\operatorname{GL}(n)\)
defines a representation of \(T\), which is necessarily
semisimple; this can be seen after passing to an algebraic closure
\(\overline{k}\) of \(k\) and applying \SGARefLink{sga3:i:proposition:4-7-3}{Exposé I, 4.7.3}.  Thus, if
\(X\in\mathfrak t\), then \(X\) is a semisimple endomorphism in
\(\mathfrak{gl}(n)\).  It follows at once that
\[
  \operatorname{ad}X:\ Y\longmapsto[X,Y]
\]
is a semisimple endomorphism of \(\mathfrak{gl}(n)\), and hence of
\(\mathfrak g\).  This endomorphism is also nilpotent, since
\(\mathfrak g\) is nilpotent.  Therefore
\(\operatorname{ad}X=0\), so \(X\) is central.  Consequently
\(\mathfrak t\) is central and consists of semisimple
endomorphisms; by hypothesis it is zero.  A fortiori, \(T\) is the
trivial group.

\SGAStatement{Remark 4.1.4}{%
\begin{enumerate}
\item[(a)] \SGARefTarget{sga3:xvii:remark:4:1:4:item:a}
  We shall give below (\SGARefLink{sga3:xvii:proposition:4-3-1}{4.3.1}) an infinitesimal characterization of
  unipotent groups in characteristic \(p>0\) that is independent of
  an embedding in \(\operatorname{GL}(n)\).
\item[(b)] \SGARefTarget{sga3:xvii:remark:4:1:4:item:b}
  When \(k\) is perfect, conditions (ii) and (iii) of \SGARefLink{sga3:xvii:proposition:4-1-1}{4.1.1}
  simplify by virtue of the following lemma.
\end{enumerate}}
\SGARefTarget{sga3:xvii:remark:4:1:4}

\SGAStatement{Lemma 4.1.5}{If \(k\) is a perfect field, every
\(k\)-algebraic group \(G\) that is a form of
\((\mathbb G_a)^r\) is isomorphic to \((\mathbb G_a)^r\).}
\SGARefTarget{sga3:xvii:lemma:4:1:5}

% END inputs/XVII/tex/components/12_expose_XVII_corollary4_1_3_through_lemma4_1_5_page014.tex

% BEGIN inputs/XVII/tex/components/13_expose_XVII_lemma4_1_5_proof_page015.tex
% Authority: Expo17.pdf, local pp.14--15 / combined pp.976--977 /
% printed pp.560--561, stopping before the two-case degree argument.

The lemma follows from \SGARefLink{sga3:xvii:lemma:3:9:ter}{3.9 ter} when the characteristic \(p\) of
\(k\) is zero, and from \SGARefLink{sga3:xvii:lemma:2:3:bis}{2.3 bis} when \(r=1\).  In the general case
\((p>0)\), realize \(G\) as an algebraic subgroup of
\((\mathbb G_a)^n\), for a suitable integer \(n\) (\SGARefLink{sga3:xvii:lemma:4:1:2}{4.1.2}), and
argue by induction on \(n-r\).  If \(r=n\), then indeed
\(G=(\mathbb G_a)^r\).  Otherwise, the quotient
\((\mathbb G_a)^n/G\) is a nonzero smooth, connected, unipotent
group.  By \SGARefLink{sga3:xvii:proposition:4-1-1}{4.1.1}, \SGARefLink{sga3:xvii:proposition:4:1:1:item:i}{(i)}\(\Rightarrow\)\SGARefLink{sga3:xvii:proposition:4:1:1:item:ii}{(ii)}, and \SGARefLink{sga3:xvii:lemma:2:3:bis}{2.3 bis}, it has a
composition series whose quotients are isomorphic to
\(\mathbb G_a\).  It follows that there is a smooth, connected
algebraic subgroup \(G_1\) of \((\mathbb G_a)^n\), containing \(G\),
such that
\[
  (\mathbb G_a)^n/G_1=\mathbb G_a.
\]
By induction, it suffices to show that \(G_1\) is isomorphic to
\((\mathbb G_a)^{n-1}\).

It is immediate that a homomorphism
\[
  \operatorname{Spec}k[X_1,\ldots,X_n]=(\mathbb G_a)^n
  \longrightarrow
  \mathbb G_a=\operatorname{Spec}k[T]
\]
is defined by an additive polynomial of the form
\[
  \sum_{i,j}a_{i,j}X_i^{p^j}.
\]
Since \(G_1\) is smooth, the linear part of this polynomial is
nonzero.  After a linear change of the coordinates \(X_i\), we may
assume that \(G_1\) is the algebraic subgroup of
\((\mathbb G_a)^n\) defined by
\[
  \tag{\(*\)}
  P(X)=-X_1+\sum_{j=1}^{q}a_jX_1^{p^j}
    +Q(X_2,\ldots,X_n)=0,
\]
\SGARefTarget{sga3:xvii:lemma:4:1:5:eq:star}
where
\[
  Q(X_2,\ldots,X_n)
    =\sum_{\substack{i>1\\j>0}}b_{i,j}X_i^{p^j}.
\]

We now argue by induction on the degree of \(P\).  If
\(\deg P=1\), it is clear that
\[
  G_1\xrightarrow{\ \sim\ }(\mathbb G_a)^{n-1}.
\]
Otherwise, since \(k\) is perfect, one has
\(Q(X)=Q_1(X)^p\), and we may define an endomorphism \(v\) of
\((\mathbb G_a)^n\) by
\[
\begin{aligned}
  X_i&\longmapsto X_i &&\text{for \(i>1\)},\\
  X_1&\longmapsto
    \sum_{j=1}^{q}a_j^{1/p}X_1^{p^{j-1}}+Q_1(X).
\end{aligned}
\]
\SGARefTarget{sga3:xvii:lemma:4:1:5:map:v}
It is clear that \(v\) induces an isomorphism on \(G_1\), and that
\(v(G_1)\) is defined in \((\mathbb G_a)^n\) by the equation
\[
  \tag{\(*\)\(_1\)}
  P_1(X)=-X_1+\sum_{j=1}^{q}a_j^{1/p}X_1^{p^j}
    +Q_1(X_2,\ldots,X_n)=0.
\]
\SGARefTarget{sga3:xvii:lemma:4:1:5:eq:star:1}

% END inputs/XVII/tex/components/13_expose_XVII_lemma4_1_5_proof_page015.tex

% BEGIN inputs/XVII/tex/components/14_expose_XVII_lemma4_1_5_completion_pages015_016.tex
% Authority: Expo17.pdf, local pp.15--16 / combined pp.977--978 /
% printed pp.561--562, completing the proof of Lemma 4.1.5.

We distinguish two cases.

\medskip\noindent\textit{First case.}\quad
\(\deg Q=\deg P>p^q\).  Then \(\deg P_1<\deg P\), and the
induction hypothesis applies.

\medskip\noindent\textit{Second case.}\quad
\(\deg P=p^q\), with \(a_q\ne0\).  Then
\[
  \deg P_1=\deg P=p^q
  \qquad\text{and}\qquad
  \deg Q_1<p^q.
\]
We cannot have \(Q=0\), for otherwise \(G_1\) would not be
connected.  If the polynomial \(Q_1\) has no linear part, the
preceding transformation may be repeated.  Continuing in this
fashion, we eventually obtain an equation of the form
\[
  \tag{\(*\)\(_s\)}\label{sga3:xvii:lemma:4-1-5:equation:star-s}
  P_s(X)=-X_1+\sum_{j=1}^{q}a_j^{1/p^s}X_1^{p^j}+Q_s(X),
\]
\SGARefTarget{sga3:xvii:lemma:4:1:5:eq:star:s}
where
\[
  Q_s(X)=Q(X_2,\ldots,X_n)^{1/p^s}
\]
is an additive polynomial with a nonzero linear part, and moreover
\(\deg Q_s<p^q\).  Suppose, for example, that the coefficient of
\(X_2\) in \(Q_s\) is nonzero, and let \(-L\) be the linear part
of \(P_s(X)\).  After a linear change of coordinates, the equation
of \(G_1\) becomes
\[
  P'(X)=-L+\sum_{j=1}^{q}a_j^{1/p^s}X_1^{p^j}
    +Q'(L,X_3,\ldots,X_n),
\]
where \(Q'\) is an additive polynomial without a linear part and
\(\deg Q'<p^q\).  We are thereby reduced to the first
case.\footnote[6]{Editor's note: replace \(X_1\) by \(L\).}


% END inputs/XVII/tex/components/14_expose_XVII_lemma4_1_5_completion_pages015_016.tex

% BEGIN inputs/XVII/tex/components/15_expose_XVII_subsection4_2_through_height1_case_pages016_017.tex
% Authority: Expo17.pdf, local pp.16--17 / combined pp.978--979 /
% printed pp.562--563 and the top of printed p.564.

\SGASubsection{4.2}{Radicial Groups}
\SGARefTarget{sga3:xvii:subsection:4:2}

\SGAStatement{\SGAInlineRefTarget{sga3:xvii:proposition:4-2-1}Proposition 4.2.1}{Let \(G\) be a radicial algebraic
group defined over a field \(k\) of characteristic \(p>0\).  The
following conditions are equivalent:
\begin{enumerate}
\item[(i)] \SGARefTarget{sga3:xvii:proposition:4:2:1:item:i}
  \(G\) is unipotent.
\item[(ii)] \SGARefTarget{sga3:xvii:proposition:4:2:1:item:ii}
  \(G\) has a central composition series whose successive
  quotients are isomorphic to \(\boldsymbol\alpha_p\).
\item[(iii)] \SGARefTarget{sga3:xvii:proposition:4:2:1:item:iii}
  \(G\) has a central, characteristic composition series whose
  successive quotients are isomorphic to
  \((\boldsymbol\alpha_p)^r\).
\item[(iv)] \SGARefTarget{sga3:xvii:proposition:4:2:1:item:iv}
  \(G_{\overline{k}}\) contains no subgroup isomorphic to
  \(\boldsymbol\mu_p\).
\item[(v)] \SGARefTarget{sga3:xvii:proposition:4:2:1:item:v}
  \(\mathfrak g=\operatorname{Lie}G\) is a unipotent
  \(p\)-Lie algebra (\SGARefLink{sga3:xvii:definition:3:6}{3.6}).
\end{enumerate}}
\SGARefTarget{sga3:xvii:proposition:4:2:1}

\SGAProof{Proof}
The implications
\SGARefLink{sga3:xvii:proposition:4:2:1:item:iii}{(iii)}\(\Rightarrow\)\SGARefLink{sga3:xvii:proposition:4:2:1:item:ii}{(ii)}\(\Rightarrow\)\SGARefLink{sga3:xvii:proposition:4:2:1:item:i}{(i)} are clear;
\SGARefLink{sga3:xvii:proposition:4:2:1:item:i}{(i)}\(\Rightarrow\)\SGARefLink{sga3:xvii:proposition:4:2:1:item:iii}{(iii)} is \SGARefLink{sga3:xvii:lemma:3:9:item:i}{3.9(i)}, and
\SGARefLink{sga3:xvii:proposition:4:2:1:item:i}{(i)}\(\Rightarrow\)\SGARefLink{sga3:xvii:proposition:4:2:1:item:iv}{(iv)} follows from \SGARefLink{sga3:xvii:proposition:2:4:item:i}{2.4(i)}.

We shall need the following lemma on abelian \(p\)-Lie algebras.

\SGAStatement{\SGAInlineRefTarget{sga3:xvii:lemma:4-2-2}Lemma 4.2.2}{Let \(\mathfrak g\) be an abelian,
finite-dimensional \(p\)-Lie algebra over a perfect field \(k\).
Then \(\mathfrak g\) can be written uniquely as a direct sum
\[
  \mathfrak g=\mathfrak r\oplus\mathfrak u
\]
of a \(p\)-Lie subalgebra \(\mathfrak r\) on which the \(p\)-th
power is bijective and a unipotent \(p\)-Lie subalgebra
\(\mathfrak u\) (\SGARefLink{sga3:xvii:definition:3:6}{3.6}).  The algebra \(\mathfrak r\) is called the
\emph{reductive part} of \(\mathfrak g\), and \(\mathfrak u\) its
\emph{unipotent part}.  The formation of \(\mathfrak r\) and
\(\mathfrak u\) is compatible with extension of the field \(k\).
If, moreover, \(k\) is algebraically closed, then
\(\mathfrak r\) has a basis \(e_i\) such that
\[
  e_i^{(p)}=e_i.
\]}
\SGARefTarget{sga3:xvii:lemma:4:2:2}

The proof of this lemma is easy and is left to the reader; see
Bourbaki, \emph{Groupes et algèbres de Lie}, Chapter I, \S1,
Exercise 23.  We merely note that \(\mathfrak u\) is the kernel of
a suitable iterate of the map \(X\mapsto X^{(p)}\), while
\(\mathfrak r\) is the image of that same iterate.

We now prove \SGARefLink{sga3:xvii:proposition:4:2:1:item:iv}{(iv)}\(\Rightarrow\)\SGARefLink{sga3:xvii:proposition:4:2:1:item:v}{(v)}.  After extending the base
field if necessary, we may suppose that \(k\) is algebraically
closed.  Let \(X\in\mathfrak g\), and let \(\mathfrak h\) be the
\(p\)-Lie subalgebra of \(\mathfrak g\) generated by \(X\).  The
algebra \(\mathfrak h\) is plainly commutative, and its reductive
part (\SGARefLink{sga3:xvii:lemma:4-2-2}{4.2.2}) is zero.  Otherwise, by the cited result,
\(\mathfrak h\) would contain an element \(Y\ne0\) such that
\(Y^{(p)}=Y\), and consequently, by Appendix II, \SGARefLink{sga3:xvii:proposition:2:1}{2.1} and \SGARefLink{sga3:xvii:proposition:2:2}{2.2},
\(G\) would contain a subgroup isomorphic to
\(\boldsymbol\mu_p\), contrary to the hypothesis.  Thus
\(\mathfrak h\), and hence \(\mathfrak g\), is a unipotent
\(p\)-Lie algebra.

It remains to prove \SGARefLink{sga3:xvii:proposition:4:2:1:item:v}{(v)}\(\Rightarrow\)\SGARefLink{sga3:xvii:proposition:4:2:1:item:i}{(i)}, the least immediate
implication in \SGARefLink{sga3:xvii:proposition:4-2-1}{4.2.1}.

\medskip\noindent\textit{(a) The case in which \(G\) has height
\(1\)} (Exposé VII\textsubscript{A}, \SGARefLink{sga3:VIIA:4.1.3}{4.1.3}).
Since \(G\) is radicial, it is affine and hence isomorphic to an
algebraic subgroup of a linear group \(\operatorname{GL}(V)\)
(Exposé VI\textsubscript{B}, \S11).  This embedding identifies
\(\mathfrak g\) with a \(p\)-Lie subalgebra of
\(\operatorname{End}(V)\), and the \(p\)-th power of \(X\) in
\(\mathfrak g\) agrees with the \(p\)-th power of the endomorphism
\(X\) (Exposé VII\textsubscript{A}, \SGARefLink{sga3:VIIA:6.4.4}{6.4.4}).  Since
\(\mathfrak g\) is unipotent by hypothesis, it is therefore an
algebra of nilpotent endomorphisms of \(V\).  By Engel's theorem
(Bourbaki, \emph{Groupes et algèbres de Lie}, Chapter I, \S4,
Theorem 1), it is a Lie subalgebra of the Lie algebra
\(\mathfrak h\) of the group of strictly upper-triangular matrices
\(\operatorname{Trigstr}(n)\), relative to a suitable basis of
\(V\).  Since \(G\) has height \(1\), Appendix II, \SGARefLink{sga3:ii:subsection:2-2}{2.2} then shows
that \(G\) itself is an algebraic subgroup of
\(\operatorname{Trigstr}(n)\), and hence is unipotent by
\SGARefLink{sga3:xvii:theorem:3:5:item:v}{3.5(v)}\(\Rightarrow\)\SGARefLink{sga3:xvii:proposition:4:2:1:item:i}{(i)}.


% END inputs/XVII/tex/components/15_expose_XVII_subsection4_2_through_height1_case_pages016_017.tex

% BEGIN inputs/XVII/tex/components/16_expose_XVII_proposition4_2_1_general_case_page017.tex
% Authority: Expo17.pdf, local p.17 / combined p.979 /
% printed p.564, completing Proposition 4.2.1 and stopping before 4.3.

\medskip\noindent\textit{\SGARefLink{sga3:xvii:section:b}{(b)} The general case.}\quad
We argue by induction on the height \(h\) of
\(G\).{\renewcommand{\thefootnote}{*}\footnote{That is, the least
integer \(h\) such that \(F^h=\operatorname{id}_G\); see
Appendix II, 1.}}
The case \(h=1\) has just been treated.  Suppose that \(h>1\), and
put
\[
  G'={}_{F}G,
  \qquad
  G''=G/G'.
\]
The group \(G'\) has height \(1\), and
\(\operatorname{Lie}G'=\operatorname{Lie}G\) is unipotent.
Consequently \(G'\) is unipotent by \SGARefLink{sga3:xvii:appendix:A}{(a)}.  To prove that \(G\) is
unipotent, it is therefore enough to show that \(G''\) is unipotent
(\SGARefLink{sga3:xvii:proposition:2:2}{2.2}).  But \(G''\) has height \(h-1\), so the induction
hypothesis reduces us to proving that \(\operatorname{Lie}G''\) is
unipotent.

In view of (iv)\(\Rightarrow\)(v), it suffices to prove that
\(G''_{\overline{k}}\) contains no subgroup isomorphic to
\(\boldsymbol\mu_p\).  Suppose, then, that such a subgroup has been
given, and let \(H\) be its inverse image in
\(G_{\overline{k}}\).  Since \(G'\) is unipotent, we shall prove
in \S5 that the extension
\[
  e\longrightarrow G'_{\overline{k}}
   \longrightarrow H
   \longrightarrow\boldsymbol\mu_p
   \longrightarrow e
\]
is necessarily trivial; the proof of that fact is independent of
the results of the present section.  In short,
\(\boldsymbol\mu_p\) lifts to \(H\).  Being of height \(1\), it is
then necessarily contained in
\[
  G'_{\overline{k}}={}_{F}(G_{\overline{k}}),
\]
which is impossible because \(G'\) is unipotent.

% END inputs/XVII/tex/components/16_expose_XVII_proposition4_2_1_general_case_page017.tex

% BEGIN inputs/XVII/tex/components/17_expose_XVII_proposition4_3_1_and_lemma4_3_2_pages017_018.tex
% Authority: Expo17.pdf, local pp.17--18 / combined pp.979--980 /
% printed pp.564--566, through Lemma 4.3.2.

\SGASubsection{4.3}{Connected Affine Groups in Characteristic \(p>0\)}
\SGARefTarget{sga3:xvii:subsection:4:3}

\SGAStatement{\SGAInlineRefTarget{sga3:xvii:proposition:4-3-1}Proposition 4.3.1}{Let \(G\) be a connected affine
algebraic group over a field \(k\) of characteristic \(p>0\).
The following conditions are equivalent:
\begin{enumerate}
\item[(i)] \SGARefTarget{sga3:xvii:proposition:4:3:1:item:i}
  \(G\) is unipotent.
\item[(ii)] \SGARefTarget{sga3:xvii:proposition:4:3:1:item:ii}
  \(G\) has a composition series whose successive quotients are
  isomorphic to \(\boldsymbol\alpha_p\) and \(\mathbb G_a\),
  taken in that order.
\item[(iii)] \SGARefTarget{sga3:xvii:proposition:4:3:1:item:iii}
  \(G\) has a characteristic composition series whose successive
  quotients are isomorphic to \((\boldsymbol\alpha_p)^r\) and
  \((\mathbb G_a)^s\), taken in that order.
\item[(iv)] \SGARefTarget{sga3:xvii:proposition:4:3:1:item:iv}
  \(G_{\overline{k}}\) contains no subgroup isomorphic to
  \(\boldsymbol\mu_p\).
\item[(v)] \SGARefTarget{sga3:xvii:proposition:4:3:1:item:v}
  \(\mathfrak g=\operatorname{Lie}G\) is unipotent (\SGARefLink{sga3:xvii:definition:3:6}{3.6}).
\item[(vi)] \SGARefTarget{sga3:xvii:proposition:4:3:1:item:vi}
  \(\mathfrak g\) is nilpotent, and the reductive part of the
  center of \(\mathfrak g\) (\SGARefLink{sga3:xvii:lemma:4-2-2}{4.2.2}) is trivial.
\item[(vi bis)]
  \SGARefTarget{sga3:xvii:proposition:4:3:1:item:vi-bis}
  \(\mathfrak g\) is nilpotent, and every subgroup of
  multiplicative type of the identity component of the center of
  \(G\) is trivial.
\item[(vi ter)]
  \SGARefTarget{sga3:xvii:proposition:4:3:1:item:vi-ter}
  \(G\) is nilpotent, and every subgroup of multiplicative type in
  the identity component of the center of \(G\) is trivial.
\end{enumerate}}
\SGARefTarget{sga3:xvii:proposition:4:3:1}

\SGAProof{Proof}
The implications
\SGARefLink{sga3:xvii:proposition:4:3:1:item:iii}{(iii)}\(\Rightarrow\)\SGARefLink{sga3:xvii:proposition:4:3:1:item:ii}{(ii)}\(\Rightarrow\)\SGARefLink{sga3:xvii:proposition:4:3:1:item:i}{(i)} are clear.  To prove
\SGARefLink{sga3:xvii:proposition:4:3:1:item:i}{(i)}\(\Rightarrow\)\SGARefLink{sga3:xvii:proposition:4:3:1:item:iii}{(iii)}, we shall need the following lemma.

\SGAStatement{\SGAInlineRefTarget{sga3:xvii:lemma:4-3-2}Lemma 4.3.2}{Let \(k\) be a field of characteristic
\(p>0\), let \(n\in\mathbb N\), and let \(k'\) be a radicial
extension of \(k\) such that \((k')^{p^n}\) is contained in \(k\).
For every \(k\)-prescheme \(X\) (respectively, every
\(k'\)-prescheme \(X'\)), denote by \(X^{(p^n)}\) (respectively,
\((X')^\phi\)) the \(k\)-prescheme obtained from \(X\)
(respectively, from \(X'\)) by the base change
\[
  F^n:k\longrightarrow k,\qquad x\longmapsto x^{p^n},
  \qquad
  \bigl(\text{respectively, }\;
  \phi:k'\longrightarrow k,\qquad x'\longmapsto(x')^{p^n}\bigr).
\]
Then, for every \(k\)-prescheme \(X\), there is a functorial
isomorphism
\[
  (X_{k'})^\phi\xrightarrow{\sim}X^{(p^n)}.
\]
Consequently, if \(X\) and \(Y\) are two \(k\)-preschemes for
which there exists a \(k'\)-isomorphism
\[
  u':X_{k'}\xrightarrow{\sim}Y_{k'},
\]
then there exists a \(k\)-isomorphism
\[
  v:X^{(p^n)}\xrightarrow{\sim}Y^{(p^n)}.
\]
If, moreover, \(X\) and \(Y\) are endowed with structures of
\(k\)-group preschemes and \(u'\) is a \(k'\)-homomorphism, then
\(v\) is a \(k\)-homomorphism.}
\SGARefTarget{sga3:xvii:lemma:4:3:2}

The lemma follows immediately from the transitivity of base
changes and from the fact that the composite
\[
  k\longrightarrow k'\xrightarrow{\phi}k
\]
is \(F^n\).

% END inputs/XVII/tex/components/17_expose_XVII_proposition4_3_1_and_lemma4_3_2_pages017_018.tex

% BEGIN inputs/XVII/tex/components/18_expose_XVII_proposition4_3_1_proof_page018.tex
% Authority: Expo17.pdf, local p.18 / combined p.980 /
% printed pp.566--567, through (iv) => (i).

\medskip\noindent\textit{Continuation of the proof of \SGARefLink{sga3:xvii:proposition:4-3-1}{4.3.1}.}

\medskip\noindent
\textit{(i)\(\Rightarrow\)\SGARefLink{sga3:xvii:proposition:4:3:1:item:iii}{(iii)}.}\quad
We argue by induction on \(\dim G\).  If \(\dim G=0\), then,
because \(G\) is connected, it is radicial, and \SGARefLink{sga3:xvii:lemma:3:9:item:i}{3.9(i)} applies.
If \(\dim G>0\), there exists an integer \(m\geq0\) such that the
quotient \(G/F_mG\) is a smooth group (Appendix II, \SGARefLink{sga3:xvii:lemma:3-1}{3.1}), plainly
connected and nonzero.  Applying
\SGARefLink{sga3:xvii:proposition:4:1:1:item:i}{4.1.1(i)}\(\Rightarrow\)\SGARefLink{sga3:xvii:lemma:3-1:item:iii}{(iii)} to this quotient, we see that \(G\)
has an algebraic subgroup \(G'\) which is characteristic and
connected, and for which the quotient
\[
  G''=G/G'
\]
is a form of \((\mathbb G_a)^r\), where \(r>0\).  By \SGARefLink{sga3:xvii:lemma:4:1:5}{4.1.5}, if
\(K\) is a perfect closure of \(k\), then
\[
  G''_K\xrightarrow{\sim}(\mathbb G_{a,K})^r.
\]
Since \(G''\) is of finite type over \(k\), there exists a finite
radicial extension \(k'\) of \(k\) such that
\[
  G''_{k'}\xrightarrow{\sim}(\mathbb G_{a,k'})^r
\]
(Exposé VI\textsubscript{B}, \S10).  Choose \(n>0\) such that
\((k')^{p^n}\subset k\).  Retaining the notation of \SGARefLink{sga3:xvii:lemma:4-3-2}{4.3.2}, we
deduce a \(k\)-isomorphism of algebraic groups
\[
  (\mathbb G_{a,k})^r
   =\bigl((\mathbb G_{a,k'})^r\bigr)^\phi
   \xrightarrow{\sim}(G'')^{(p^n)}.
\]
Now consider the Frobenius homomorphism relative to \(G''\)
(Exposé VII\textsubscript{A}, \S4):
\[
  F^n:G''\longrightarrow(G'')^{(p^n)}.
\]
Since \(G''\), and hence also \((G'')^{(p^n)}\), is smooth over
\(k\), while \(F^n\) is radicial, \(F^n\) is an epimorphism for
the fpqc topology.  Thus \((G'')^{(p^n)}\) is identified with
\(G''/F_n(G'')\), and we have shown that
\[
  G''/F_n(G'')\xrightarrow{\sim}(\mathbb G_a)^r
\]
as algebraic groups.  The inverse image \(G'_n\) of \(F_n(G'')\)
in \(G\) is a connected, characteristic subgroup of \(G\), of
dimension strictly less than that of \(G\); the induction
hypothesis therefore applies to \(G'_n\).

\medskip\noindent
\textit{(i)\(\Rightarrow\)(iv).}\quad
This follows from \SGARefLink{sga3:VIIA:2.4:item:i}{2.4(i)}.

\medskip\noindent
\textit{(iv)\(\Rightarrow\)(i).}\quad
Regard \(G\) as an extension of a smooth connected group \(G''\)
by a radicial group \(G'\) (Appendix II, \SGARefLink{sga3:xvii:lemma:3-1}{3.1}).  The group \(G'\)
is unipotent by \SGARefLink{sga3:xvii:proposition:4:2:1:item:iv}{4.2.1(iv)}\(\Leftrightarrow\)\SGARefLink{sga3:xvii:lemma:3-1:item:i}{(i)}.  It is therefore
enough to prove that \(G''\) is unipotent, and for this it is
enough to show that \(G''_{\overline{k}}\) contains no subgroup
isomorphic to \(\boldsymbol\mu_p\), by
\SGARefLink{sga3:xvii:proposition:4:1:1:item:i}{4.1.1(i)}\(\Leftrightarrow\)(iv).  If
\(G''_{\overline{k}}\) contained such a subgroup, it would lift
to \(G_{\overline{k}}\) by the result \SGARefLink{sga3:xvii:subsection:5:1}{5.1}, already used above,
whose proof is given in \S5.  This would contradict (iv).


% END inputs/XVII/tex/components/18_expose_XVII_proposition4_3_1_proof_page018.tex

% BEGIN inputs/XVII/tex/components/19_expose_XVII_proposition4_3_1_completion_through_section4_5_opening_page019.tex
% Authority: Expo17.pdf, local p.19 / combined p.981 /
% printed pp.568--569, through the statement in 4.5.

\medskip\noindent\textit{Continuation of the proof of \SGARefLink{sga3:xvii:proposition:4-3-1}{4.3.1}.}

\medskip\noindent
\textit{(i)\(\Rightarrow\)\SGARefLink{sga3:xvii:proposition:4:3:1:item:v}{(v)}.}\quad
This follows from \SGARefLink{sga3:xvii:corollary:3:7}{3.7}.

\medskip\noindent
\textit{(v)\(\Rightarrow\)\SGARefLink{sga3:xvii:proposition:4:3:1:item:vi}{(vi)}.}\quad
Indeed, since
\[
  (\operatorname{ad}X)^{p^r}
    =\operatorname{ad}\bigl(X^{(p^r)}\bigr)
  \qquad\text{(Exposé VII\textsubscript{A}, 5.2),}
\]
\(\operatorname{ad}X\) is nilpotent when \(\mathfrak g\) is
unipotent; hence \(\mathfrak g\) is nilpotent by Engel's theorem.
Moreover, if \(\mathfrak g\) is unipotent, then its center plainly
is as well, and the reductive part of that center is therefore
trivial (\SGARefLink{sga3:xvii:lemma:4-2-2}{4.2.2}).

\medskip\noindent
\textit{(vi)\(\Rightarrow\)(iv).}\quad
Indeed, if \(G_{\overline{k}}\) contained a subgroup isomorphic to
\(\boldsymbol\mu_p\), its Lie algebra would contain a nonzero
element \(X\) such that
\[
  X^{(p)}=X
  \qquad\text{(Appendix II, 2.1),}
\]
and hence \(X^{(p^r)}=X\) for every \(r>0\).
Since \(\operatorname{ad}X\) is nilpotent because \(\mathfrak g\)
is nilpotent, and
\[
  (\operatorname{ad}X)^{p^r}
    =\operatorname{ad}\bigl(X^{(p^r)}\bigr),
\]
we necessarily have \(\operatorname{ad}X=0\).  Thus \(X\) belongs
to the reductive part of the center of
\(\mathfrak g_{\overline{k}}\), contradicting (vi).

\medskip\noindent
\textit{(i)\(\Rightarrow\)(vi ter).}\quad
This follows from \SGARefLink{sga3:VIIA:2.4:item:i}{2.4(i)} and \SGARefLink{sga3:xvii:theorem:3:5:item:i}{3.5(i)}\(\Rightarrow\)(iii).

\medskip\noindent
\textit{(vi ter)\(\Rightarrow\)(vi bis).}\quad
Indeed, if \(G\) is nilpotent, then so is its subgroup
\({}_{F}G\).  Appendix II, \SGARefLink{sga3:ii:subsection:2-2}{2.2} shows that \({}_{F}G\) is
nilpotent if and only if
\[
  \operatorname{Lie}({}_{F}G)=\operatorname{Lie}G
  \qquad\text{(Exposé VII\textsubscript{A})}
\]
is nilpotent.

\medskip\noindent
\textit{(vi bis)\(\Rightarrow\)(vi).}\quad
Let \(Z\) be the identity component of the center of \(G\), and
let \(\mathfrak r\) be the reductive part of the center of
\(\mathfrak g\).  We must prove that \(\mathfrak r=0\).
It is immediate that \(\mathfrak r\) is a characteristic
\(p\)-Lie subalgebra of
\[
  \mathfrak g=\operatorname{Lie}({}_{F}G);
\]
that is, it is stable under the functor
\(\operatorname{Aut}_{p\text{-Lie}}(\mathfrak g)\).
Consequently \(\mathfrak r\) is the Lie algebra of a
characteristic radicial subgroup \(R\) of \({}_{F}G\)
(Appendix II, \SGARefLink{sga3:ii:subsection:2-2}{2.2}).  The last assertion of \SGARefLink{sga3:xvii:lemma:4-2-2}{4.2.2} and Appendix II,
\SGARefLink{sga3:xvii:proposition:2:1}{2.1} show moreover that \(R\) is a form of
\((\boldsymbol\mu_p)^r\).  Since \(R\) is characteristic in
\({}_{F}G\), which is itself characteristic in \(G\)
(Appendix II, 1), it is invariant in \(G\), and hence central
because \(G\) is connected (\SGARefLink{sga3:ix:corollary:5:5}{Exposé IX, 5.5}).  By (vi bis),
\(R=0\), and therefore \(\mathfrak r=0\).

\SGASubsection{4.4}{Étale Groups}
\SGARefTarget{sga3:xvii:subsection:4:4}

The following proposition is an easy consequence of the Sylow
theorems and of the structure of finite \(q\)-groups
(cf.\ J.-P. Serre, \emph{Corps locaux}, Chapter IX, \S1).

\SGAStatement{Proposition 4.4.1}{Let \(G\) be a finite étale
algebraic group over an algebraically closed field.  Then \(G\)
is unipotent if and only if, for every prime number \(q\) distinct
from the characteristic \(p\) of the field, \(G\) contains no
subgroup isomorphic to \(\boldsymbol\mu_q\).}
\SGARefTarget{sga3:xvii:proposition:4:4:1}

\SGASubsection{4.5}{Abelian Varieties}
\SGARefTarget{sga3:xvii:subsection:4:5}

Let \(G\) be an abelian variety over an algebraically closed field
\(k\).  The following conditions are equivalent:
\begin{enumerate}
\item[(i)] \SGARefTarget{sga3:xvii:subsection:4:5:item:i}
  \(G\) is unipotent.
\item[(ii)] \SGARefTarget{sga3:xvii:subsection:4:5:item:ii}
  \(G=0\).
\item[(iii)] \SGARefTarget{sga3:xvii:subsection:4:5:item:iii}
  There exists an integer \(n\), relatively prime to the
  characteristic \(p\) of \(k\), such that \(G\) contains no
  subgroup isomorphic to \(\boldsymbol\mu_n\).
\end{enumerate}

% END inputs/XVII/tex/components/19_expose_XVII_proposition4_3_1_completion_through_section4_5_opening_page019.tex

% BEGIN inputs/XVII/tex/components/20_expose_XVII_section4_5_proof_through_theorem5_1_1_pages019_021.tex
% Authority: Expo17.pdf, local pp.19--21 / combined pp.981--983 /
% printed pp.569--574, from the proof in 4.5 through the Galois
% cohomology notation following Theorem 5.1.1.

Indeed, if \(G\) is an abelian variety of dimension \(d\), it is
known (cf.\ S. Lang, \emph{Abelian Varieties}, Chapter IV, \S3,
\SGARefLink{sga3:xvii:section:6:site:da2ef89f51}{Theorem 6}) that the group \({}_nG(k)\), for \(n\) prime to \(p\),
is isomorphic to
\[
  (\mathbb Z/n\mathbb Z)^{2d},
\]
and hence to \((\boldsymbol\mu_n)^{2d}\).  Thus
(iii)\(\Rightarrow\)(ii), while
(ii)\(\Rightarrow\)(i)\(\Rightarrow\)(iii) are clear.

\SGASubsection{4.6}{The General Case}
\SGARefTarget{sga3:xvii:subsection:4:6}

If \(G\) and \(H\) are two algebraic groups defined over an
algebraically closed field \(k\), we shall denote by \(P(G,H)\)
the property
\[
  \text{``there is no algebraic subgroup of \(G\) isomorphic to
  \(H\).''}
\]
We then obtain the following characterizations of unipotent
groups.

\SGAStatement{\SGAInlineRefTarget{sga3:xvii:theorem:4-6-1}Theorem 4.6.1}{Let \(G\) be an algebraic group
defined over an algebraically closed field \(k\) of characteristic
\(p\).  Then:
\begin{enumerate}
\item[(i)] \SGARefTarget{sga3:xvii:theorem:4:6:1:item:i}
  If \(G\) is smooth, affine, and connected, then
  \[
    \begin{split}
    G\text{ is unipotent}
      &\Longleftrightarrow
        \text{there exists \(n>1\) such that \(P(G,\boldsymbol\mu_n)\)
        holds}\\
      &\Longleftrightarrow P(G,\mathbb G_m)\text{ holds}.
    \end{split}
  \]
\item[(ii)] \SGARefTarget{sga3:xvii:theorem:4:6:1:item:ii}
  If \(G\) is smooth and connected, then
  \[
    G\text{ is unipotent}
      \Longleftrightarrow
    \text{there exists \(n\) prime to \(p\) such that
      \(P(G,\boldsymbol\mu_n)\) holds}.
  \]
\item[(iii)] \SGARefTarget{sga3:xvii:theorem:4:6:1:item:iii}
  If \(G\) is smooth, then
  \[
    G\text{ is unipotent}
      \Longleftrightarrow
    P(G,\boldsymbol\mu_n)\text{ holds for every prime \(n\ne p\)}.
  \]
\item[(iv)] \SGARefTarget{sga3:xvii:theorem:4:6:1:item:iv}
  If \(G\) is affine and connected and \(p>0\), then
  \[
    G\text{ is unipotent}
      \Longleftrightarrow P(G,\boldsymbol\mu_p)\text{ holds}.
  \]
\item[(v)] \SGARefTarget{sga3:xvii:theorem:4:6:1:item:v}
  If \(G\) is connected and \(p>0\), then
  \[
    \begin{split}
    G\text{ is unipotent}
      \Longleftrightarrow{}&
      \text{there exists \(n\) prime to \(p\) such that}\\[-2pt]
      &P(G,\boldsymbol\mu_n)\text{ and }
        P(G,\boldsymbol\mu_p)\text{ both hold}.
    \end{split}
  \]
\item[(vi)] \SGARefTarget{sga3:xvii:theorem:4:6:1:item:vi}
  For an arbitrary algebraic group \(G\),
  \[
    G\text{ is unipotent}
      \Longleftrightarrow
    P(G,\boldsymbol\mu_n)\text{ holds for every prime \(n\)}.
  \]
\end{enumerate}}
\SGARefTarget{sga3:xvii:theorem:4:6:1}

\SGAProof{Proof}
Assertion (i) follows from \SGARefLink{sga3:xvii:proposition:4-1-1}{4.1.1}, and (iv) from \SGARefLink{sga3:xvii:proposition:4-3-1}{4.3.1}.  We shall
prove (vi); assertions (ii), (iii), and (v) are proved in the same
way and are left to the reader.

Let \(G\) be an algebraic group.  If \(G\) is unipotent, then
\(P(G,\boldsymbol\mu_n)\) holds for every \(n>1\) by 2.4(i).
Conversely, suppose that \(P(G,\boldsymbol\mu_n)\) holds for every
prime \(n\), and let us prove that \(G\) is unipotent.  Write
\(G^0\) for the identity component of \(G\).  If \(G^0\) is
smooth, a classical theorem of
Chevalley\begingroup\renewcommand{\thefootnote}{*}\footnote[0]{%
\emph{Séminaire Bourbaki}, no.\ 145, 1956/57.}\endgroup
shows that \(G^0\) is an extension
of an abelian variety \(A\) by a smooth connected affine group
\(L\).  If \(G\) is not smooth---which entails \(p>0\)
(Exposé VI\textsubscript{B}, \SGARefLink{cor:sga3-VIB-1.6.1}{1.6.1})---there is an integer \(n>0\)
such that
\[
  G''=G^0/{}_{F^n}(G)
\]
is smooth (Appendix II, \SGARefLink{sga3:xvii:lemma:3-1}{3.1}).  Thus \(G''\) is an extension of an
abelian variety \(A\) by a smooth connected linear group \(L''\).
Let \(L\) be the inverse image of \(L''\) in \(G^0\).  It is still
affine and connected, since \({}_{F^n}(G)\) is radicial.  In every
case, therefore, \(G\) has a composition series
\[
  0\subset L\subset G^0\subset G
\]
such that \(L\) is affine and connected, \(G^0/L=A\) is an
abelian variety, and \(G/G^0\) is an étale group.

If \(P(G,\boldsymbol\mu_n)\) holds, then a fortiori
\(P(L,\boldsymbol\mu_n)\) holds, and hence \(L\) is unipotent
(\SGARefLink{sga3:xvii:proposition:4-1-1}{4.1.1} and \SGARefLink{sga3:xvii:proposition:4-3-1}{4.3.1}).  If \(A\ne0\), there are a prime number \(n\)
and a subgroup of \(A\) isomorphic to \(\boldsymbol\mu_n\) (\SGARefLink{sga3:xvii:subsection:4:5}{4.5}).
By \SGARefLink{sga3:xvii:subsection:5:1}{5.1} below, this subgroup lifts to \(G\), contradicting
\(P(G,\boldsymbol\mu_n)\).  Hence \(A=0\).  Finally, if
\(G/G^0\) is not unipotent, there are an integer \(q\) and a
subgroup of \(G/G^0\) isomorphic to \(\boldsymbol\mu_q\) (\SGARefLink{sga3:xvii:proposition:4:4:1}{4.4.1}).
As above, it follows that \(P(G,\boldsymbol\mu_q)\) does not hold.
Thus \(G/G^0\) is unipotent, and consequently so is \(G\).

\SGASection{5}{Extensions of Groups of Multiplicative Type by
Unipotent Groups}
\SGARefTarget{sga3:xvii:section:5}

\SGASubsection{5.1}{Statement of the Theorem}
\SGARefTarget{sga3:xvii:subsection:5:1}

\SGAStatement{\SGAInlineRefTarget{sga3:xvii:definition:5-1-0}Definition 5.1.0}{Let \(k\) be a field and \(G\) an
algebraic \(k\)-group.  Following the terminology introduced by
Rosenlicht (\emph{Questions of rationality for solvable algebraic
groups over non perfect fields}, Annali di Matematica 61 (1963)),
we shall say that \(G\) is \emph{\(k\)-solvable} if it satisfies
the following equivalent conditions:
\begin{enumerate}
\item[(i)] \SGARefTarget{sga3:xvii:definition:5:1:0:item:i}
  \(G\) has a composition series whose successive quotients are
  isomorphic to \(\mathbb G_a\).
\item[(ii)] \SGARefTarget{sga3:xvii:definition:5:1:0:item:ii}
  \(G\) has a characteristic composition series \((G_i)\) such
  that the successive quotients \(G_i/G_{i+1}\) are commutative
  and themselves have composition series whose successive
  quotients are isomorphic to \(\mathbb G_a\).
\end{enumerate}}
\SGARefTarget{sga3:xvii:definition:5:1:0}

The implication \SGARefLink{sga3:xvii:definition:5:1:0:item:i}{(i)}\(\Rightarrow\)\SGARefLink{sga3:xvii:definition:5:1:0:item:ii}{(ii)} is proved in
\emph{loc.\ cit.}  In fact, for the composition series \((G_i)\)
one may take the series introduced in the proof of \SGARefLink{sga3:xvii:lemma:3:9:item:ii}{3.9(ii)}.

\SGAStatement{\SGAInlineRefTarget{sga3:xvii:theorem:5-1-1}Theorem 5.1.1}{Let \(k\) be a field, \(U\) a
unipotent algebraic \(k\)-group, \(H\) a \(k\)-group of
multiplicative type, and \(E\) an algebraic \(k\)-group which is
an extension of \(H\) by \(U\), so that there is an exact sequence
\[
  1\longrightarrow U\longrightarrow E\longrightarrow H
    \longrightarrow1.
\]
Then:
\begin{enumerate}
\item[(i)] \SGARefTarget{sga3:xvii:theorem:5:1:1:item:i}
  The extension \(E\) is trivial in each of the following cases:
  \begin{enumerate}
  \item[(a)] \SGARefTarget{sga3:xvii:theorem:5:1:1:item:i:a}
    \(k\) is algebraically closed;
  \item[(b)] \SGARefTarget{sga3:xvii:theorem:5:1:1:item:i:b}
    \(k\) is perfect and one of the groups \(U\) or \(H\) is
    connected;
  \item[(c)] \SGARefTarget{sga3:xvii:theorem:5:1:1:item:i:c}
    \(U\) is \(k\)-solvable;
  \item[(d)] \SGARefTarget{sga3:xvii:theorem:5:1:1:item:i:d}
    \(U\) is smooth and \(H\) is connected.
  \end{enumerate}
\item[(ii)] \SGARefTarget{sga3:xvii:theorem:5:1:1:item:ii}
  If \(H'\) and \(H''\) are two lifts of \(H\) in \(E\), then
  they are conjugate by an element of \(U(k)\) in each of the
  following cases:
  \begin{enumerate}
  \item[(a)] \SGARefTarget{sga3:xvii:theorem:5:1:1:item:ii:a}
    \(k\) is algebraically closed and \(H\) is smooth;
  \item[(b)] \SGARefTarget{sga3:xvii:theorem:5:1:1:item:ii:b}
    \(k\) is algebraically closed and \(U\) is smooth.
  \end{enumerate}
\end{enumerate}
(During the proof we shall indicate other cases in which the
conclusion of \SGARefLink{sga3:xvii:theorem:5:1:1:item:ii}{(ii)} is true without assuming that \(k\) is
algebraically closed.)}
\SGARefTarget{sga3:xvii:theorem:5:1:1}

If \(U\) is an algebraic group (respectively, a commutative
algebraic group) defined over a field \(k\), we denote by
\(H^1(k,U)\) (respectively, by \(H^i(k,U)\), \(i\geq0\)) the
pointed set (respectively, the \(i\)-th group) of Galois
cohomology of \(k\) with values in \(U\); see J.-P. Serre,
\emph{Cohomologie galoisienne}, Lecture Notes in Mathematics
no.\ 5, Springer.

% END inputs/XVII/tex/components/20_expose_XVII_section4_5_proof_through_theorem5_1_1_pages019_021.tex

% BEGIN inputs/XVII/tex/components/21_expose_XVII_subsection5_2_through_lemma5_2_4_pages022_023.tex
% Authority: Expo17.pdf, local pp.22--23 / combined pp.984--985 /
% printed pp.575--578, through the proof of Lemma 5.2.4(b).

If \(S\) is a prescheme, \(H\) a commutative \(S\)-group
prescheme, and \(G\) an \(S\)-group prescheme acting on \(H\), we
denote by \(H^i(G,H)\)\footnote{Editor's note: \(H_0^i\) has been
replaced by \(H^i\), in agreement with the notation of Appendix I
and Exposé I.} the \(i\)-th Hochschild cohomology group of \(G\)
with values in \(H\) (Appendix I, 1).

To prove \SGARefLink{sga3:xvii:theorem:5-1-1}{5.1.1}, we shall proceed in several steps.

\SGASubsection{5.2}{\texorpdfstring{Proof of \SGARefLink{sga3:xvii:theorem:5:1:1:item:i}{5.1.1(i)} and \SGARefLink{sga3:xvii:theorem:5:1:1:item:ii}{(ii)}
when \(U\) is smooth and \(H\) is étale}{Proof of \SGARefLink{sga3:xvii:theorem:5:1:1:item:i}{5.1.1(i)} and
\SGARefLink{sga3:xvii:theorem:5:1:1:item:ii}{(ii)} when U is smooth and H is etale}}
\SGARefTarget{sga3:xvii:subsection:5:2}

\SGAStatement{Lemma 5.2.1}{With the notation of \SGARefLink{sga3:xvii:theorem:5-1-1}{5.1.1}, suppose
that \(H\) is étale.  The canonical morphism \(E\to H\) has a
section \(s:H\to E\), defined over \(k\), in each of the following
cases:
\begin{enumerate}
\item[(a)] \SGARefTarget{sga3:xvii:lemma:5:2:1:item:a}
  \(k\) is algebraically closed;
\item[(b)] \SGARefTarget{sga3:xvii:lemma:5:2:1:item:b}
  \(k\) is perfect and \(U\) is smooth and connected;
\item[(c)] \SGARefTarget{sga3:xvii:lemma:5:2:1:item:c}
  \(U\) is \(k\)-solvable.
\end{enumerate}}
\SGARefTarget{sga3:xvii:lemma:5:2:1}

Assertion \SGARefLink{sga3:xvii:appendix:A}{(a)} follows simply from the surjectivity of
\(E(k)\to H(k)\).  We have \SGARefLink{sga3:xvii:section:b}{(b)}\(\Rightarrow\)\SGARefLink{sga3:xvii:appendix:C}{(c)} by \SGARefLink{sga3:xvii:lemma:4:1:2:item:b}{4.1.2(b)}
and \SGARefLink{sga3:xvii:lemma:2:3:bis}{2.3 bis}, so it is enough to treat case \SGARefLink{sga3:xvii:appendix:C}{(c)}.  Let \(x\) be a
point of \(H\).  It is an open and closed subset of \(H\), and
the induced subscheme is isomorphic to \(\operatorname{Spec}K\),
where \(K\) is a finite separable extension of \(k\).  Let \(X\)
be the inverse image in \(E\) of the scheme \(x\).  The
\(K\)-scheme \(X\) is a principal homogeneous \(K\)-space under
\(U_K\).  But \(U_K\) has a composition series whose successive
quotients are isomorphic to \((\mathbb G_a)_K\); hence
\[
  H^1(K,U_K)=0
\]
(J.-P. Serre, \emph{Cohomologie galoisienne}, Chapter III,
\SGARefLink{sga3:xvii:section:6}{Proposition 6}).  Consequently \(X\) is trivial and has a
\(K\)-rational point.  Thus, over every point \(x\) of \(H\), we
obtain a \(k\)-section of \(E\to H\), and these give a section
\(H\to E\).

\SGAStatement{\SGAInlineRefTarget{sga3:xvii:lemma:5-2-3}Lemma 5.2.3}{\footnote{Editor's note: there is no
item numbered 5.2.2.}With the notation of \SGARefLink{sga3:xvii:theorem:5-1-1}{5.1.1}, suppose that
\(H\) is étale.  Then:
\begin{enumerate}
\item[(i)] \SGARefTarget{sga3:xvii:lemma:5:2:3:item:i}
  The extension \(E\) is trivial in each of the following cases:
  \begin{enumerate}
  \item[(a)] \SGARefTarget{sga3:xvii:lemma:5:2:3:item:i:a}
    \(U\) is commutative and \(E\to H\) has a section;
  \item[(b)] \SGARefTarget{sga3:xvii:lemma:5:2:3:item:i:b}
    \(k\) is algebraically closed;
  \item[(c)] \SGARefTarget{sga3:xvii:lemma:5:2:3:item:i:c}
    \(k\) is perfect and \(U\) is connected;
  \item[(d)] \SGARefTarget{sga3:xvii:lemma:5:2:3:item:i:d}
    \(U\) is \(k\)-solvable.
  \end{enumerate}
  Moreover, in each of the preceding cases, two lifts \(H'\) and
  \(H''\) of \(H\) in \(E\) are conjugate by an element of
  \(U(k)\).
\item[(ii)] \SGARefTarget{sga3:xvii:lemma:5:2:3:item:ii}
  Let \(R\) be the \(k\)-functor
  \[
    (\operatorname{Sch}/k)^\circ\longrightarrow\operatorname{Ens}
  \]
  for which, for every \(k\)-prescheme \(S\), \(R(S)\) is the set
  of lifts of \(H_S\) in \(E_S\).  If \(U\) is commutative, then
  \(R\) is representable by a nonempty affine \(k\)-scheme.  The
  group \(U\) acts on \(R\) by inner automorphisms, and this action
  makes \(R\) a homogeneous space under \(U\) for the fpqc
  topology (cf.\ Exposé IV).
\end{enumerate}}
\SGARefTarget{sga3:xvii:lemma:5:2:3}

\SGAProof{Proof of \SGARefLink{sga3:xvii:lemma:5:2:3:item:i}{(i)}}
We shall reduce cases \SGARefLink{sga3:xvii:section:b}{(b)}, \SGARefLink{sga3:xvii:appendix:C}{(c)}, and \SGARefLink{sga3:xvii:lemma:5:2:3:item:i:d}{(d)} to case \SGARefLink{sga3:xvii:appendix:A}{(a)}.

\medskip\noindent\textit{Case \SGARefLink{sga3:xvii:section:b}{(b)}.}\quad
Since \(U\) has a characteristic composition series with
commutative successive quotients
(\SGARefLink{sga3:xvii:theorem:3-5}{3.5}, \SGARefLink{sga3:xvii:theorem:3:5:item:i}{(i)}\(\Leftrightarrow\)\SGARefLink{sga3:xvii:theorem:3:5:item:iv}{(iv)}), we immediately reduce to the
case in which \(U\) is commutative.  Moreover, \(E\to H\) has a
section by \SGARefLink{sga3:xvii:lemma:5:2:1:item:a}{5.2.1(a)}, so we are reduced to case \SGARefLink{sga3:xvii:appendix:A}{(a)}.

\medskip\noindent\textit{Case \SGARefLink{sga3:xvii:lemma:5:2:3:item:i:d}{(d)}.}\quad
We argue in the same way, using \SGARefLink{sga3:xvii:definition:5:1:0:item:ii}{5.1.0(ii)} and \SGARefLink{sga3:xvii:lemma:5:2:1:item:c}{5.2.1(c)}.

\medskip\noindent\textit{Case \SGARefLink{sga3:xvii:appendix:C}{(c)}.}\quad
Let \(E_{\mathrm{red}}\) be the reduced subgroup associated with
\(E\).  Since \(H\) is smooth, \(E_{\mathrm{red}}\) is an
extension of \(H\) by
\[
  U_0=U\cap E_{\mathrm{red}},
\]
and we may replace \(E\) by \(E_{\mathrm{red}}\).  The group \(U\),
and therefore \(U_0\), is connected, and \(U_0\) is plainly the
identity component of \(E_{\mathrm{red}}\); in particular it is
smooth.  Thus \(U_0\) is smooth and connected and \(k\) is
perfect, so \(U_0\) is \(k\)-solvable by \SGARefLink{sga3:xvii:remark:4:1:4:item:b}{4.1.4(b)}, reducing us to
case \SGARefLink{sga3:xvii:lemma:5:2:3:item:i:d}{(d)}.

The preceding reductions show that in cases \SGARefLink{sga3:xvii:section:b}{(b)}, \SGARefLink{sga3:xvii:appendix:C}{(c)}, and \SGARefLink{sga3:xvii:lemma:5:2:3:item:i:d}{(d)} we
may suppose that \(U\) has a characteristic composition series
\((U_i)\) such that \(U_i/U_{i+1}\) is commutative and the maps
\[
  U_i(k)\longrightarrow (U_i/U_{i+1})(k)
\]
are surjective.  (In cases \SGARefLink{sga3:xvii:appendix:C}{(c)} and \SGARefLink{sga3:xvii:lemma:5:2:3:item:i:d}{(d)}, the latter assertion
follows from \(H^1(k,\mathbb G_a)=0\).)  An immediate dévissage
then shows that it is enough to prove the conjugacy of two lifts
\(H'\) and \(H''\) in \(E\) when \(U\) is commutative.  In short,
it is enough to prove \SGARefLink{sga3:xvii:lemma:5:2:3:item:i}{(i)}\SGARefLink{sga3:xvii:lemma:5:2:4:item:a}{(a)}.  In that case the extension \(E\)
is trivial provided that \(H^2(H,U)=0\) (Appendix I, \SGARefLink{sga3:i:definition:3-1}{3.1}), and
\(H'\) and \(H''\) are conjugate provided that \(H^1(H,U)=0\).
We now use the following lemma.

\SGAStatement{Lemma 5.2.4}{Let \(S\) be a prescheme, \(U\) a
commutative \(S\)-group prescheme, and \(H\) an étale finite
\(S\)-group prescheme of rank \(n\) acting on \(U\).  Then the
groups \(H^i(H,U)\), for \(i>0\), are annihilated by \(n\) in
each of the following cases:
\begin{enumerate}
\item[(a)] \SGARefTarget{sga3:xvii:lemma:5:2:4:item:a}
  \(H\) is a constant \(S\)-group (\SGARefLink{sga3:i:subsection:4-1}{Exposé I, 4.1});
\item[(b)] \SGARefTarget{sga3:xvii:lemma:5:2:4:item:b}
  \(S\) is the spectrum of a field.
\end{enumerate}}
\SGARefTarget{sga3:xvii:lemma:5:2:4}

\SGAProof{Proof of \SGARefLink{sga3:xvii:appendix:A}{(a)}}
By hypothesis, \(H\) is the constant group associated with an
ordinary group \(\{H\}\) of order \(n\).  It is
clear\footnote{Editor's note: this is Proposition III.6.4.2 of
M. Demazure and P. Gabriel, \emph{Groupes algébriques I},
Masson \& North-Holland (1970).} that \(H^i(H,U)\) is isomorphic
to the \(i\)-th cohomology group
\[
  H^i(\{H\},U(S))
\]
of the group \(\{H\}\) with values in the ordinary group \(U(S)\).
It is classical that these groups are annihilated by \(n\)
(J.-P. Serre, \emph{Corps locaux}, Chapter VIII, \SGARefLink{sga3:viii:section:4}{Proposition 4},
Corollary 1).

\SGAProof{Proof of \SGARefLink{sga3:xvii:section:b}{(b)}}
Let \(x\in H^i(H,U)\), where \(i>0\), and represent it by an
\(i\)-cocycle
\[
  f:H^{(i)}\longrightarrow U,
\]
where \(H^{(i)}\) denotes the product over \(k\) of \(i\) copies
of \(H\).  If \(K\) is a finite Galois extension of \(k\) which
splits \(H\), assertion \SGARefLink{sga3:xvii:lemma:5:2:4:item:a}{(a)} shows that \(nf_K\) is a coboundary.
More precisely, a direct computation shows that, if one defines
the \(K\)-morphism
\[
  F_K:H_K^{(i-1)}\longrightarrow U_K
\]
by
\[
  F_K(h_1,\ldots,h_{i-1})
    =\sum_{h\in H(K)}
      f_K(h_1,\ldots,h_{i-1},h),
\]
then, up to sign,
\[
  d(F_K)=nf_K,
  \qquad d\text{ denoting the coboundary operator}.
\]
An immediate Galois argument shows that \(F_K\) descends to a
\(k\)-morphism \(F:H^{(i-1)}\to U\), and consequently
\[
  d(F)=nf.
\]

% END inputs/XVII/tex/components/21_expose_XVII_subsection5_2_through_lemma5_2_4_pages022_023.tex

% BEGIN inputs/XVII/tex/components/22_expose_XVII_corollary5_2_4_bis_through_proposition5_3_1_opening_page024.tex
% Authority: Expo17.pdf, local p.24 / combined p.986 /
% printed pp.579--580, from Corollary 5.2.4 bis through the
% opening inverse system in the proof of Proposition 5.3.1.

\SGAStatement{Corollary 5.2.4 bis}{With the notation of \SGARefLink{sga3:xvii:lemma:5:2:4}{5.2.4},
suppose in addition that \(U\) is flat and of finite presentation
over \(S\), with unipotent fibres, and that \(n\) is prime to the
residue characteristics of \(S\).  Then, in cases (a) and (b)
above,
\[
  H^i(H,U)=0\qquad(i>0).
\]}
\SGARefTarget{sga3:xvii:corollary:5:2:4:bis}

\SGAProof{Proof}
It is enough to show that the \(n\)-th power map on \(U\) is an
isomorphism, for it will then follow that multiplication by \(n\)
on \(H^i(H,U)\) is both an isomorphism and the zero morphism, and
hence that \(H^i(H,U)=0\).  Under the hypotheses made on \(U\), it
is enough to check that the \(n\)-th power map is an isomorphism
on the fibres of \(U\) (EGA IV, 17.9.5).  We are thus reduced to
the case in which \(S\) is the spectrum of a field \(k\) of
characteristic \(p\).  Since \((n,p)=1\), the \(n\)-th power map
on \(U\) is étale (Exposé VII); it is also a monomorphism (2.4(i)),
and hence an open immersion (EGA IV, 17.9.1).  This already shows
that its restriction to the identity component \(U^0\) is an
isomorphism.  Since \(U/U^0\) is a finite \(p\)-group, the proof
is complete.

This completes the proof of \SGARefLink{sga3:xvii:lemma:5:2:3:item:i}{5.2.3(i)}\SGARefLink{sga3:xvii:appendix:A}{(a)}, since an étale group of
multiplicative type \(H\) has order prime to \(p\).

\SGAProof{Proof of \SGARefLink{sga3:xvii:lemma:5:2:3:item:ii}{5.2.3(ii)}}
It is clear that \(R\) is a sheaf for the fpqc topology.
Moreover, in view of descent for affine schemes, the assertions
of \SGARefLink{sga3:xvii:lemma:5:2:3:item:ii}{5.2.3(ii)} are local for the fpqc topology.  We may therefore
suppose that \(k\) is algebraically closed.  The group \(H\) is
then split and the extension \(E\) is trivial by (i)\SGARefLink{sga3:xvii:section:b}{(b)}; let
\(H'\) be a lift of \(H\) in \(E\).  For every \(k\)-prescheme
\(S\) and every lift \(H''\) of \(H_S\) in \(E_S\), the groups
\(H'_S\) and \(H''\) are conjugate by an element of \(U(S)\),
since
\[
  H^1(H_S,U_S)=0
\]
by \SGARefLink{sga3:xvii:corollary:5:2:4:bis}{5.2.4 bis}.

Let \(U^{H'}\) be the sheaf of invariants of \(U\) under \(H'\).
It is representable by an algebraic subgroup of \(U\)
(\SGARefLink{sga3:viii:examples:6:5:item:d:01}{Exposé VIII, 6.5(d)}).  The preceding remarks show that the
\(k\)-morphism
\[
  U\longrightarrow R,
  \qquad
  u\longmapsto \operatorname{int}(u)H'
  \quad (u\in U(S))
\]
induces a \(k\)-isomorphism
\[
  U/U^{H'}\xrightarrow{\ \sim\ }R.
\]
This proves that \(R\) is representable and affine (\SGARefLink{sga3:xvii:proposition:2:1}{2.1}).

\SGAStatement{Remark 5.2.5}{The assertions of \SGARefLink{sga3:xvii:lemma:5:2:3:item:ii}{5.2.3(ii)} can also
be proved when \(U\) is not commutative, but we shall not need
this in order to prove \SGARefLink{sga3:xvii:theorem:5-1-1}{5.1.1}.}
\SGARefTarget{sga3:xvii:remark:5:2:5}

\SGASubsection{5.3}{\texorpdfstring{The Case Where \(H\) Is
Smooth}{The Case Where H Is Smooth}}
\SGARefTarget{sga3:xvii:subsection:5:3}

\SGAStatement{Proposition 5.3.1}{The assertions of \SGARefLink{sga3:xvii:lemma:5:2:3:item:i}{5.2.3(i)}
remain true when the hypothesis ``\(H\) is étale'' is replaced by
``\(H\) is smooth.''}
\SGARefTarget{sga3:xvii:proposition:5:3:1}

Proceeding as in the proof of \SGARefLink{sga3:xvii:lemma:5:2:3:item:i}{5.2.3(i)}, we reduce to the case in
which \(U\) is also commutative.

Let \(N'\) be the set of positive integers prime to \(p\), ordered
by divisibility.  For every \(n\in N'\), the group \({}_nH\) is
étale, and the family of groups \({}_nH\), \(n\in N'\), is
schematically dense in \(H\), since \(H\) is smooth
(\SGARefLink{sga3:ix:remark:4:10}{Exposé IX, 4.10}).  Let \(E_n\) be the inverse image of
\({}_nH\) in \(E\), so that \(E_n\) is an extension of
\({}_nH\) by \(U\), and let \(R_n\) be the \(k\)-functor of
lifts of \({}_nH\) in \(E_n\) (cf.\ \SGARefLink{sga3:xvii:lemma:5:2:3:item:ii}{5.2.3(ii)}).  If \(n\)
divides \(m\), there is plainly a natural \(k\)-morphism
\[
  R_m\longrightarrow R_n,
\]
so the \(R_n\) form an inverse system of \(k\)-functors.

% END inputs/XVII/tex/components/22_expose_XVII_corollary5_2_4_bis_through_proposition5_3_1_opening_page024.tex

% BEGIN inputs/XVII/tex/components/23_expose_XVII_proposition5_3_1_completion_and_proposition5_4_1_opening_page025.tex
% Authority: Expo17.pdf, local p.25 / combined p.987 /
% printed pp.581--582, completing Proposition 5.3.1 and beginning
% Proposition 5.4.1 through the Frobenius exact sequence.

By \SGARefLink{sga3:xvii:lemma:5:2:3:item:ii}{5.2.3(ii)}, \(R_n\) is representable by a nonempty affine
\(k\)-scheme.  Since a filtered direct limit of nonzero rings is
nonzero, the functor
\[
  R=\varprojlim_n R_n
\]
is representable by a nonempty affine \(k\)-scheme
(EGA IV, 8 and 1.9.1).  There are therefore an extension \(K\) of
\(k\) and a point \(u\in R(K)\).  The image \(u_n\) of \(u\) in
\(R_n(K)\) corresponds to a lift \(H'_n\) of \(({}_nH)_K\) in
\((E_n)_K\).  By construction,
\[
  H'_n={}_n(H'_m)\qquad\text{if \(n\) divides \(m\)}.
\]
Put
\[
  U_n=(U_K)^{H'_n}.
\]
The choice of \(H'_n\) identifies \((R_n)_K\) with \(U_K/U_n\).
The family of groups \(H'_n\) is filtered increasing, and hence
the family \(U_n\) is filtered decreasing.  It is therefore
stationary for sufficiently large \(n\), since \(U_K\) is
noetherian.  It follows that the family \((R_n)_K\) is stationary,
and consequently so is the family \(R_n\).  In short,
\[
  R_n=R
\]
for sufficiently large \(n\).

Under the hypotheses in force, \SGARefLink{sga3:xvii:lemma:5:2:3:item:i}{5.2.3(i)} shows that \(R_n(k)\) is
nonempty.  We may therefore find a coherent system of lifts
\(H'_n\) of \({}_nH\), for \(n\in N'\).  Let \(H'\) be the
smallest closed subscheme of \(E\) containing \(H'_n\) for every
\(n\) (Exposé VI\textsubscript{B}, \S7).  The argument of
\SGARefLink{sga3:xv:lemma:4:6}{Exposé XV, 4.6} shows that \(H'\) is a smooth commutative
algebraic group whose formation commutes with extension of the
base field.  To prove that \(H'\) is a lift of \(H\) in \(E\), we
may therefore suppose that \(k\) is algebraically closed.
According to BIBLE 4, Theorem 4, \(H'\) is then the direct product
of a smooth group \(M\) of multiplicative type and a unipotent
group \(V\).  The groups \(H'_n\) are necessarily contained in
\(M\) by \SGARefLink{sga3:xvii:proposition:2:4}{2.4}; the definition of \(H'\) then gives \(H'=M\).
Thus \(H'\) is of multiplicative type, and hence
\[
  H'\cap U=0.
\]
The morphism \(H'\to H\) is consequently a monomorphism, and the
density theorem of \SGARefLink{sga3:ix:remark:4:10}{Exposé IX, 4.10} shows that it is an
epimorphism.  It is therefore an isomorphism.

Now let \(H'\) and \(H''\) be two lifts of \(H\) in \(E\).  For
each \(n\in N'\), let
\[
  T_n=\operatorname{Transp}({}_nH',{}_nH'')
\]
be the transporter of \({}_nH'\) into \({}_nH''\).  It is
representable by a closed subscheme of \(U\)
(\SGARefLink{sga3:viii:examples:6:5:item:e:01}{Exposé VIII, 6.5(e)}).  The \(T_n\) form a filtered decreasing
family of closed subschemes of \(U\), and they are nonempty by
\SGARefLink{sga3:xvii:lemma:5:2:3:item:i}{5.2.3(i)}\SGARefLink{sga3:xvii:appendix:A}{(a)}.  Let \(T\) be their stationary value.  Under the
hypotheses of \SGARefLink{sga3:xvii:lemma:5:2:3:item:i}{5.2.3(i)}, \(T_n(k)\) is nonempty.  Thus there is an
element \(u\in U(k)\) such that
\[
  {}_nH''=\operatorname{int}(u)({}_nH')
  \qquad\text{for every \(n\in N'\)}.
\]
It follows that
\[
  H''=\operatorname{int}(u)H'
\]
by \SGARefLink{sga3:ix:corollary:4:8:item:b}{Exposé IX, 4.8(b)}.

\SGASubsection{5.4}{\texorpdfstring{The Case Where \(U\) Is
Radicial}{The Case Where U Is Radicial}}
\SGARefTarget{sga3:xvii:subsection:5:4}

\SGAStatement{Proposition 5.4.1}{If \(k\) is a perfect field of
characteristic \(p>0\), and if \(U\) is radicial, then the
extension \(E\) of \SGARefLink{sga3:xvii:theorem:5-1-1}{5.1.1} is trivial.}
\SGARefTarget{sga3:xvii:proposition:5:4:1}

Using a characteristic composition series of \(U\), we may
restrict ourselves to the case
\[
  U=(\boldsymbol\alpha_p)^r
\]
by \SGARefLink{sga3:xvii:lemma:3-9}{3.9}.

It follows from Appendix II, \SGARefLink{sga3:ii:subsection:2-2}{2.2} and \SGARefLink{sga3:xvii:proposition:2:1}{2.1} that there are
isomorphisms of \(k\)-functors
\[
  \operatorname{Aut}_{k\text{-}\mathrm{gr}}
    \bigl((\boldsymbol\alpha_p)^r\bigr)
  \xrightarrow{\ \sim\ }
  \operatorname{Aut}_{p\text{-}\mathrm{Lie}}
    \bigl(\operatorname{Lie}((\boldsymbol\alpha_p)^r)\bigr)
  \xrightarrow{\ \sim\ }
  \operatorname{GL}
    \bigl(\operatorname{Lie}((\boldsymbol\alpha_p)^r)\bigr).
\]
Let \(V\) be a \(k\)-vector space of rank \(r\), let
\(\mathbf W(V)\) be the corresponding vector \(k\)-group
(\SGARefLink{sga3:i:subsection:4-6}{Exposé I, 4.6}), and identify \((\boldsymbol\alpha_p)^r\) with
\({}_F V\).  The preceding remarks show that the action of \(H\)
on \({}_F V\) defined by the extension \(E\) comes from a linear
representation \(v\) of \(H\) on \(V\).  Consider the exact
sequence
\[
  \tag{*}\label{sga3:xvii:proposition:5-4-1:equation:star}
  0\longrightarrow {}_F V\longrightarrow V
    \xrightarrow{F}V^{(p)}\longrightarrow0,
\]
where \(F\) is the Frobenius morphism.

% END inputs/XVII/tex/components/23_expose_XVII_proposition5_3_1_completion_and_proposition5_4_1_opening_page025.tex

% BEGIN inputs/XVII/tex/components/24_expose_XVII_proposition5_4_1_completion_through_proposition5_6_1_statement_pages026_028.tex
% Authority: Expo17.pdf, local pp.26--28 / combined pp.988--990 /
% printed pp.583--587, completing Proposition 5.4.1, proving the
% relevant cases of Theorem 5.1.1(i), and giving the complete
% statement of Proposition 5.6.1.  The next cursor is its proof.

The sequence \((*)\) is an exact sequence of \(H\)-groups,
provided that \(H\) acts on \(V^{(p)}\) through the linear
representation
\[
  H\xrightarrow{v}\operatorname{GL}(V)
    \xrightarrow{F}\operatorname{GL}(V^{(p)}).
\]
Since \(k\) is perfect, the morphism \(F:V\to V^{(p)}\) induces
a surjective map
\[
  V(k)\longrightarrow V^{(p)}(k).
\]
It follows from Appendix I, \SGARefLink{sga3:i:subsection:2-1}{2.1} that the exact sequence \((*)\)
defines an exact sequence
\[
  \tag{**}\label{sga3:xvii:proposition:5-4-1:equation:double-star}
  H^1(H,V^{(p)})\longrightarrow
  \operatorname{Ext}_{\mathrm{alg}}(H,{}_F V)
  \longrightarrow
  \operatorname{Ext}_{\mathrm{alg}}(H,V).
\]

We shall show that
\[
  \operatorname{Ext}_{\mathrm{alg}}(H,V)=0.
\]
Let \(E'\) be an algebraic group which is an extension of \(H\)
by \(V\):
\[
  1\longrightarrow V\longrightarrow E'
    \xrightarrow{h}H\longrightarrow1.
\]
The scheme \(E'\) is a torsor over \(H\) under
\(\mathbb G_a^r\), and hence defines an element of
\[
  H^1(H,\mathcal O_H^r)
\]
in coherent-sheaf cohomology.  Since \(H\) is affine,
\[
  H^1(H,\mathcal O_H^r)=0
\]
by EGA III, \S1.  Thus \(E'\to H\) has a section.  Consequently,
\(\operatorname{Ext}_{\mathrm{alg}}(H,V)\) is isomorphic to
\(H^2(H,V)\) by Appendix I, \SGARefLink{sga3:i:definition:3-1}{3.1}.  But
\[
  H^i(H,V)=H^i(H,\mathbf W(V))=0\qquad(i>0)
\]
by \SGARefLink{sga3:ix:theorem:3:1}{Exposé IX, 3.1}.  The exact sequence \((**)\) now gives
\[
  \operatorname{Ext}_{\mathrm{alg}}(H,{}_F V)=0,
\]
and hence \(E\) is a trivial extension.

\SGASubsection{5.5}{Proof of \SGARefLink{sga3:xvii:theorem:5:1:1:item:i}{5.1.1(i)}}
\SGARefTarget{sga3:xvii:subsection:5:5}

If \(k\) has characteristic \(0\), then \(U\) is \(k\)-solvable
by \SGARefLink{sga3:xvii:corollary:4:1:3}{4.1.3} and \(H\) is smooth.  It follows from \SGARefLink{sga3:xvii:proposition:5:3:1}{5.3.1} and
\SGARefLink{sga3:xvii:lemma:5:2:3:item:i:d}{5.2.3(d)} that \(E\) is a trivial extension.  The same argument
shows that two lifts of \(H\) in \(E\) are conjugate by an
element of \(U(k)\).

Henceforth we suppose that \(k\) is a field of characteristic
\(p>0\).

\SGAProof{Proof of \SGARefLink{sga3:xvii:lemma:5:5:1:item:i}{(i)}\SGARefLink{sga3:xvii:section:b}{(b)}: the case where \(k\) is perfect and
\(U\) is connected}
We shall reduce to the case in which \(U\) is radicial.  Since
\(k\) is perfect, \(H_{\mathrm{red}}\) is smooth; let \(E'\) be
its inverse image in \(E\).  By \SGARefLink{sga3:xvii:proposition:5:3:1}{5.3.1} and \SGARefLink{sga3:xvii:lemma:5:2:3:item:i}{5.2.3(i)}\SGARefLink{sga3:xvii:appendix:C}{(c)}, the
extension
\[
  1\longrightarrow U\longrightarrow E'
    \longrightarrow H_{\mathrm{red}}\longrightarrow1
\]
is trivial.  Let \(H_1\) be a lift of \(H_{\mathrm{red}}\) in
\(E\).  By Appendix II, \SGARefLink{sga3:xvii:lemma:3-1}{3.1}, there is an integer \(n>0\) such
that
\[
  E^{(n)}=E/{}_{F^n}(E)
\]
is smooth.  Let \(E''\) be the algebraic subgroup of \(E\)
generated by \(H_1\) and \({}_{F^n}(E)\); equivalently, \(E''\)
is the inverse image in \(E\) of the image of \(H_1\) in
\(E^{(n)}\).  Let \(H''\) be the image of \(E''\) in \(H\).  We
claim that \(H''=H\).

Indeed, let \(R\) be the image of \({}_{F^n}(E)\) in \(H\), so
that \(H''\) is generated by \(R\) and \(H_{\mathrm{red}}\).
The group \(H/R\) is a quotient of \(E^{(n)}\), and is therefore
smooth.  Consequently the canonical morphism
\[
  H_{\mathrm{red}}\longrightarrow H/R
\]
is an epimorphism.  Thus \(H\) is generated by \(R\) and
\(H_{\mathrm{red}}\), and hence equals \(H''\).  We have obtained
an exact sequence
\[
  \tag{\(\dagger\)}\label{sga3:xvii:subsection:5-5:equation:dagger}
  1\longrightarrow U''=U\cap E''
    \longrightarrow E''\longrightarrow H\longrightarrow1.
\]
But \(E''\) has the same underlying space as \(H_1\), so \(U''\)
is radicial.  It is plainly enough to prove that the extension
\((\dagger)\) is trivial, and this follows from \SGARefLink{sga3:xvii:proposition:5:4:1}{5.4.1}.

\SGAProof{Proof of \SGARefLink{sga3:xvii:lemma:5:5:1:item:i}{(i)}\SGARefLink{sga3:xvii:section:b}{(b)}: the case where \(k\) is perfect and
\(H\) is connected}
The group \(U\) is an extension of an étale group by its identity
component \(U^0\).  The case where \(U\) is connected having just
been treated, it is enough to consider the case where \(U\) is
étale.  We shall use the following more precise lemma.

\SGAStatement{Lemma 5.5.1}{With the notation of \SGARefLink{sga3:xvii:theorem:5-1-1}{5.1.1}, suppose
in addition that \(U\) is étale.  Then:
\begin{enumerate}
\item[(i)] \SGARefTarget{sga3:xvii:lemma:5:5:1:item:i}
  If \(H\) is connected, it has a unique lift in \(E\), namely
  the identity component \(E^0\) of \(E\).
\item[(ii)] \SGARefTarget{sga3:xvii:lemma:5:5:1:item:ii}
  If \(k\) is algebraically closed, then \(E\) is trivial, and
  any two lifts of \(H\) in \(E\) are conjugate by an element of
  \(U(k)\).
\end{enumerate}}
\SGARefTarget{sga3:xvii:lemma:5:5:1}

\SGAProof{Proof of \SGARefLink{sga3:xvii:lemma:5:5:1:item:i}{(i)}}
Since formation of the identity component commutes with extension
of the base field, we may restrict ourselves to the case where
\(k\) is algebraically closed.  If \(H\) is a torus, then \(E\)
is trivial by \SGARefLink{sga3:xvii:proposition:5:3:1}{5.3.1} and \SGARefLink{sga3:xvii:lemma:5:2:3:item:i}{5.2.3(i)}\SGARefLink{sga3:xvii:section:b}{(b)}, and \(E^0\) is plainly the
unique lift of \(H\).  This already proves that in the general
case \(E^0\cap U\) is radicial.  It is also étale, since \(U\) is
étale, and is therefore zero.  The morphism
\[
  E^0\longrightarrow H
\]
is consequently a monomorphism; it is flat because \(E^0\) is
open in \(E\), and it is surjective because \(H\) is connected.
It is therefore an isomorphism.  If \(H'\) is any other lift of
\(H\), then \(H'\) is connected, hence contained in \(E^0\), and
thus equal to \(E^0\).

\SGAProof{Proof of \SGARefLink{sga3:xvii:lemma:5:5:1:item:ii}{(ii)}}
Let \(H^0\) be the identity component of \(H\).  By \SGARefLink{sga3:xvii:lemma:5:5:1:item:i}{(i)}, \(E^0\)
is the unique lift of \(H^0\) in \(E\).  Put
\[
  E'=E/E^0,\qquad H'=H/H^0.
\]
There is then an extension
\[
  1\longrightarrow U\longrightarrow E'
    \longrightarrow H'\longrightarrow1.
\]
Since \(H'\) is étale, this extension is trivial by
\SGARefLink{sga3:xvii:lemma:5:2:3:item:i}{5.2.3(i)}\SGARefLink{sga3:xvii:section:b}{(b)}.  If \(H'_1\) is a lift of \(H'\) in \(E'\), and
\(H_1\) is its inverse image in \(E\), then \(H_1\) plainly
lifts \(H\) in \(E\).  If \(H_2\) is a second lift of \(H\) in
\(E\), then it contains \(E^0\).  By \SGARefLink{sga3:xvii:lemma:5:2:3:item:i}{5.2.3(i)}\SGARefLink{sga3:xvii:section:b}{(b)}, the image of
\(H_2\) in \(E'\) is conjugate to \(H'_1\) by an element
\(u\in U(k)\), and hence
\[
  H_2=\operatorname{int}(u)H_1.
\]

\SGAStatement{Remark 5.5.2}{Under the hypotheses of \SGARefLink{sga3:xvii:lemma:5:5:1:item:i}{5.5.1(i)},
the group \(E^0\) centralizes \(U\).}
\SGARefTarget{sga3:xvii:remark:5:5:2}

\SGAProof{Proof of \SGARefLink{sga3:xvii:theorem:5:1:1:item:i}{5.1.1(i)}\SGARefLink{sga3:xvii:appendix:A}{(a)}}
Using the composition series
\[
  1\longrightarrow U^0\longrightarrow U
    \longrightarrow U/U^0\longrightarrow1,
\]
assertion \SGARefLink{sga3:xvii:lemma:5:5:1:item:i}{(i)}\SGARefLink{sga3:xvii:appendix:A}{(a)} follows by combining \SGARefLink{sga3:xvii:lemma:5:5:1:item:i}{(i)}\SGARefLink{sga3:xvii:section:b}{(b)} with \SGARefLink{sga3:xvii:lemma:5:5:1:item:ii}{5.5.1(ii)}.

Before proving \SGARefLink{sga3:xvii:theorem:5:1:1:item:i}{5.1.1(i)}\SGARefLink{sga3:xvii:appendix:C}{(c)} and (d), we shall first establish
\SGARefLink{sga3:xvii:theorem:5:1:1:item:ii}{5.1.1(ii)}.

\SGAProof{Proof of \SGARefLink{sga3:xvii:theorem:5:1:1:item:ii}{5.1.1(ii)}\SGARefLink{sga3:xvii:appendix:A}{(a)}}
For want of a satisfactory general statement, we shall describe
a number of cases in which, when \(H\) is smooth, two lifts of
\(H\) in \(E\) are conjugate.

\SGAStatement{\SGAInlineRefTarget{sga3:xvii:proposition:5-6-1}Proposition 5.6.1}{With the notation of \SGARefLink{sga3:xvii:theorem:5-1-1}{5.1.1},
suppose in addition that \(H\) is smooth, and let \(H'\) and
\(H''\) be two lifts of \(H\) in \(E\).  Then \(H'\) and \(H''\)
are conjugate by an element of \(U(k)\) in each of the following
cases:
\begin{enumerate}
\item[(a)] \SGARefTarget{sga3:xvii:proposition:5:6:1:item:a}
  \(k\) is algebraically closed.
\item[(b)] \SGARefTarget{sga3:xvii:proposition:5:6:1:item:b}
  \(U\) is commutative.
\item[(c)] \SGARefTarget{sga3:xvii:proposition:5:6:1:item:c}
  \(k\) is perfect and \(U\) is connected.
\item[(d)] \SGARefTarget{sga3:xvii:proposition:5:6:1:item:d}
  \(U\) is \(k\)-solvable.
\item[(e)] \SGARefTarget{sga3:xvii:proposition:5:6:1:item:e}
  The centralizer of \(H'\) in \(U\) is \(k\)-solvable.
\item[(f)] \SGARefTarget{sga3:xvii:proposition:5:6:1:item:f}
  The group of multiplicative type \(H\) is split by a finite
  Galois extension \(K/k\) of degree prime to \(p\).
\end{enumerate}}
\SGARefTarget{sga3:xvii:proposition:5:6:1}

% END inputs/XVII/tex/components/24_expose_XVII_proposition5_4_1_completion_through_proposition5_6_1_statement_pages026_028.tex

% BEGIN inputs/XVII/tex/components/25_expose_XVII_proposition5_6_1_proof_and_section5_7_pages028_031.tex
% Authority: Expo17.pdf, local pp.28--31 / combined pp.990--993 /
% printed pp.587--593, completing Proposition 5.6.1 and all of
% subsection 5.7.  The next cursor is subsection 5.8.

\SGAProof{Proof}
Parts \SGARefLink{sga3:xvii:appendix:A}{(a)}, \SGARefLink{sga3:xvii:section:b}{(b)}, \SGARefLink{sga3:xvii:appendix:C}{(c)}, and (d) follow from \SGARefLink{sga3:xvii:proposition:5:3:1}{5.3.1}.

\smallskip
\noindent
(e) Let \(Z\) be the centralizer of \(H'\) in \(U\), and let \(T\)
be the transporter of \(H'\) into \(H''\).  By \SGARefLink{sga3:xvii:appendix:A}{(a)}, \(T\) is
nonempty; hence \(T\) is a principal homogeneous space under \(Z\).
The hypothesis on \(Z\) therefore implies that \(T\) is trivial.

\smallskip
\noindent
(f) Proceeding as in \SGARefLink{sga3:xvii:proposition:5:3:1}{5.3.1}, we see that it is enough to consider the
case where \(H\) is étale.  Suppose first that \(H\) is
diagonalizable, defined by an abstract group \(M\) of order \(m\)
prime to \(p\).  The two lifts \(H'\) and \(H''\) determine a
one-cocycle of \(M\) with values in \(U(k)\), that is, a map
\[
  h:M\longrightarrow U(k)
\]
such that
\[
  h(mn)=h(m)\,{}^{m}\!h(n)
\]
for every pair \(m,n\) of elements of \(M\).  The groups \(H'\) and
\(H''\) are conjugate by an element of \(U(k)\) if and only if there
exists \(a\in U(k)\) such that
\[
  h(m)=a^{-1}({}^{m}\!a).
\]
Now the abstract group \(U(k)\) has a composition series whose
successive quotients are commutative and annihilated by a power of
\(p\) (in view of \SGARefLink{sga3:xvii:proposition:5:6:1:item:c}{5.6.1(c)}, we may suppose that \(p>0\)).  It follows
at once in this case that \(h\) is a coboundary.

Consider now the general case.  Again let \(Z\) denote the
centralizer of \(H'\) in \(U\), and \(T\) the transporter of \(H'\)
into \(H''\); thus \(T\) is a torsor under \(Z\).  By hypothesis there
is a finite Galois extension \(K/k\), with Galois group \(G\) of
order prime to \(p\), which splits \(H\).  By the preceding case,
\(H'_K\) and \(H''_K\) are conjugate by an element of \(U(K)\), and
therefore \(T_K\) is trivial.  Consequently \(T\) is defined by an
element of
\[
  H^1(G,Z).
\]
For the same reasons as above, the hypothesis on \(G\) implies that
\[
  H^1(G,Z)=e;
\]
hence \(T\) is trivial.

\SGASubsection{5.7}{Proof of \SGARefLink{sga3:xvii:theorem:5:1:1:item:ii}{5.1.1(ii)}\SGARefLink{sga3:xvii:section:b}{(b)}}
\SGARefTarget{sga3:xvii:subsection:5:7}

\SGAStatement{Lemma 5.7.1}{Let \(S\) be a prescheme, let \(G\) be an
\(S\)-group prescheme, separated and smooth over \(S\), and let \(H\)
be an \(S\)-group of multiplicative type acting on \(G\).  Then the
\(S\)-functor
\[
  Z=G^H
\]
of invariants of \(G\) under \(H\) is representable by a closed
smooth \(S\)-group subprescheme of \(G\).}
\SGARefTarget{sga3:xvii:lemma:5:7:1}

The representability of \(G^H\) by a closed group subprescheme of
\(G\) follows from \SGARefLink{sga3:viii:examples:6:5:item:d:01}{Exposé VIII, 6.5(d)}.  Consider the semidirect
product
\[
  K=G\rtimes_S H.
\]
The centralizer of \(H\) in \(K\) is then \(Z\times_S H\).  To prove
that \(Z\) is smooth, we must verify the following.  Suppose that
\(S\) is affine, that \(S_0\) is a closed subscheme of \(S\) defined
by a square-zero ideal, and that \(u_0\in Z(S_0)\).  We must lift
\(u_0\) to an element of \(Z(S)\).

Since \(G\) is smooth over \(S\), there is an element \(u\in G(S)\)
lifting \(u_0\).  Let \(i\) be the canonical immersion of \(H\) into
\(K\), and consider the two \(S\)-group morphisms
\[
  H\rightrightarrows K,\qquad i
  \quad\hbox{and}\quad
  j=\operatorname{int}(u)i.
\]
Since \(u_0\in Z(S_0)\), we have \(i_{S_0}=j_{S_0}\).  By
\SGARefLink{sga3:ix:theorem:3:2}{Exposé IX, 3.2}, there exists \(v\in K(S)\), lifting the identity
section in \(K(S_0)\), such that
\[
  i=\operatorname{int}(v)j
   =\operatorname{int}(vu)i.
\]
Thus
\[
  vu=(u',v')\in (Z\times_S H)(S).
\]
It follows that \(u'\in Z(S)\), and \(u'\) lifts \(u_0\).

\SGAStatement{Lemma 5.7.2}{Let \(k\) be a field and \(H\) a
\(k\)-algebraic group, and let \(P(H)\) denote the following
property:
\begin{quote}
For every smooth unipotent \(k\)-group \(U\), and every extension
\(E\) of \(H\) by \(U\), any two lifts of \(H\) in \(E\) are
conjugate by an element of \(U(k)\).
\end{quote}
Suppose that \(H\) is an algebraic group which is an extension of a
group \(H''\) of multiplicative type by a group \(H'\) of
multiplicative type.  If \(P(H')\) and \(P(H'')\) hold, then
\(P(H)\) holds.}
\SGARefTarget{sga3:xvii:lemma:5:7:2}

Indeed, let \(U\) be a smooth unipotent \(k\)-group, let \(E\) be an
extension of \(H\) by \(U\), and let \(H_1,H_2\) be two lifts of
\(H\) in \(E\).  Let \(H'_1,H'_2\) be the corresponding lifts of
\(H'\).  Since \(P(H')\) holds, there exists \(u\in U(k)\) such that
\[
  H'_2=\operatorname{int}(u)H'_1.
\]
After replacing \(H_1\) by \(\operatorname{int}(u)H_1\), we may
suppose that \(H'_1=H'_2\); denote this common group simply by
\(H'\).  Put
\[
  E'=\operatorname{Centr}_E(H')
     =U^{H'}H_1
     =U^{H'}H_2.
\]
By \SGARefLink{sga3:xvii:lemma:5:7:1}{5.7.1}, \(U^{H'}=U'\) is smooth.  We then have an extension
\[
  1\longrightarrow U'\longrightarrow E'
    \longrightarrow H\longrightarrow1.
\]
By construction \(H'\) is central in \(E'\), and hence invariant.
Passing to the quotient gives an exact sequence
\[
  1\longrightarrow U'\longrightarrow E''
    \longrightarrow H''\longrightarrow1.
\]
Since \(U'\) is smooth and \(P(H'')\) holds, the images of \(H_1\)
and \(H_2\) in \(E''\) are conjugate by an element
\(u\in U'(k)\).  It follows that
\[
  H_1=\operatorname{int}(u)H_2.
\]

To prove \SGARefLink{sga3:xvii:theorem:5:1:1:item:ii}{5.1.1(ii)}\SGARefLink{sga3:xvii:section:b}{(b)}, observe that, since \(k\) is algebraically
closed, \(H\) has a composition series whose successive quotients
are smooth or, when \(p>0\), isomorphic to \(\boldsymbol\mu_p\).
Repeated application of \SGARefLink{sga3:xvii:lemma:5:7:2}{5.7.2} reduces us to the case where \(H\) is
smooth or \(H=\boldsymbol\mu_p\).  In the first case it is enough to
apply \SGARefLink{sga3:xvii:theorem:5:1:1:item:ii}{5.1.1(ii)}(a).  It remains to consider \(H=\boldsymbol\mu_p\).
Since \(U\) is smooth, \(U\) has a characteristic composition series
whose successive quotients are étale or isomorphic to
\(\mathbb G_a^r\), by \SGARefLink{sga3:xvii:lemma:3-9}{3.9}.  If \(U\) is étale, apply \SGARefLink{sga3:xvii:lemma:5:5:1}{5.5.1}.  We are
finally reduced to
\[
  H=\boldsymbol\mu_p,\qquad U=\mathbb G_a^r.
\]

We must prove that
\[
  H^1(\boldsymbol\mu_p,U)=0.
\]
The method used in \SGARefLink{sga3:xvii:proposition:5:4:1}{5.4.1} no longer applies, because
\(\boldsymbol\mu_p\) does not in general act linearly on
\(\mathbb G_a^r\).  Fix the following notation: \(H'\) denotes a
lift of \(H\) in \(E\), and
\[
  \mathfrak e=\operatorname{Lie}E,\qquad
  \mathfrak u=\operatorname{Lie}U,\qquad
  \mathfrak h=\operatorname{Lie}H,\qquad
  \mathfrak h'=\operatorname{Lie}H'.
\]
Let \(X\) be a nonzero element of \(\mathfrak h\) such that
\[
  X^{(p)}=X
\]
(Appendix II, \SGARefLink{sga3:xvii:lemma:3-1}{3.1}), and let \(X'\) be its lift in \(\mathfrak h'\).

Since \(\boldsymbol\mu_p\) is a radicial group of height \(1\), the
set of lifts of \(H\) in \(E\) is in bijective correspondence with
the set \(A\) of elements \(Y\in\mathfrak e\) such that
\[
  Y^{(p)}=Y
\]
and whose image is \(X\) (Appendix II, \SGARefLink{sga3:ii:subsection:2-2}{2.2}).  Likewise, if
\(Y\in A\) corresponds to a lift \(H''\) of \(H\) in \(E\), and if
\(u\in U(k)\), then
\[
  \operatorname{int}(u)H'=H''
  \quad\Longleftrightarrow\quad
  \operatorname{Ad}(u)X'=Y.
\]
Let \(B\) be the subset of \(\mathfrak e\) consisting of the
elements \(\operatorname{Ad}(u)X'\), with \(u\in U(k)\).  Plainly
\(B\subset A\), and it remains only to prove \(A=B\) when \(k\) is
algebraically closed.

\smallskip
\noindent
(a) \emph{Study of \(A\).}
Since \(\mathfrak u\) is commutative and normalized by
\(\mathfrak h'\), Jacobson's formula (Exposé VII\textsubscript{A},
\SGARefLink{sga3:VIIB:5.2:theorem}{5.2}) gives, for \(u\in\mathfrak u\),
\[
  (X'+u)^{(p)}
   =X'{}^{(p)}+u^{(p)}
      +(\operatorname{ad}X')^{p-1}(u)
   =X'+(\operatorname{ad}X')^{p-1}(u).
\]
Consequently,
\[
  X'+u\in A
  \quad\Longleftrightarrow\quad
  u=(\operatorname{ad}X')^{p-1}(u).
\]
Let
\[
  \mathfrak u
   =\bigoplus_{n\in\mathbb Z/p\mathbb Z}\mathfrak u_n
   =\mathfrak u_0\oplus
      \bigoplus_{n\in(\mathbb Z/p\mathbb Z)^\times}\mathfrak u_n
\]
be the canonical decomposition of \(\mathfrak u\) under the action
of \(H'\).  If \(u\in\mathfrak u_0\), then
\(\operatorname{ad}X'(u)=0\).  If
\(u\in\mathfrak u_n\), with
\(n\in(\mathbb Z/p\mathbb Z)^\times\), then
\[
  (\operatorname{ad}X')^{p-1}(u)=u.
\]
Thus \(Y=X'+u\) belongs to \(A\) if and only if
\[
  u\in
  \bigoplus_{n\in(\mathbb Z/p\mathbb Z)^\times}\mathfrak u_n.
\]
In particular, \(A\) is the set of \(k\)-rational points of an
irreducible subscheme of \(\mathbf W(\mathfrak e)\), of dimension
\[
  \operatorname{rk}\mathfrak u-\operatorname{rk}\mathfrak u_0
  =
  \operatorname{rk}\mathfrak u
  -\operatorname{rk}\operatorname{Centr}_{\mathfrak u}(X').
\]

\smallskip
\noindent
(b) \emph{Study of \(B\).}
We shall need the following lemma.

\SGAStatement{Lemma 5.7.3 (Rosenlicht)}{Let \(U\) be a unipotent
algebraic group over a field \(k\), acting on a quasi-affine
\(k\)-scheme \(X\).  Then the orbit of every point \(x\in X(k)\) is
closed in \(X\).  Here, by the orbit of \(x\), we mean the subset of
\(\operatorname{ens}(X)\) which is the image of \(U\) under the
morphism \(g\mapsto g\cdot x\).}
\SGARefTarget{sga3:xvii:lemma:5:7:3}

By functoriality, \(U\) acts on the affine envelope of \(X\), that
is,
\[
  \operatorname{Spec}\Gamma(X,\mathcal O_X),
\]
so we may suppose that \(X\) is affine.  We may then suppose that
\(k\) is algebraically closed, \(X\) is reduced, and \(U\) is smooth
(observe that \(U_{\mathrm{red}}\) acts on \(X_{\mathrm{red}}\) when
\(k\) is perfect).  Let \(Y\) be the schematic image of \(U\)
(EGA I, 9.5.1) under the morphism \(g\mapsto g\cdot x\).  It is a
closed reduced subscheme of \(X\), on which \(U\) acts.  It follows
readily from EGA IV, 1.8.6 that the orbit of \(x\) is an open subset
\(Z\) of \(Y\), dense in \(Y\).  We must prove that
\[
  Z=\operatorname{ens}(Y).
\]
Let \(F\) be the closed reduced subscheme of \(Y\) whose underlying
space is \(Y\setminus Z\).  Thus \(F=V(J)\), where \(J\) is a
nonzero ideal of \(\Gamma(Y,\mathcal O_Y)\).  Since \(U\) is smooth,
\(U\) acts on \(F\), hence on \(J\), and therefore
\[
  J^U\ne0
\]
by \SGARefLink{sga3:xvii:proposition:3:2}{3.2}.  If \(a\) is a nonzero element of \(J^U\), then \(a\) is
constant on the orbit \(Z\), and hence on \(Y\), because \(Z\) is
dense in \(Y\).  The ideal \(J\) consequently contains \(k\), and
\(F=\varnothing\).

Apply the preceding lemma to the group \(U\) acting on the affine
space \(\mathbf W(\mathfrak e)\) through the adjoint representation.
The orbit of \(X'\) is the underlying set of a closed subscheme of
\(\mathbf W(\mathfrak e)\).  Moreover, the stabilizer \(Z\) of
\(X'\) is its centralizer in \(U\), and there is a closed immersion
\[
  U/Z\longrightarrow\mathbf W(\mathfrak e).
\]
Thus the orbit of \(X'\) is the underlying space of a closed
subscheme of \(\mathbf W(\mathfrak e)\) of dimension
\[
  \dim U-\dim Z.
\]

\smallskip
\noindent
\SGARefLink{sga3:xvii:appendix:C}{(c)} \emph{End of the proof of \SGARefLink{sga3:xvii:theorem:5:1:1:item:ii}{5.1.1(ii)}\SGARefLink{sga3:xvii:section:b}{(b)}.}
When \(k\) is algebraically closed, the canonical map
\[
  U(k)\longrightarrow(U/Z)(k)
\]
is surjective.  By \SGARefLink{sga3:xvii:section:b}{(b)}, \(B\) is therefore the set of \(k\)-rational
points of a closed subscheme of \(\mathbf W(\mathfrak e)\) of
dimension \(\dim U-\dim Z\).  In view of \SGARefLink{sga3:xvii:appendix:A}{(a)}, to prove \(A=B\) it is
enough to show that
\[
  \operatorname{rk}\mathfrak u
  -\operatorname{rk}\operatorname{Centr}_{\mathfrak u}(X')
  \leq
  \dim U-\dim\operatorname{Centr}_U(X').
\]
Since \(U\) is smooth,
\[
  \dim U=\operatorname{rk}\mathfrak u.
\]
On the other hand, \SGARefLink{sga3:ii:paragraph:5-3-3}{Exposé II, 5.3.3} gives
\[
  \dim\operatorname{Centr}_U(X')
  \leq
  \operatorname{rk}\operatorname{Centr}_{\mathfrak u}(X'),
\]
which proves the assertion.  Notice in fact that
\[
  \dim\operatorname{Centr}_U(X')
  =
  \operatorname{rk}\operatorname{Centr}_{\mathfrak u}(X');
\]
this also proves once more that \(Z\) is smooth, as in \SGARefLink{sga3:xvii:lemma:5:7:1}{5.7.1}.

% END inputs/XVII/tex/components/25_expose_XVII_proposition5_6_1_proof_and_section5_7_pages028_031.tex

% BEGIN inputs/XVII/tex/components/26_expose_XVII_section5_8_pages031_032.tex
% Authority: Expo17.pdf, local pp.31--32 / combined pp.993--994 /
% printed pp.593--597, all of subsection 5.8.  The next cursor is
% subsection 5.9.

\SGASubsection{5.8}{End of the proof of \SGARefLink{sga3:xvii:theorem:5:1:1:item:i}{5.1.1(i)}}
\SGARefTarget{sga3:xvii:subsection:5:8}

It remains to prove (i)\SGARefLink{sga3:xvii:appendix:C}{(c)} and (i)(d).

\SGAProof{Proof of \SGARefLink{sga3:xvii:theorem:5:1:1:item:i}{5.1.1(i)}(d): \(U\) smooth and \(H\) connected}
We shall need the following lemma.

\SGAStatement{Lemma 5.8.1}{With the notation of \SGARefLink{sga3:xvii:theorem:5-1-1}{5.1.1}, suppose that
\(U\) is smooth and \(H\) is radicial, and let \(H_1\) be an
algebraic subgroup of \(H\) which has a lift \(H'_1\) in \(E\).
Then \(H\) has a lift \(H'\) in \(E\) which contains \(H'_1\).}
\SGARefTarget{sga3:xvii:lemma:5:8:1}

\SGAProof{Proof}
By induction on the height of \(H/H_1\), we may suppose that
\(H/H_1\) has height \(1\).  Put
\[
  C'=\operatorname{Centr}_E(H'_1).
\]
We claim that the canonical morphism \(C'\longrightarrow H\) is an
epimorphism.  To prove this, we may suppose that \(k\) is
algebraically closed.  The extension \(E\) is then trivial, by
\SGARefLink{sga3:xvii:theorem:5:1:1:item:i}{5.1.1(i)}\SGARefLink{sga3:xvii:appendix:A}{(a)}.  Let \(H''\) be a lift of \(H\) in \(E\), and let
\(H''_1\) be the inverse image of \(H_1\) in \(H''\).  The groups
\(H'_1\) and \(H''_1\) are two lifts of \(H_1\) in \(E\); hence they
are conjugate by an element of \(U(k)\), since \(k\) is
algebraically closed and \(U\) is smooth, by \SGARefLink{sga3:xvii:theorem:5:1:1:item:ii}{5.1.1(ii)}\SGARefLink{sga3:xvii:section:b}{(b)}.
Consequently, to prove the assertion about \(C'\), it is enough to
prove it for
\[
  C''=\operatorname{Centr}_E(H''_1).
\]
But \(C''\) contains \(H''\), so the assertion is clear.

It also follows from \SGARefLink{sga3:xvii:lemma:5:7:1}{5.7.1} that \(C'\cap U\) is smooth.  We may
therefore replace \(E\) by \(C'\), and thus suppose that \(H'_1\) is
central.  Passing then to the quotient by \(H'_1\), we may suppose
that \(H_1=0\) and that \(H\) has height \(1\).

Since \(U\) is smooth, we have the following exact sequence of
\(p\)-Lie algebras (Appendix II, \SGARefLink{sga3:xvii:proposition:3:2}{3.2}):
\[
  (*)\qquad
  0\longrightarrow\operatorname{Lie}U
   \longrightarrow\operatorname{Lie}E
   \longrightarrow\operatorname{Lie}H
   \longrightarrow0.
\]
In view of Appendix II, \SGARefLink{sga3:ii:subsection:2-2}{2.2}, to say that \(E\) is trivial is
equivalent to saying that \((*)\) is a trivial extension of
\(p\)-Lie algebras.  Suppose that \(H\ne0\), and hence that
\(\operatorname{Lie}H\ne0\).  Suppose also that we have found a
nonzero \(p\)-Lie subalgebra
\[
  \mathfrak h_1\subset\mathfrak h=\operatorname{Lie}H
\]
which lifts to a \(p\)-Lie subalgebra
\(\mathfrak h'_1\subset\operatorname{Lie}E\).  By the cited result,
there is a subgroup \(H_1\) of \(H\), with
\(\operatorname{Lie}H_1=\mathfrak h_1\), and a lift \(H'_1\) of
\(H_1\) in \(E\), with
\(\operatorname{Lie}H'_1=\mathfrak h'_1\).  Applying once more the
reduction just described, we are brought back to the same problem
with \(H\) replaced by \(H/H_1\).  Since \(H\) has height \(1\),
\[
  \operatorname{Lie}(H/H_1)
   =\operatorname{Lie}H/\operatorname{Lie}H_1
\]
(loc.\ cit.), and consequently
\[
  \operatorname{rk}\operatorname{Lie}(H/H_1)
  \leq \operatorname{rk}\operatorname{Lie}H-1.
\]
Thus, by induction on the rank of \(\operatorname{Lie}H\), it is
enough, when \(\mathfrak h\ne0\), to find a nonzero Lie subalgebra
of \(\operatorname{Lie}H\) which lifts to \(\operatorname{Lie}E\).
This follows from the next lemma.

\SGAStatement{\SGAInlineRefTarget{sga3:xvii:lemma:5-8-2}Lemma 5.8.2}{Let \(k\) be a field of characteristic
\(p>0\), and let
\[
  \varphi:\mathfrak g\longrightarrow\mathfrak h
\]
be a surjective morphism of finite-rank \(p\)-Lie algebras over
\(k\).  Then:
\begin{enumerate}
\item[\SGAInlineRefTarget{sga3:xvii:lemma:5-8-2:item:i}(i)] If \(\mathfrak h_{\overline{k}}\) is reductive
  (\SGARefLink{sga3:xvii:lemma:4-2-2}{4.2.2}) and nonzero, there exists a reductive Lie subalgebra
  \(\mathfrak h_1\) of \(\mathfrak g\) whose image in
  \(\mathfrak h\) is nonzero.
\item[\SGAInlineRefTarget{sga3:xvii:lemma:5-8-2:item:ii}(ii)] If \(k\) is perfect and \(u\) is a unipotent element of
  \(\mathfrak h\), that is, if \(u^{(p^n)}=0\) for some \(n\), then
  \(u\) lifts to a unipotent element of \(\mathfrak g\).
\end{enumerate}}
\SGARefTarget{sga3:xvii:lemma:5:8:2}

\SGAProof{Proof}
Let \(X\) be a nonzero element of \(\mathfrak h\) in case \SGARefLink{sga3:xvii:lemma:5-8-2:item:i}{(i)}, and
let \(X=u\) in case \SGARefLink{sga3:xvii:lemma:5-8-2:item:ii}{(ii)}; choose a lift \(X'\) of \(X\) in
\(\mathfrak g\).  The \(p\)-Lie subalgebra \(\mathfrak h'\) of
\(\mathfrak g\) generated by \(X'\) is abelian (Exposé VII).

\smallskip
\noindent
\emph{Case \SGARefLink{sga3:xvii:lemma:5-8-2:item:i}{(i)}.}
The description given in \SGARefLink{sga3:xvii:lemma:4-2-2}{4.2.2} shows that the reductive part
(loc.\ cit.) of \(\mathfrak h'_{\overline{k}}\) is already defined
over \(k\); denote it by \(\mathfrak r'\).  We claim that the image
of \(\mathfrak r'\) in \(\mathfrak h\) is nonzero.  We may suppose
that \(k\) is algebraically closed, so that
\[
  \mathfrak h'=\mathfrak r'\oplus\mathfrak u',
\]
where \(\mathfrak u'\) is the unipotent part of
\(\mathfrak h'\).  If the image of \(\mathfrak r'\) in
\(\mathfrak h\) were zero, then the image of \(\mathfrak h'\) in
\(\mathfrak h\) would be unipotent, and hence zero because
\(\mathfrak h\) is reductive (\SGARefLink{sga3:VIIA:2.4:item:ii}{2.4(ii)}); but by construction that
image contains \(X\), a contradiction.

\smallskip
\noindent
\emph{Case \SGARefLink{sga3:xvii:lemma:5-8-2:item:ii}{(ii)}.}
The same argument applies, with the roles of \(\mathfrak r'\) and
\(\mathfrak u'\) interchanged.

\smallskip
\noindent
\emph{End of the proof of \SGARefLink{sga3:xvii:lemma:5:8:1}{5.8.1}.}
Suppose that \(H\ne0\).  By \SGARefLink{sga3:xvii:lemma:5-8-2}{5.8.2} there exists a nonzero reductive
\(p\)-Lie subalgebra
\(\mathfrak h'_1\subset\operatorname{Lie}E\).  Since
\(\operatorname{Lie}U\) is unipotent, by \SGARefLink{sga3:xvii:proposition:4:3:1:item:i}{4.3(i)}, we necessarily
have
\[
  (\operatorname{Lie}U)\cap\mathfrak h'_1=0.
\]
Thus \(\mathfrak h'_1\) is a lift of a \(p\)-Lie subalgebra of
\(\operatorname{Lie}H\).

\smallskip
\noindent
\emph{End of the proof of \SGARefLink{sga3:xvii:theorem:5:1:1:item:i}{5.1.1(i)}(d).}
By \SGARefLink{sga3:xvii:lemma:5:8:1}{5.8.1} there is a family of algebraic subgroups
\(H'_n\) (\(n\in\mathbb N\)) of \(E\) such that
\(H'_{n+1}\) contains \(H'_n\), and such that \(H'_n\) lifts
\({}_{F^n}H\).  The descending sequence of subgroups
\[
  \operatorname{Centr}_E(H'_n)
\]
is stationary; let \(C\) be its stationary value.  The center \(Z\)
of \(C\) contains \(H'_n\) for every \(n\).  Hence the image of
\(Z\) in \(H\) contains \({}_{F^n}(H)\) for every \(n\), and is
therefore an open subgroup of \(H\) (Exposé VII\textsubscript{A},
\S 4).  It is consequently equal to \(H\), since \(H\) is
connected.  To prove that \(E\) is trivial, we may thus replace
\(E\) by \(Z\), and hence suppose that \(E\) is commutative.  We
shall see in \SGARefLink{sga3:xvii:theorem:7:2:1}{7.2.1} that \(E\) then contains a maximal subgroup \(M\)
of multiplicative type, whose formation commutes with extension of
the base field.  Since \(E_{\overline{k}}\) is a trivial extension,
by \SGARefLink{sga3:xvii:theorem:5:1:1:item:i}{5.1.1(i)}\SGARefLink{sga3:xvii:appendix:A}{(a)}, and \(U\) is unipotent, it is clear that \(M\) is
the unique lift of \(H\) in \(E\).

\SGAProof{Proof of \SGARefLink{sga3:xvii:theorem:5:1:1:item:i}{5.1.1(i)}\SGARefLink{sga3:xvii:appendix:C}{(c)}: \(U\) \(k\)-solvable}
Since \((H/H^0)(\overline{k})\) has order prime to \(p\), duality
immediately shows that there exists an integer \(n\) such that
\({}_nH\) is an étale subgroup and the canonical morphism
\[
  {}_nH\longrightarrow H/H^0
\]
is an epimorphism.  Thus \(H/{}_nH\) is connected.  By \SGARefLink{sga3:xvii:lemma:5:2:3:item:i:d}{5.2.3(d)},
there exists a lift \(H'\) of \({}_nH\) in \(E\).  As at the
beginning of the proof of \SGARefLink{sga3:xvii:lemma:5:8:1}{5.8.1}, one sees that
\[
  C=\operatorname{Centr}_E(H')
\]
is a subgroup of \(E\) such that \(C\longrightarrow H\) is an
epimorphism and \(C\cap U\) is smooth.  Replacing \(E\) by \(C\)
and passing to the quotient by \(H'\), we are reduced to the case
where \(U\) is smooth and \(H\) is connected, namely case \SGARefLink{sga3:xvii:lemma:5-8-2:item:i}{(i)}(d).

This completes the proof of \SGARefLink{sga3:xvii:theorem:5-1-1}{5.1.1}.

% END inputs/XVII/tex/components/26_expose_XVII_section5_8_pages031_032.tex

% BEGIN inputs/XVII/tex/components/27_expose_XVII_section5_9_counterexamples_pages032_034.tex
% Authority: Expo17.pdf, local pp.32--34 / combined pp.994--996 /
% printed pp.597--601, all of subsection 5.9.  The next cursor is
% section 6, local p.35 / combined p.997 / printed p.602.

\SGASubsection{5.9}{Counterexamples}
\SGARefTarget{sga3:xvii:subsection:5:9}

Let us first indicate a procedure for obtaining nontrivial
extensions of a \(k\)-group \(H\) of multiplicative type by a
unipotent \(k\)-group \(U\).  Suppose that an action of \(H\) on
\(U\) is given, and let
\[
  E=U\rtimes H
\]
be the semidirect product.  Let, moreover, a class in
\(\mathrm H^1(k,U)\) be represented by a \(1\)-cocycle \(a\).
The group \(U\) acts on \(E\) by inner automorphisms, and \(a\)
therefore defines a \(k\)-form of \(E\), denoted by \(E_a\).
Suppose in addition that \(U\) is commutative.  Then \(U\) acts
trivially on \(U\) and on the quotient \(E/U=H\), so \(E_a\) is
again an extension of \(H\) by \(U\).

For simplicity, suppose that \(H\) is an étale diagonalizable
group, so that the Galois group
\(\mathscr G=\operatorname{Gal}(\overline{k}/k)\) acts trivially
on \(H(\overline{k})\).  Its action on \(E_a(\overline{k})\) is
then given by
\[
  {}^g(u,h)
  =\bigl({}^g u+a(g)-{}^h a(g),h\bigr)
  \qquad
  \bigl(g\in\mathscr G,\ u\in U(\overline{k}),\
        h\in H(\overline{k})\bigr).
\]
If \(h\in H(k)\), and \(X\) is its inverse image in \(E_a\), then
\(X\) is a torsor under \(U\) defined by the class of the
\(1\)-cocycle
\[
  g\longmapsto a(g)-{}^h a(g)
\]
of \(\mathscr G\) with values in \(U\).  It follows that if there
is a point \(h\in H(k)\) for which this cocycle is nontrivial,
then the extension \(E_a\) is nontrivial.  We apply this
construction in the following two particular cases.

\smallskip
\noindent
\emph{\SGARefLink{sga3:xvii:appendix:A}{(a)} A nontrivial extension of an étale diagonalizable group
\(H\) by \(\mathbb Z/p\mathbb Z\).}
Take
\[
  H=(\mathbb Z/p\mathbb Z)^\times,
  \qquad U=\mathbb Z/p\mathbb Z,
\]
with \(H\) acting on \(U\) by multiplication.  Let \(k\) be a
field of characteristic \(p\), let \(K/k\) be a Galois extension
whose Galois group \(\mathscr G\) is isomorphic to
\(\mathbb Z/p\mathbb Z\), and take a nonzero element
\[
  a\in\mathrm H^1(\mathscr G,U)
   =\operatorname{Hom}(\mathscr G,U).
\]
Then \(E_a\) is the required example.

\smallskip
\noindent
\emph{\SGARefLink{sga3:xvii:section:b}{(b)} A nontrivial extension of an étale diagonalizable group
\(H\) by a smooth connected unipotent group \(U\).}
Choose a nonperfect field \(k\) admitting a \(k\)-form \(U\) of
\(\mathbb G_a\) such that \(\mathrm H^1(k,U)\ne0\).  For example
(see J.-P.\ Serre, \emph{Cohomologie Galoisienne}), one may take
for \(k\) the fraction field of a discrete valuation ring of
equal characteristic \(p>0\), and for \(U\) the algebraic
subgroup of \(\mathbb G_a\times_k\mathbb G_a\) defined by
\[
  X^p+X+tY^p=0,
\]
where \(t\) is a uniformizer of \(k\).  Indeed, suppose for
simplicity that \(k\) contains the \((p-1)\)-st roots of unity.
There is then an exact sequence of algebraic groups
\[
  0\longrightarrow\mathbb Z/p\mathbb Z
   \longrightarrow U
   \longrightarrow\mathbb G_a
   \longrightarrow0,
  \qquad
  (x,y)\longmapsto y,
\]
and hence an exact cohomology sequence
\[
  \mathbb G_a(k)\xrightarrow{\,d\,}
  \mathrm H^1(k,\mathbb Z/p\mathbb Z)
  \longrightarrow\mathrm H^1(k,U)
  \longrightarrow0,
\]
where \(d\) associates to \(x\) the homogeneous space under
\(\mathbb Z/p\mathbb Z\) defined by
\[
  X^p+X+tx^p=0.
\]
Moreover,
\[
  \mathrm H^1(k,\mathbb Z/p\mathbb Z)\simeq k/P(k),
  \qquad P(x)=x^p+x,
\]
and therefore
\[
  \mathrm H^1(k,U)\simeq k/\bigl(P(k)+tk^p\bigr).
\]
If, in addition, \(p\ne2\), then \(t^2\) does not belong to
\(P(k)+tk^p\); consequently \(\mathrm H^1(k,U)\ne0\).

The group \(\boldsymbol\mu_{p-1}\) acts on \(U\) by
\[
  (h,x,y)\longmapsto(hx,hy).
\]
Let \(\mathscr G\) be the Galois group of the extension \(K/k\)
defined by
\[
  X^p+X+t^2=0,
\]
and let \(a\in\mathrm H^1(\mathscr G,U)\) be the nonzero class
described above.  It follows immediately that \(E_a\) is a
nontrivial extension of \(\boldsymbol\mu_{p-1}\) by \(U\).

\smallskip
\noindent
\emph{\SGAInlineRefTarget{sga3:xvii:subsection:5-9:item:c}\SGARefLink{sga3:xvii:appendix:C}{(c)} A nontrivial extension of \(\mathbb G_m\) by
\(\boldsymbol\alpha_p\).}
By \SGARefLink{sga3:xvii:theorem:5:1:1:item:i}{5.1.1(i)}\SGARefLink{sga3:xvii:section:b}{(b)}, such an extension can exist only over a
nonperfect field.  Let \(k\) therefore be nonperfect, and let
\(G\) be the semidirect product of \(U=\mathbb G_a\) by
\(H=\mathbb G_m\), where \(\mathbb G_m\) acts on
\(\mathbb G_a\) by homotheties.  Since \(U\) is invariant in
\(G\), so is \({}_{F}U\).  Put
\[
  G'=G/{}_{F}U.
\]
The group \(G'\) is isomorphic to
\(\mathbb G_a\rtimes\mathbb G_m\), where \(\mathbb G_m\) acts on
\(\mathbb G_a\) by
\[
  (h,u)\longmapsto h^p u.
\]
The functor \(\mathscr T_G\) (respectively
\(\mathscr T_{G'}\)) of subtori of \(G\) (respectively \(G'\))
(see Exposé XV) is isomorphic to \(\mathbb G_a\), and the morphism
\[
  \mathscr T_G\longrightarrow\mathscr T_{G'}
\]
induced by \(G\longrightarrow G'\) identifies with
\[
  u\longmapsto u^p.
\]
It follows that if \(T'\) is a subtorus of \(G'\) corresponding
to a point \(x\in k\simeq\mathbb G_a(k)\) such that
\(x^{1/p}\notin k\), then its inverse image \(E\) in \(G\) is an
extension of the torus \(T'\) by
\({}_{F}U=\boldsymbol\alpha_p\).  It has no maximal torus defined
over \(k\), and hence is nontrivial.  As a subgroup of
\[
  G=\operatorname{Spec}k[U,T,T^{-1}],
\]
the group \(E\) is defined by
\[
  U^p=x-xT^p.
\]

\SGAStatement{Remark 5.9.1}{This last example shows that a
nonsmooth algebraic group defined over a nonperfect field need
not possess a maximal torus, and thus answers the question posed
in \SGARefLink{sga3:xiv:remark:1-5:item:b}{Exposé XIV, 1.5(b)}.}
\SGARefTarget{sga3:xvii:remark:5:9:1}

\smallskip
\noindent
\emph{(d) A trivial extension with nonconjugate lifts.}
We now give an example of a trivial extension \(E\) of a group
\(H\) of multiplicative type by a unipotent group \(U\), and of
two lifts \(H'\) and \(H''\) of \(H\) which are not conjugate by
an element of \(U(k)\).

Take for \(E\) the semidirect product of \(U=\mathbb G_a\) by
\[
  \boldsymbol\mu_p
  =\operatorname{Spec}k[T]/(T^p-1),
\]
where the action of \(\boldsymbol\mu_p\) on
\(\mathbb G_a=\operatorname{Spec}k[U]\) is defined by the
comorphism
\[
  U\longmapsto(U+U^p)T-U^p.
\]
The centralizer of \(\boldsymbol\mu_p\) in \(U\) is then the
étale group \(Z\) defined by
\[
  U+U^p=0.
\]
It follows that if \(k\) is not algebraically closed, the
canonical map
\[
  U(k)\longrightarrow(U/Z)(k)
\]
is not in general surjective.  In view of \SGARefLink{sga3:xvii:theorem:5:1:1:item:ii}{5.1.1(ii)}\SGARefLink{sga3:xvii:section:b}{(b)}, two
lifts of \(\boldsymbol\mu_p\) in \(E\) therefore need not be
conjugate by an element of \(U(k)\).

Here is another example, with \(k\) algebraically closed of
characteristic \(p>0\).  Let \(G\) be the radicial semidirect
product in which \(\boldsymbol\mu_p\) acts on
\(\boldsymbol\alpha_p\) by homotheties.  There is an exact
sequence of algebraic groups with \(\boldsymbol\mu_p\)-operators
\[
  0\longrightarrow\boldsymbol\alpha_p
   \longrightarrow\mathbb G_a
   \xrightarrow{\,x\longmapsto x^p\,}\mathbb G_a
   \longrightarrow0,
\]
where \(\boldsymbol\mu_p\) acts by homotheties on the first
\(\mathbb G_a\) and trivially on the second.  The exact
cohomology sequence (Appendix I, \SGARefLink{sga3:xvii:proposition:a-1-1}{Proposition 11}) gives
\[
  0\longrightarrow\mathbb G_a(k)
   \longrightarrow\mathrm H^1(\boldsymbol\mu_p,\boldsymbol\alpha_p)
   \longrightarrow\mathrm H^1(\boldsymbol\mu_p,\mathbb G_a).
\]
Since the last term is zero by I, \SGARefLink{sga3:i:theorem:5-3-3}{5.3.3}, the group
\(\mathrm H^1(\boldsymbol\mu_p,\boldsymbol\alpha_p)\) is nonzero.
Thus two lifts of \(\boldsymbol\mu_p\) in \(G\) need not be
conjugate.

% END inputs/XVII/tex/components/27_expose_XVII_section5_9_counterexamples_pages032_034.tex

% BEGIN inputs/XVII/tex/components/28_expose_XVII_section6_1_and_proofs_A_B_pages035_037.tex
% Authority: Expo17.pdf, local pp.35--37 / combined pp.997--999 /
% printed pp.602--608.  This component covers section 6 through the
% proof of Theorem 6.1.1(B), and stops before the proof of 6.1.1(C).

\SGAInlineRefTarget{sga3:xvii:section:6:site:da2ef89f51}\SGASection{6}{Extensions of a Unipotent Group by a Group of
Multiplicative Type}
\SGARefTarget{sga3:xvii:section:6}

\SGASubsection{6.1}{Statement of the Theorem}
\SGARefTarget{sga3:xvii:subsection:6:1}

\SGAStatement{\SGAInlineRefTarget{sga3:xvii:theorem:6-1-1}Theorem 6.1.1}{Let \(k\) be a field, \(U\) a
unipotent algebraic \(k\)-group, \(H\) a \(k\)-group of
multiplicative type, and \(E\) an algebraic \(k\)-group which is
an extension of \(U\) by \(H\), so that there is an exact sequence
\[
  1\longrightarrow H\longrightarrow E\longrightarrow U
    \longrightarrow1.
\]
Then the extension \(E\) is trivial and there exists a unique
lift of \(U\) in \(E\) in each of the following cases:
\begin{enumerate}
\item[(A)] \SGARefTarget{sga3:xvii:theorem:6:1:1:item:A}
  The group \(U\) is smooth and one of the following conditions
  holds:
  \begin{enumerate}
  \item[(i)] \SGARefTarget{sga3:xvii:theorem:6:1:1:item:A:i}
    \(U\) is connected and the canonical morphism
    \(E\longrightarrow U\) admits a section.
  \item[(ii)] \SGARefTarget{sga3:xvii:theorem:6:1:1:item:A:ii}
    \(U\) has a composition series whose successive quotients
    are isomorphic to \(\mathbb G_a\).
  \item[(iii)] \SGARefTarget{sga3:xvii:theorem:6:1:1:item:A:iii}
    \(H\) is étale.
  \item[(iv)] \SGARefTarget{sga3:xvii:theorem:6:1:1:item:A:iv}
    \(k\) is perfect.
  \end{enumerate}
\item[(B)] \SGARefTarget{sga3:xvii:theorem:6:1:1:item:B}
  \(U=\boldsymbol\alpha_p\) and \(k\) is perfect.
\item[(C)] \SGARefTarget{sga3:xvii:theorem:6:1:1:item:C}
  \(E\) is commutative and \(k\) is perfect.
\end{enumerate}}
\SGARefTarget{sga3:xvii:theorem:6:1:1}

\SGASubsection{6.2}{Proof of \SGARefLink{sga3:xvii:theorem:6:1:1:item:A}{6.1.1(A)}}
\SGARefTarget{sga3:xvii:subsection:6:2}

We first establish three lemmas.

\SGAStatement{Lemma 6.2.1}{Let \(S\) be a prescheme, \(E\) an
\(S\)-group prescheme which is an extension of an \(S\)-group
prescheme \(U\), with connected fibres, by an \(S\)-group \(H\)
of multiplicative type and finite type (that is, \(U\) is the
quotient of \(E\) by \(H\) for the fpqc topology).  Then \(E\) is
a central extension.}
\SGARefTarget{sga3:xvii:lemma:6:2:1}

\SGAProof{Proof}
Since \(H\) is commutative, \(U\) acts on \(H\) by inner
automorphisms, through an \(S\)-morphism of groups
\[
  u:U\longrightarrow\operatorname{Aut}_{S\text{-}\mathrm{gr}}(H).
\]
The functor
\(\operatorname{Aut}_{S\text{-}\mathrm{gr}}(H)\) is representable
by an étale \(S\)-scheme (\SGARefLink{sga3:x:corollary:5-10}{Exposé X, 5.10}), and consequently its
identity section is both an open and a closed immersion.  Since
\(U\) has connected fibres, it follows that \(u\) is the unit
morphism.

\SGAStatement{Lemma 6.2.2}{With the notation of \SGARefLink{sga3:xvii:theorem:6-1-1}{6.1.1}, if \(E\)
is trivial, there exists a unique lift of \(U\) in \(E\).}
\SGARefTarget{sga3:xvii:lemma:6:2:2}

\SGAProof{Proof}
Let \(U'\) and \(U''\) be two lifts of \(U\) in \(E\).  To prove
that \(U'=U''\), we may suppose that \(k\) is algebraically closed,
and it is enough to prove that \(\mathrm H^1(U,H)=0\).  If \(U\)
is connected, then \(U\) centralizes \(H\) by \SGARefLink{sga3:xvii:lemma:6:2:1}{6.2.1}, and hence
\[
  \mathrm H^1(U,H)
    =\operatorname{Hom}_{k\text{-}\mathrm{gr}}(U,H)=0
\]
by 2.4(ii).

In the general case, let \(U_1'\) denote the unique lift in \(E\)
of the identity component \(U^0\) of \(U\), and put
\[
  N=\operatorname{Norm}_{E}(U_1').
\]
If \(g\in E(k)\), then
\(\operatorname{int}(g)U_1'\) is a lift of \(U^0\), and is
therefore equal to \(U_1'\); thus \(N(k)\supset E(k)\).
Moreover, \(N\) contains \(H\), by \SGARefLink{sga3:xvii:lemma:6:2:1}{6.2.1}, as well as \(U_1'\),
and it follows at once that \(N=E\).  On passing to the quotient
by \(U_1'\), we are reduced to the case in which \(U\) is étale.
Since \(k\) is algebraically closed, in this case
\(\mathrm H^1(U,H)\) identifies with the ordinary cohomology
group\footnote{Editor's note: see, for example, Proposition
III.6.4.2 of M.\ Demazure and P.\ Gabriel,
\emph{Groupes algébriques I}, Masson \& North-Holland (1970).}
\[
  \mathrm H^1\bigl(U(k),H(k)\bigr),
\]
and is therefore zero, since \(U(k)\) is a finite \(p\)-group and
\(H(k)\) is uniquely \(p\)-divisible.

\SGAStatement{Corollary 6.2.3}{To prove that \(E\) is trivial, it
is enough to prove that it becomes trivial after a finite
separable extension of the base field.  In particular, we may
suppose that \(H\) is diagonalizable.}
\SGARefTarget{sga3:xvii:corollary:6:2:3}

\SGAStatement{Lemma 6.2.4}{In order that every central extension
\(E\) of \(U\) by a diagonalizable \(k\)-group \(H\) be trivial,
it is enough that every central extension of \(U\) by
\(\mathbb G_m\) be trivial.}
\SGARefTarget{sga3:xvii:lemma:6:2:4}

\SGAProof{Proof}
Induction on \(r\) first shows that the hypothesis implies that
\(E\) is trivial when \(H=\mathbb G_m^r\).  In the general case,
\(H\) embeds in \(\mathbb G_m^r\) for a suitable integer \(r\)
(this is immediate by duality); put
\[
  H''=\mathbb G_m^r/H.
\]
We obtain the exact sequence (Appendix I, \SGARefLink{sga3:i:subsection:2-1}{2.1})
\[
  \mathrm Z^1(U,H'')
    \longrightarrow\operatorname{Ext}_{\mathrm{alg}}(U,H)
    \longrightarrow
      \operatorname{Ext}_{\mathrm{alg}}(U,\mathbb G_m^r)=0,
\]
where \(U\) acts trivially on \(H\), \(H''\), and
\(\mathbb G_m^r\).  But
\[
  \mathrm Z^1(U,H'')
    =\operatorname{Hom}_{k\text{-}\mathrm{gr}}(U,H'')=0
\]
by 2.4(ii), and hence
\(\operatorname{Ext}_{\mathrm{alg}}(U,H)=0\).

\SGAProof{Proof of \SGARefLink{sga3:xvii:theorem:6:1:1:item:A}{6.1.1(A)}\SGARefLink{sga3:xvii:theorem:6:1:1:item:A:i}{(i)}}
Since \(E\longrightarrow U\) admits a section, \(E\) defines an
element \(f\) of \(\mathrm H^2(U,H)\) (Appendix I, \SGARefLink{sga3:i:definition:3-1}{3.1}).  We must
prove that \(f\) is zero.  For this it is enough to prove that a
\(2\)-cocycle
\[
  f:U\times_kU\longrightarrow H
\]
is a constant morphism, which follows from the next lemma.

\SGAStatement{Lemma 6.2.5}{Let \(U\) be a smooth connected
unipotent algebraic \(k\)-group, and \(H\) a \(k\)-group of
multiplicative type.  Then every \(k\)-morphism of preschemes
\[
  f:U\longrightarrow H
\]
is constant.}
\SGARefTarget{sga3:xvii:lemma:6:2:5}

\SGAProof{Proof}
We may suppose that \(k\) is algebraically closed.  We argue by
induction on \(\dim U\).  If \(\dim U>0\), then \(U\) has a
composition series (see \SGARefLink{sga3:xvii:lemma:3-9}{3.9})
\[
  1\longrightarrow U'\longrightarrow U
    \xrightarrow{\ \pi\ }U''\longrightarrow1,
  \qquad
  U'\simeq\mathbb G_a,\quad \dim U''<\dim U.
\]
It is enough to show that \(f\) factors through \(U''\).  Since
the graph of the equivalence relation defined by \(\pi\) is
smooth over \(k\), and hence reduced, it is enough to show that
if \(x,y\in U(k)\) have the same image \(z\) in \(U''(k)\), then
\(f(x)=f(y)\).  Now \(\pi^{-1}(z)\) is isomorphic to
\(\mathbb G_a\), so the restriction of \(f\) to
\(\pi^{-1}(z)\) factors through a reduced irreducible component
of \(H\), and hence through a \(k\)-scheme isomorphic to
\(\mathbb G_m^r\).  It remains only to note that every morphism
from \(\mathbb G_a\) to \(\mathbb G_m^r\) is constant, since
every invertible regular function on \(\mathbb G_a\) is
constant.

\SGAProof{Proof of \SGARefLink{sga3:xvii:theorem:6:1:1:item:A}{6.1.1(A)}\SGARefLink{sga3:xvii:theorem:6:1:1:item:A:ii}{(ii)}}
By \SGARefLink{sga3:xvii:lemma:6:2:1}{6.2.1}, \SGARefLink{sga3:xvii:corollary:6:2:3}{6.2.3}, and \SGARefLink{sga3:xvii:lemma:6:2:4}{6.2.4}, we may suppose that
\(H=\mathbb G_m\).  By hypothesis, \(U\) has a composition series
\[
  U\supset U_1\supset\cdots
\]
such that \(U_i/U_{i+1}\) is isomorphic to \(\mathbb G_a\).  Let
\(E_1\) be the inverse image of \(U_1\) in \(E\).  By induction
on \(\dim U\), we may suppose that the extension \(E_1\) is
trivial; let \(U_1'\) be the unique lift of \(U_1\).  Arguing as
in the normalizer argument of \SGARefLink{sga3:xvii:lemma:6:2:2}{6.2.2}, we find that \(U_1'\) is
invariant in \(E\).  After taking the quotient by \(U_1'\), we
are reduced to the case \(U=\mathbb G_a\).  The \(S\)-scheme
\(E\), where \(S=U\), is then a torsor under the \(S\)-group
\(\mathbb G_m\times_kS\), and therefore admits a section, since
\(\operatorname{Pic}(S)=0\) (\SGARefLink{sga3:viii:corollary:4:3}{Exposé VIII, 4.3}).  The extension
\(E\) is then trivial by (A)\SGARefLink{sga3:xvii:theorem:6:1:1:item:A:i}{(i)}.

\SGAProof{Proof of \SGARefLink{sga3:xvii:theorem:6:1:1:item:A}{6.1.1(A)}\SGARefLink{sga3:xvii:theorem:6:1:1:item:A:iii}{(iii)} (\(U\) smooth, \(H\) étale)}
Suppose first that \(U\) is connected.  The group
\(U_{\overline{k}}\) then has a composition series whose
successive quotients are isomorphic to \(\mathbb G_a\), so
\(E_{\overline{k}}\) is trivial by (A)\SGARefLink{sga3:xvii:theorem:6:1:1:item:A:ii}{(ii)}.  Since \(H\) is
étale, the unique lift of \(U_{\overline{k}}\) in
\(E_{\overline{k}}\) is plainly the identity component of
\(E_{\overline{k}}\); this lift is therefore already defined
over \(k\).  In the general case, after quotienting by the
identity component of \(E\), we are reduced to the case where
\(E\) is étale, and then to the case where \(E\) is completely
split (\SGARefLink{sga3:xvii:corollary:6:2:3}{6.2.3}).  Since \(U(k)\) is a \(p\)-group and \(H(k)\) has
order prime to \(p\), we may take the Sylow \(p\)-subgroup of
\(E(k)\) as a lift of \(U(k)\).

\SGAProof{Proof of \SGARefLink{sga3:xvii:theorem:6:1:1:item:A}{6.1.1(A)}\SGARefLink{sga3:xvii:theorem:6:1:1:item:A:iv}{(iv)} (\(U\) smooth, \(k\) perfect)}
If \(U\) is connected, it has a composition series whose
successive quotients are isomorphic to \(\mathbb G_a\)
(\SGARefLink{sga3:xvii:lemma:4:1:2:item:b}{4.1.2(b)}), and we apply (A)\SGARefLink{sga3:xvii:theorem:6:1:1:item:A:ii}{(ii)}.  In the general case, the
preceding argument reduces us to the case in which \(U\) is
étale, and then to the case in which \(U\) is completely split
and \(H\) is diagonalizable (\SGARefLink{sga3:xvii:corollary:6:2:3}{6.2.3}).  Using a characteristic
composition series of \(H\),
\[
  H\supset H^0\supset{}_{F^n}(H)\supset0,
\]
we are reduced to the case in which \(H\) is of one of the
following three types:
\begin{enumerate}
\item[(a)] \(H\) is étale;
\item[(b)] \(H=\mathbb G_m^r\);
\item[(c)] \(H\) is radicial.
\end{enumerate}
In case \SGARefLink{sga3:xvii:appendix:A}{(a)} we apply (A)\SGARefLink{sga3:xvii:theorem:6:1:1:item:A:iii}{(iii)}, and in case \SGARefLink{sga3:xvii:appendix:C}{(c)} we apply \SGARefLink{sga3:xvii:lemma:1:6}{1.6}.
Finally, in case \SGARefLink{sga3:xvii:section:b}{(b)}, Hilbert's \SGARefLink{sga3:viii:corollary:4:5}{Theorem 90} shows that
\(E\longrightarrow U\) admits a section, so it is enough to show
that
\[
  \mathrm H^2(U,\mathbb G_m^r)
    =\mathrm H^2\bigl(U(k),\mathbb G_m^r(k)\bigr)=0.
\]
But \(U(k)\) is a finite \(p\)-group, whereas
\(\mathbb G_m^r(k)\) is uniquely \(p\)-divisible, since \(k\) is
perfect.

\SGASubsection{6.3}{Proof of \SGARefLink{sga3:xvii:theorem:6:1:1:item:B}{6.1.1(B)} and (C)}
\SGARefTarget{sga3:xvii:subsection:6:3}

By \SGARefLink{sga3:xvii:lemma:6:2:1}{6.2.1}, \SGARefLink{sga3:xvii:corollary:6:2:3}{6.2.3}, and \SGARefLink{sga3:xvii:lemma:6:2:4}{6.2.4}, it is enough to prove (B) when
\(H=\mathbb G_m\).  We therefore have an exact sequence
\[
  1\longrightarrow\mathbb G_m\longrightarrow E
    \longrightarrow\boldsymbol\alpha_p\longrightarrow1.
\]
Since \(\mathbb G_m\) is smooth, it gives an exact sequence of
restricted Lie algebras (Appendix II, \SGARefLink{sga3:xvii:proposition:3:2}{3.2})
\begin{equation}
  0\longrightarrow\operatorname{Lie}\mathbb G_m
    \longrightarrow\operatorname{Lie}E
    \longrightarrow\operatorname{Lie}\boldsymbol\alpha_p
    \longrightarrow0.
  \tag{*}\label{sga3:xvii:subsection:6-3:equation:star}
\end{equation}
The group \(\boldsymbol\alpha_p\) has height \(1\), so the
extension \(E\) is trivial if and only if the exact sequence
\((*)\) of restricted Lie algebras splits (Appendix II, \SGARefLink{sga3:ii:subsection:2-2}{2.2}).
The algebra \(\operatorname{Lie}\boldsymbol\alpha_p\) is
generated by an element \(X\ne0\) such that \(X^{(p)}=0\)
(Appendix II, \SGARefLink{sga3:xvii:proposition:2:1}{2.1}).  It is therefore enough to show that \(X\)
lifts to an element \(Z\) of \(\operatorname{Lie}E\) for which
\(Z^{(p)}=0\).  By \SGARefLink{sga3:xvii:lemma:5-8-2:item:ii}{5.8.2(ii)}, there is a unipotent element \(Z\)
of \(\operatorname{Lie}E\) which lifts \(X\).  Since the unipotent
part of \(\operatorname{Lie}E\) plainly has dimension at most
\(1\), it follows necessarily that \(Z^{(p)}=0\).

% END inputs/XVII/tex/components/28_expose_XVII_section6_1_and_proofs_A_B_pages035_037.tex

% BEGIN inputs/XVII/tex/components/29_expose_XVII_proof_C_and_section6_4_examples_page038.tex
\SGAProof{Proof of \SGARefLink{sga3:xvii:theorem:6:1:1:item:C}{6.1.1(C)}}
If \(U_1'\) is a lift in \(E\) of an algebraic subgroup \(U_1\)
of \(U\), then \(U_1'\) is invariant in \(E\), since \(E\) is
assumed commutative, and we may pass to the quotient by \(U_1'\).
Using a composition series of \(U\) (\SGARefLink{sga3:xvii:theorem:3:5:item:ii}{3.5(ii)}), the preceding
remark reduces us to the case in which \(U\) is smooth or equal
to \(\boldsymbol\alpha_p\).  But then \(E\) is trivial by
(A)(iv) and (B).

\SGASubsection{6.4}{Examples of nontrivial extensions of a
unipotent group \(U\) by a group \(H\) of multiplicative type}
\SGARefTarget{sga3:xvii:subsection:6:4}

In view of \SGARefLink{sga3:xvii:theorem:6:1:1:item:A}{6.1.1(A)}(iv), the question arises only in
characteristic \(p>0\).

\begin{enumerate}
\item[(a)] \SGARefTarget{sga3:xvii:subsection:6:4:item:a}
  \(H=\mathbb G_m\), and \(U\) is a nontrivial form of
  \(\mathbb G_a\) (cf.\ Appendix III, \S5).

  Let \(k\) be a nonperfect field of characteristic \(2\), and
  let \(u\in k\) be such that \(u^{1/2}\notin k\).  Consider the
  affine group \(E\) with coordinate ring
  \[
    k\bigl[X,Y,(X^2+uY^2)^{-1}\bigr],
  \]
  whose multiplication is given by the comorphism
  \[
    (X,Y)\longmapsto
      \bigl(XX'+uYY',\,XY'+YX'\bigr).
  \]
  The group \(E\) is smooth, connected, commutative, and of
  dimension \(2\); the subscheme \(Y=0\) defines a subgroup
  \(H\simeq\mathbb G_m\).  The kernel \(K\) of the squaring map
  on \(E\) has equation
  \[
    X^2+uY^2=1,
  \]
  and hence has dimension \(1\).  The group \(K\) contains the
  unipotent radical of \(E\), defined over \(\bar k\), as well as
  the contribution of \(H\), which is isomorphic to
  \(\boldsymbol\mu_2\).  Since \(K\) is reduced over \(\bar k\),
  the unipotent radical of \(E\) is not defined over \(k\), and
  \(U=E/H\) does not lift in \(E\).

  One checks immediately that \(U\) is the form of
  \(\mathbb G_a\) with coordinate ring
  \[
    k[V,W]/(V+uV^2+W^2),
  \]
  the morphism \(E\longrightarrow U\) corresponding to the
  comorphism
  \[
    V\longmapsto\frac{Y^2}{X^2+uY^2},
    \qquad
    W\longmapsto\frac{XY}{X^2+uY^2}.
  \]

\item[(b)] \SGARefTarget{sga3:xvii:subsection:6:4:item:b}
  \(H=\mathbb G_m\), \(U=\mathbb Z/2\mathbb Z\), and \(k\) is
  nonperfect of characteristic \(2\).

  Choose \(k\) and \(u\) as in \SGARefLink{sga3:xvii:appendix:A}{(a)}, and consider the subgroup
  \(E\) of \(\mathrm{GL}_2\) generated by the element
  \[
    X=
    \begin{pmatrix}0&1\\u&0\end{pmatrix}
    =
    \begin{pmatrix}u^{1/2}&0\\0&u^{1/2}\end{pmatrix}
    \begin{pmatrix}0&u^{-1/2}\\u^{1/2}&0\end{pmatrix}.
  \]
  The group \(E\) is an extension of \(\mathbb Z/2\mathbb Z\)
  by \(\mathbb G_m\), but this extension is not trivial because
  the unipotent part of \(X\) is not defined over \(k\).

\item[(c)] \SGARefTarget{sga3:xvii:subsection:6:4:item:c}
  \(H=\boldsymbol\mu_p\), \(U=\boldsymbol\alpha_p\), and \(k\)
  is nonperfect.

  Let \(\mathfrak e\) be the commutative restricted Lie algebra
  generated by two elements \(X\) and \(Y\) such that
  \[
    X^{(p)}=X,\qquad Y^{(p)}=aX.
  \]
  By Appendix II, \SGARefLink{sga3:ii:subsection:2-2}{2.2}, \(\mathfrak e\) is the restricted Lie
  algebra of an algebraic group \(E\) which is an extension of
  \(\boldsymbol\alpha_p\) by \(\boldsymbol\mu_p\).  This
  extension is trivial if and only if there exists \(b\in k\)
  such that \(b^p=a\), for then
  \[
    (Y-bX)^{(p)}=0.
  \]

\item[(d)] \SGARefTarget{sga3:xvii:subsection:6:4:item:d}
  \(H=\boldsymbol\mu_2\), \(U=\boldsymbol\alpha_2\times
  \boldsymbol\alpha_2\), \(E\) noncommutative, and \(k\) a field
  of characteristic \(2\).

  Consider the special linear group \(\mathrm{SL}_{2,k}\), and
  put \(E=F(\mathrm{SL}_{2,k})\).  The group \(E\) is a radicial
  group of height \(1\), whose Lie algebra is generated by three
  elements \(X,Y,Z\) satisfying
  \[
    [X,Y]=Z,\qquad [X,Z]=[Y,Z]=0,
  \]
  \[
    X^{(p)}=Y^{(p)}=0,\qquad Z^{(p)}=Z.
  \]
\end{enumerate}

% END inputs/XVII/tex/components/29_expose_XVII_proof_C_and_section6_4_examples_page038.tex

% BEGIN inputs/XVII/tex/components/30_expose_XVII_section6_4_completion_and_proposition7_1_1_page039.tex
% Authority: Expo17.pdf, local p.39 / combined p.1001 /
% printed pp.611--612.  This component completes 6.4(d), begins
% section 7, and carries the proof of Proposition 7.1.1 through
% part ii(a).  It stops before part ii(b) on authority local p.40.

Consequently, \(E\) is a central extension of
\(U\simeq\boldsymbol\alpha_2\times\boldsymbol\alpha_2\) by
\(\boldsymbol\mu_2\).  Each factor \(\boldsymbol\alpha_2\) of
\(U\) lifts uniquely in \(E\), but \(U\) itself does not lift in
\(E\), since \([X,Y]\ne0\).

\SGASection{7}{Nilpotent Affine Algebraic Groups}
\SGARefTarget{sga3:xvii:section:7}

\SGASubsection{7.1}{Extensions of Groups of Multiplicative Type}
\SGARefTarget{sga3:xvii:subsection:7:1}

\SGAStatement{Proposition 7.1.1}{Let \(S\) be a prescheme, let
\(H\) and \(K\) be two \(S\)-group preschemes of multiplicative
type and finite type, and let \(E\) be an \(S\)-group prescheme
which is an extension of \(K\) by \(H\) (that is, \(K\) is the
quotient of \(E\) by \(H\) for the fpqc topology).  Then \(E\) is
of multiplicative type in either of the following two cases:
\begin{enumerate}
\item[(a)] \SGARefTarget{sga3:xvii:proposition:7:1:1:item:a}
  \(E\) is commutative.
\item[(b)] \SGARefTarget{sga3:xvii:proposition:7:1:1:item:b}
  The fibres of \(K\) are connected.
\end{enumerate}}
\SGARefTarget{sga3:xvii:proposition:7:1:1}

\SGAProof{Proof}
\begin{enumerate}
\item[\SGAInlineRefTarget{sga3:xvii:proposition:7:1:1:item:i}(i)] Suppose first that \(S\) is the spectrum of a field
  \(k\).  The assertion is local for the fpqc topology, so we may
  suppose that \(k\) is algebraically closed and hence that \(H\)
  and \(K\) are diagonalizable.  Notice that \(K\) acts trivially
  on \(H\) by inner automorphisms: this is clear in case \SGARefLink{sga3:xvii:appendix:A}{(a)}, and
  follows from \SGARefLink{sga3:xvii:lemma:6:2:1}{6.2.1} in case \SGARefLink{sga3:xvii:section:b}{(b)}.  Induction on the length of a
  suitable composition series of \(K\) reduces us to the case
  where \(K\) is of one of the following three types:
  \[
    \text{(a) }K=\mathbb G_m,\qquad
    \text{(b) }K=\boldsymbol\mu_p,\qquad
    \text{(c) }K=\boldsymbol\mu_q,\quad(q,p)=1,
  \]
  and in case \SGARefLink{sga3:xvii:proposition:7:1:1:item:c}{(c)}, \(E\) is commutative.  Now embed \(H\) in
  \(\mathbb G_m^r\).  It follows that \(E\) embeds in an extension
  of \(K\) by \(\mathbb G_m^r\), and we may therefore suppose that
  \(H=\mathbb G_m^r\).

  \begin{enumerate}
  \item[(a)] If \(K=\mathbb G_m\), then \(E\) is a smooth,
    connected, affine algebraic group (Exposé VI\textsubscript{B},
    \SGARefLink{prop:sga3-VIB-9.2}{9.2}) of unipotent rank zero, and is therefore a torus.

  \item[(b)] Let \(K=\boldsymbol\mu_p\).  In this case \(E\) is a
    trivial extension.  Indeed, since \(H\) is smooth, the
    canonical morphism
    \[
      \operatorname{Lie}E\longrightarrow
      \operatorname{Lie}\boldsymbol\mu_p
    \]
    is surjective (Appendix II, \SGARefLink{sga3:xvii:proposition:3:2}{3.2}), and the assertion follows
    from \SGARefLink{sga3:xvii:lemma:5-8-2:item:i}{5.8.2(i)}, in view of Appendix II, \SGARefLink{sga3:ii:subsection:2-2}{2.2}.

  \item[\SGAInlineRefTarget{sga3:xvii:proposition:7:1:1:item:c}(c)] Let \(K=\boldsymbol\mu_q\), where \((q,p)=1\), and
    suppose that \(E\) is commutative.  Again the extension \(E\)
    is trivial.  Indeed, let \(x\in E(\bar k)\) lift a generator
    \(\bar x\) of \(\boldsymbol\mu_q(\bar k)\).  The element
    \(x^q\) belongs to \(H(\bar k)\), and hence has the form
    \(y^q\), with \(y\in H(\bar k)\) (notice that
    \(\mathbb G_m^r(\bar k)\) is \(q\)-divisible).  Since \(E\)
    is commutative, \(y^{-1}x\) is a lift of \(\bar x\) having
    order \(q\).
  \end{enumerate}

\item[(ii)] Consider now the general case.  The groups \(H\) and
  \(K\) are flat, affine, and of finite presentation over \(S\)
  (\SGARefLink{sga3:ix:proposition:2:1}{Exposé IX, 2.1}), and consequently so is \(E\) (Exposé
  VI\textsubscript{B}, \SGARefLink{prop:sga3-VIB-9.2}{9.2}).  Using the general technique of
  Exposé VI\textsubscript{B}, \S10, we reduce to the case where
  \(S\) is noetherian.  To prove that \(E\) is of multiplicative
  type, it then suffices to prove that \(E\) is commutative and
  that \({}_nE\) is finite over \(S\) for every \(n\) (Exposé X,
  4.8(b)).

  \begin{enumerate}
  \item[(a)] \(E\) is commutative.  We must verify that the
    morphism
    \[
      E\times_S E\longrightarrow E,\qquad
      (x,y)\longmapsto[x,y]=xyx^{-1}y^{-1}
    \]
    factors through the unit section of \(E\), and it is enough
    to verify this when \(S\) is the spectrum of an artinian
    local ring.  But then \(E\) is of multiplicative type by \SGARefLink{sga3:xvii:proposition:7:1:1:item:i}{(i)}
    and \SGARefLink{sga3:x:corollary:2:3}{Exposé X, 2.3}.

% END inputs/XVII/tex/components/30_expose_XVII_section6_4_completion_and_proposition7_1_1_page039.tex

% BEGIN inputs/XVII/tex/components/31_expose_XVII_proposition7_1_1_completion_and_theorem7_2_1_page040.tex
% Authority: Expo17.pdf, local p.40 / combined p.1002 /
% printed pp.613--614.  This component completes Proposition 7.1.1,
% states Theorem 7.2.1, and carries its proof into the descent
% argument.  The following component completes that argument.

  \item[(b)] \SGARefTarget{sga3:xvii:proposition:7:1:1:proof:ii:item:b}
    The group \({}_nE\) is finite over \(S\).  Indeed, there is an
    exact sequence
    \[
      0\longrightarrow{}_nH\longrightarrow{}_nE
      \longrightarrow{}_nK\xrightarrow{u}H_n,
    \]
    where \(H_n\) is the cokernel of the \(n\)-th power map on
    \(H\).  We know that \(H_n\) is of multiplicative type (\SGARefLink{sga3:ix:proposition:2:7}{Exposé
    IX, 2.7}), and hence separated, and that \({}_nH\) and
    \({}_nK\) are finite over \(S\).  The group \(\ker u\) is a
    closed subgroup of \({}_nK\), and hence is finite over \(S\);
    consequently \({}_nE\) is finite over \(S\), as an extension
    of one finite group by another (Exposé VI\textsubscript{B},
    \SGARefLink{prop:sga3-VIB-9.2}{9.2}).
  \end{enumerate}
\end{enumerate}

\SGASubsection{7.2}{Structure of Commutative Affine Algebraic Groups}
\SGARefTarget{sga3:xvii:subsection:7:2}

\SGAStatement{Theorem 7.2.1}{Let \(k\) be a field and \(G\) a
commutative affine algebraic \(k\)-group.  Then:
\begin{enumerate}
\item[(a)] \SGARefTarget{sga3:xvii:theorem:7:2:1:item:a}
  \(G\) contains a largest subgroup \(M\) of multiplicative type.
  The group \(M\) is characteristic in \(G\), the quotient \(G/M\)
  is unipotent, and the formation of \(M\) commutes with extension
  of the ground field \(k\).
\item[(b)] \SGARefTarget{sga3:xvii:theorem:7:2:1:item:b}
  If \(k\) is perfect, \(G\) is the direct product of \(M\) and a
  unipotent algebraic subgroup \(U\), uniquely so.
\end{enumerate}}
\SGARefTarget{sga3:xvii:theorem:7:2:1}

\SGAProof{Proof}
\begin{enumerate}
\item[(i)] Suppose that \(k\) is algebraically closed.  When \(G\)
  is smooth, \SGARefLink{sga3:xvii:theorem:7:2:1:item:b}{7.2.1(b)} is well known (BIBLE~4, Theorem~4).  If
  \(G\) is radicial of height \(1\), the decomposition of
  \(\operatorname{Lie}G\) described in \SGARefLink{sga3:xvii:lemma:4-2-2}{4.2.2} corresponds, in view
  of \SGARefLink{sga3:xvii:proposition:4:3:1:item:v}{4.3.1(v)} and Appendix II, \SGARefLink{sga3:ii:subsection:2-2}{2.2}, to a decomposition of \(G\) of
  the kind asserted in \SGARefLink{sga3:xvii:theorem:7:2:1:item:b}{7.2.1(b)}.  In general, \(G\) has a
  composition series whose successive quotients are smooth or
  radicial of height \(1\) (Appendix II, \SGARefLink{sga3:xvii:lemma:3-1}{3.1}).  To prove \SGARefLink{sga3:xvii:theorem:7:2:1:item:b}{7.2.1(b)},
  it therefore suffices to observe the following.  Suppose there
  is an exact sequence of commutative algebraic groups
  \[
    0\longrightarrow G'\longrightarrow G\longrightarrow
    G''\longrightarrow0,
  \]
  in which \(G'\), respectively \(G''\), is the product of a group
  of multiplicative type and a unipotent group,
  \[
    G'=M'\cdot U',\qquad G''=M''\cdot U''.
  \]
  Then the same is true of \(G\).  Indeed, consider the exact
  sequence
  \[
    0\longrightarrow G'/U'\longrightarrow G/U'
      \longrightarrow G''\longrightarrow0.
  \]
  By \SGARefLink{sga3:xvii:proposition:7:1:1:item:a}{7.1.1(a)}, the inverse image of \(M''\) in \(G/U'\) is a
  subgroup \(M_1\) of multiplicative type.  The group \(M_1\)
  lifts to a subgroup \(M\) of \(G\) by \SGARefLink{sga3:xvii:theorem:5:1:1:item:i}{5.1.1(i)}\SGARefLink{sga3:xvii:appendix:A}{(a)}.  Likewise,
  now using \SGARefLink{sga3:xvii:theorem:6:1:1:item:C}{6.1.1(C)}, one proves that there is a unipotent
  subgroup \(U\) of \(G\) which is an extension of \(U''\) by
  \(U'\).  Clearly \(G=M\cdot U\).

\item[(ii)] Let \(k\) now be arbitrary.  By \SGARefLink{sga3:xvii:subsection:6:4:item:i}{(i)},
  \(G_{\bar k}\) is the direct product of a group
  \(M_{\bar k}\) of multiplicative type and a unipotent group
  \(U_{\bar k}\).  Put
  \[
    S=\bar k\times_k\bar k,
  \]
  and let \(M_1\) and \(M_2\) be the two inverse images of
  \(M_{\bar k}\) under the two projections
  \(S\rightrightarrows\bar k\).  The group
  \[
    G_S/M_2=(G/M_{\bar k})_S
  \]
  has unipotent fibres; hence the image of \(M_1\) in \(G_S/M_2\)
  is trivial by \SGARefLink{sga3:xvii:proposition:2:4:item:i}{2.4(i)}, and \(M_1\subset M_2\).  The converse
  inclusion follows in the same way, so \(M_1=M_2\).  By fpqc
  descent, \(M_{\bar k}\) therefore descends to an algebraic
  subgroup \(M\) of \(G\).  It is clear that \(M\) is of
  multiplicative type, that \(G/M\) is unipotent, and that the
  formation of \(M\) is compatible with every extension of
  \(k\).  For every \(k\)-prescheme \(S\), every subgroup \(H\) of
  multiplicative type in \(G_S\) is contained in \(M_S\).
  Indeed, its image in the group with unipotent fibres
  \[
    (G/M)_S=G_S/M_S
  \]
  is trivial by \SGARefLink{sga3:xvii:lemma:2:5}{2.5}.

% END inputs/XVII/tex/components/31_expose_XVII_proposition7_1_1_completion_and_theorem7_2_1_page040.tex

% BEGIN inputs/XVII/tex/components/32_expose_XVII_theorem7_2_1_completion_through_lemma7_3_2_page041.tex
% Authority: Expo17.pdf, local p.41 / combined p.1003 /
% printed pp.615--616.  This component completes Theorem 7.2.1,
% gives Remark 7.2.2, states Theorem 7.3.1 and Lemma 7.3.2, and
% carries the lemma's proof through the construction in case (i).
% It stops before the concluding sentence about the generators of N.

  Taking \(H\), in particular, to be the transform of \(M_S\)
  under an automorphism of \(G_S\), we conclude that \(M\) is
  characteristic in \(G\).

  Finally, if \(k\) is perfect, \(G/M\) lifts in \(G\) to a
  unipotent group, uniquely so by \SGARefLink{sga3:xvii:theorem:6:1:1:item:C}{6.1.1(C)}.
\end{enumerate}

\SGAStatement{Remark 7.2.2}{%
\begin{enumerate}
\item[\SGAInlineRefTarget{sga3:xvii:subsection:6:4:item:i}(i)] \SGARefTarget{sga3:xvii:remark:7:2:2:item:i}
  If \(k\) is not perfect, the unipotent component of
  \(G_{\bar k}\) need not be defined over \(k\), as \SGARefLink{sga3:xvii:subsection:6:4:item:a}{Example 6.4(a)}
  shows.

\item[(ii)] \SGARefTarget{sga3:xvii:remark:7:2:2:item:ii}
  In contrast to the multiplicative-type component \(M\), the
  unipotent component \(U\) is not in general characteristic in
  \(G\), whatever the characteristic of \(k\).  Of course, the
  uniqueness of the decomposition in \SGARefLink{sga3:xvii:theorem:7:2:1:item:b}{7.2.1(b)} implies that \(U\)
  is invariant under every \(k\)-automorphism of \(G\).  But if
  \(U\) and \(M\) are such that there exist a \(k\)-prescheme
  \(S\) and a nonzero \(S\)-homomorphism
  \[
    h:U_S\longrightarrow M_S
  \]
  (cf.\ \SGARefLink{sga3:xvii:remark:2-6}{2.6}), then one obtains an \(S\)-automorphism of \(G_S\),
  \[
    (u,m)\longmapsto\bigl(u,h(u)+m\bigr),
  \]
  which does not leave \(U_S\) invariant.

\item[(iii)] \SGARefTarget{sga3:xvii:remark:7:2:2:item:iii}
  If \(G\) is finite over \(k\), then under Cartier duality (\SGARefLink{sga3:xvii:remark:2-6}{2.6})
  the quotient \(G/M\) corresponds to the connected component of
  the identity in the dual \(D(G)\) of \(G\).
\end{enumerate}}
\SGARefTarget{sga3:xvii:remark:7:2:2}

\SGASubsection{7.3}{Structure of Nilpotent Affine Algebraic Groups}
\SGARefTarget{sga3:xvii:subsection:7:3}

\SGAStatement{Theorem 7.3.1}{Let \(k\) be a field and \(G\) a
nilpotent (Exposé VI\textsubscript{B}, \S8), affine, connected
algebraic \(k\)-group.  Then \(G\) has a largest subgroup \(M\) of
multiplicative type.  The group \(M\) is central and
characteristic, and \(G/M\) is a unipotent algebraic group.}
\SGARefTarget{sga3:xvii:theorem:7:3:1}

Let \(Z\) be the centre of \(G\), and let \(M\) be the largest
subgroup of multiplicative type in \(Z\), as supplied by \SGARefLink{sga3:xvii:theorem:7:2:1}{7.2.1}.
Since \(Z\) is characteristic in \(G\), and \(M\) is
characteristic in \(Z\), the group \(M\) is characteristic in
\(G\).  It remains to show that \(G/M\) is unipotent.  By
induction on the length of the ascending central series of \(G\),
this follows from the following more general lemma.

\SGAStatement{Lemma 7.3.2 (Rosenlicht)}{Let \(G\) be a connected
algebraic \(k\)-group and let \(Z\) be its centre.  Then the centre
\(Z'\) of \(G'=G/Z\) is unipotent.}
\SGARefTarget{sga3:xvii:lemma:7:3:2}

\SGAProof{Proof}
We may suppose that \(k\) is algebraically closed.  It is then
enough to show that \(Z'\) contains no subgroup isomorphic to
\(\boldsymbol\mu_\ell\), for any prime number \(\ell\), by
\SGARefLink{sga3:xvii:theorem:4:6:1:item:vi}{4.6.1(vi)}.  Let \(\boldsymbol\mu_\ell\) be such a subgroup of
\(Z'\), and let \(N\) be its inverse image in \(G\).  Since
\(\boldsymbol\mu_\ell\) is central in \(G'\), the group \(N\) is
invariant in \(G\).

\begin{enumerate}
\item[(i)] \SGARefTarget{sga3:xvii:lemma:7:3:2:proof:item:i}
  Suppose that \((\ell,p)=1\).  We can find an element
  \(x\in N(k)\) and an integer \(n\) with the following
  properties:
  \begin{enumerate}
  \item[(a)] \SGARefTarget{sga3:xvii:lemma:7:3:2:proof:item:i:a}
    \(x\) lifts a generator \(\bar x\) of
    \(\boldsymbol\mu_\ell\);
  \item[(b)] \SGARefTarget{sga3:xvii:lemma:7:3:2:proof:item:i:b}
    \(x^{\ell^n}\in Z^0(k)\).
  \end{enumerate}
  Indeed, it suffices to choose a lift of \(\bar x\) whose image
  in \((N/Z^0)(k)\) belongs to the \(\ell\)-Sylow subgroup.
  Raising to the \(\ell\)-th power in the commutative group
  \(Z^0\) is an étale morphism, so \(Z^0(k)\) is
  \(\ell\)-divisible.  Consequently, after multiplying \(x\) by
  an element of \(Z^0(k)\), we may suppose that
  \(x^{\ell^n}=0\).

% END inputs/XVII/tex/components/32_expose_XVII_theorem7_2_1_completion_through_lemma7_3_2_page041.tex

% BEGIN inputs/XVII/tex/components/33_expose_XVII_lemma7_3_2_completion_and_appendixA_opening_page042.tex
  The group \(N\) is then generated by two commuting commutative
  groups, namely \(Z\) and the group generated by \(x\), and is
  therefore commutative.  The group \({}_{\ell^n}N\) is of
  multiplicative type and is characteristic in \(N\); hence it is
  invariant in \(G\), and consequently central because \(G\) is
  connected (\SGARefLink{sga3:ix:corollary:5:5}{Exposé IX, 5.5}).  Thus \({}_{\ell^n}N\) is contained
  in \(Z\), contradicting the fact that its image in \(G'\)
  contains \(\boldsymbol\mu_\ell\).

\item[(ii)] \SGARefTarget{sga3:xvii:lemma:7:3:2:proof:item:ii}
  Suppose that \(\ell=p\).  There is then an integer \(n\) such
  that the image of \(F_nN\) in \(G'\) contains
  \(\boldsymbol\mu_p\), and hence an exact sequence
  \[
    1\longrightarrow K\longrightarrow F_nN
      \longrightarrow\boldsymbol\mu_p\longrightarrow1. \tag{*}\label{sga3:xvii:lemma:7-3-2:equation:star}
  \]
  The group \(K\) is contained in \(Z\), and is therefore
  commutative.  It follows from \SGARefLink{sga3:xvii:theorem:7:2:1}{7.2.1}, \SGARefLink{sga3:xvii:proposition:7:1:1:proof:ii:item:b}{7.1.1(b)}, and \SGARefLink{sga3:xvii:lemma:5:5:1}{5.5.1} that
  there is a subgroup of multiplicative type in \(F_nN\) whose
  image in the quotient by \(K\) is \(\boldsymbol\mu_p\).  As in
  \SGARefLink{sga3:xvii:subsection:6:4:item:i}{(i)}, we deduce that \(F_nN\) is commutative.  Its
  multiplicative-type component (\SGARefLink{sga3:xvii:theorem:7:2:1}{7.2.1}) is characteristic in
  \(F_nN\), hence invariant in \(G\), and therefore central
  (\SGARefLink{sga3:ix:corollary:5:5}{Exposé IX, 5.5}).  Since its image in \(G'\) contains
  \(\boldsymbol\mu_p\), this is a contradiction.
\end{enumerate}

\SGASection{A}{Appendix I. Hochschild Cohomology and Extensions of
Algebraic Groups}
\SGARefTarget{sga3:xvii:appendix:A}

\SGANextSectionTarget{sga3:xvii:subsection:a-1}\SGASubsection{A.1}{Definition of the Cohomology Groups}
\SGARefTarget{sga3:xvii:appendix:A:subsection:1}

Let \(k\) be a field, let \(G\) be an algebraic \(k\)-group, and
let \((\mathrm{Ab})_G\) be the abelian category of commutative
algebraic \(k\)-groups on which \(G\) acts.  If
\(A\in\operatorname{Ob}(\mathrm{Ab})_G\), the functor
\[
  h_A:(\operatorname{Sch}/k)^\circ\longrightarrow
      \operatorname{Ens}
\]
canonically defined by \(A\) is a \(G\)-\(\mathbb Z\)-module in
the sense of \SGARefLink{sga3:i:definition:3-2}{Exposé I, 3.2}.  We may therefore consider the
standard complex \(C^\bullet(G,A)\) of algebraic cochains of \(G\)
with values in \(A\) (\SGARefLink{sga3:i:subsection:5-1}{Exposé I, 5.1}), together with the group
\(Z^i(G,A)\) of \(i\)-cocycles, the group \(B^i(G,A)\) of
\(i\)-coboundaries, and the \(i\)-th cohomology group
\(H^i(G,A)\).  As usual, \(H^0(G,A)\) is identified with the
group \(A^G(k)\), where \(A^G\) is the \(k\)-functor of invariants
of \(A\) under \(G\).  The group \(H^1(G,A)\) classifies the
trivial principal homogeneous spaces under \(A\) on which \(G\)
acts.

In general, the functor
\[
  A\longmapsto H^\bullet(G,A)
\]
from \((\mathrm{Ab})_G\) to \(\mathrm{Ab}\) is not a
cohomological functor.  Nevertheless, we have the following
proposition.

\SGAStatement{\SGAInlineRefTarget{sga3:xvii:proposition:a-1-1}Proposition A.1.1}{Let
\[
  0\longrightarrow A'\xrightarrow{u}A\xrightarrow{v}A''
    \longrightarrow0
\]
be an exact sequence in \((\mathrm{Ab})_G\).  Then:
\begin{enumerate}
\item[(a)] \SGARefTarget{sga3:xvii:proposition:A:1:1:item:a}
  if \(v\) has a section---that is, if there is a \(k\)-morphism
  of preschemes \(s:A''\to A\) such that
  \(v s=1_{A''}\)---then there is the usual exact cohomology
  sequence
  \[
    \cdots\longrightarrow H^i(G,A)\longrightarrow H^i(G,A'')
      \xrightarrow{d}H^{i+1}(G,A')\longrightarrow\cdots;
  \]
\item[(b)] \SGARefTarget{sga3:xvii:proposition:A:1:1:item:b}
  if \(A(k)\to A''(k)\) is surjective, then the sequence
  \[
    0\longrightarrow A'^G(k)\longrightarrow A^G(k)
      \longrightarrow A''^G(k)\xrightarrow{d}H^1(G,A')
      \longrightarrow H^1(G,A)\longrightarrow H^1(G,A'')
      \tag{1}\label{sga3:xvii:proposition:a-1-1:equation:1}
  \]
  is exact.
\end{enumerate}}
\SGARefTarget{sga3:xvii:proposition:A:1:1}

\SGAProof{Proof}
For \SGARefLink{sga3:xvii:appendix:A}{(a)}, observe that the existence of a section implies the
exactness of the sequence of complexes
\[
  0\longrightarrow C^\bullet(G,A')\longrightarrow C^\bullet(G,A)
    \longrightarrow C^\bullet(G,A'')\longrightarrow0.
\]

% END inputs/XVII/tex/components/33_expose_XVII_lemma7_3_2_completion_and_appendixA_opening_page042.tex

% BEGIN inputs/XVII/tex/components/34_expose_XVII_appendixA_1_completion_and_A_2_page043.tex
For (b), if \(x''\in A''{}^G(k)\), its inverse image in \(A\)
is a trivial principal homogeneous space under \(A'\) (because
\(A(k)\to A''(k)\) is assumed to be surjective) on which \(G\)
acts, and therefore defines an element
\(d(x'')\in H^1(G,A')\).  The exactness of the sequence (1) is
then immediate.

\SGANextSectionTarget{sga3:xvii:subsection:a-2}\SGASubsection{A.2}{The Group
\texorpdfstring{\(\operatorname{Ext}_{\mathrm{alg}}(G,A)\)}
{Ext alg (G,A)}}
\SGARefTarget{sga3:xvii:appendix:A:subsection:2}

Let \(A\) and \(G\) be two algebraic \(k\)-groups, and let \(E\)
and \(E'\) be two algebraic \(k\)-groups that are extensions of
\(G\) by \(A\).  These two extensions are said to be
\emph{isomorphic} if there is a \(k\)-morphism of groups
\(u:E\to E'\) that makes the following diagram commutative:
\SGARefTarget{sga3:xvii:appendix:A:diagram:A:2:isomorphic-extensions}
\[
\begin{tikzcd}[column sep=large,row sep=large]
  {} & E \arrow[dr] \arrow[d,"u"] & {} \\
  A \arrow[ur] \arrow[dr] & {} & G \\
  {} & E' \arrow[ur] & {}
\end{tikzcd}
\qquad .
\]

The group \(E\) acts on \(A\) by inner automorphisms.  If \(A\)
is commutative, this action factors through \(G\), so that \(G\)
acts on \(A\).  Conversely, if
\(A\in\operatorname{Ob}(\mathrm{Ab})_G\), we denote by
\(\operatorname{Ext}_{\mathrm{alg}}(G,A)\) the set of classes of
algebraic extensions \(E\) of \(G\) by \(A\) for which the
action of \(G\) on \(A\) defined by \(E\) agrees with the action
coming from the structure of \(A\) as an object of
\((\mathrm{Ab})_G\).

The set \(\operatorname{Ext}_{\mathrm{alg}}(G,A)\) is naturally a
bifunctor, covariant in \(A\) and contravariant in \(G\).  More
precisely:
\begin{enumerate}
\item[(a)] \SGARefTarget{sga3:xvii:appendix:A:subsection:2:item:a}
  if \(f:A\to B\) is a morphism in \((\mathrm{Ab})_G\), and if
  \(E\) represents an element of
  \(\operatorname{Ext}_{\mathrm{alg}}(G,A)\), then
  \(f_*(E)\in\operatorname{Ext}_{\mathrm{alg}}(G,B)\) is defined
  to be the class of the extension of \(G\) by \(B\) equal to
  the quotient of the semidirect product \(B\mathbin{\cdot}E\)
  (where \(E\) acts on \(B\) through \(G\)) by the algebraic
  subgroup that is the image of \(A\) under the morphism
  \((f,i)\).  This quotient is representable by Exposé VI\textsubscript{A},
  \S5, and there is therefore a commutative diagram
  \SGARefTarget{sga3:xvii:appendix:A:diagram:A:2:pushout}
  \[
  \begin{tikzcd}[column sep=large,row sep=large]
    1 \arrow[r] &
    A \arrow[r,"i"] \arrow[d,"f"'] &
    E \arrow[r] \arrow[d] &
    G \arrow[r] \arrow[d,"1_G"] &
    1 \\
    1 \arrow[r] &
    B \arrow[r] &
    f_*(E) \arrow[r] &
    G \arrow[r] &
    1
  \end{tikzcd}
  \qquad .
  \]

\item[(b)] \SGARefTarget{sga3:xvii:appendix:A:subsection:2:item:b}
  if \(g:H\to G\) is a \(k\)-morphism of algebraic
  \(k\)-groups, and if \(E\) is an extension of \(G\) by \(A\),
  the fiber product \(E\times_G H\) is naturally an extension
  of \(H\) by \(A\), denoted by \(g^*(E)\).  Thus there is a
  commutative diagram
  \SGARefTarget{sga3:xvii:appendix:A:diagram:A:2:pullback}
  \[
  \begin{tikzcd}[column sep=large,row sep=large]
    1 \arrow[r] &
    A \arrow[r] \arrow[d,"1_A"'] &
    g^*(E) \arrow[r] \arrow[d] &
    H \arrow[r] \arrow[d,"g"] &
    1 \\
    1 \arrow[r] &
    A \arrow[r] &
    E \arrow[r] &
    G \arrow[r] &
    1
  \end{tikzcd}
  \qquad .
  \]
\end{enumerate}

By adapting the proofs given in J.-P. Serre,
\emph{Groupes algébriques et corps de classe}, Chapter VII, one
endows \(\operatorname{Ext}_{\mathrm{alg}}(G,A)\) with a natural
abelian-group structure, functorial in \(A\) and \(G\).

% END inputs/XVII/tex/components/34_expose_XVII_appendixA_1_completion_and_A_2_page043.tex

% BEGIN inputs/XVII/tex/components/35_expose_XVII_appendixA_sections2_3_and_appendixB_opening_page044.tex
\SGAStatement{\SGAInlineRefTarget{sga3:xvii:proposition:a-2-1}Proposition A.2.1}{Let
\[
  0\longrightarrow A'\longrightarrow A\longrightarrow A''
    \longrightarrow0
\]
be an exact sequence in \((\mathrm{Ab})_G\).  Then:
\begin{enumerate}
\item[(a)] \SGARefTarget{sga3:xvii:proposition:A:2:1:item:a}
  there is a canonical exact sequence of abelian groups
  \[
    Z^1(G,A)\longrightarrow Z^1(G,A'')
      \xrightarrow{d}\operatorname{Ext}_{\mathrm{alg}}(G,A')
      \longrightarrow\operatorname{Ext}_{\mathrm{alg}}(G,A)
      \longrightarrow\operatorname{Ext}_{\mathrm{alg}}(G,A'');
  \]
\item[(b)] \SGARefTarget{sga3:xvii:proposition:A:2:1:item:b}
  if \(A(k)\to A''(k)\) is surjective, one deduces from \SGARefLink{sga3:xvii:appendix:A}{(a)} the
  exact sequence
  \[
  \tag{2}\label{sga3:xvii:proposition:a-2-1:equation:2}
    H^1(G,A)\longrightarrow H^1(G,A'')
      \xrightarrow{d}\operatorname{Ext}_{\mathrm{alg}}(G,A')
      \longrightarrow\operatorname{Ext}_{\mathrm{alg}}(G,A)
      \longrightarrow\operatorname{Ext}_{\mathrm{alg}}(G,A'').
  \]
\end{enumerate}}
\SGARefTarget{sga3:xvii:proposition:A:2:1}

The exact sequence in \SGARefLink{sga3:xvii:appendix:A}{(a)} generalizes the usual exact
\(\operatorname{Hom}(\,\cdot\,,\,\cdot\,)\) sequence valid for
commutative extensions (\emph{loc.\ cit.}), and is proved in the
same way.  We recall only the definition of the coboundary
\[
  d:Z^1(G,A'')\longrightarrow
    \operatorname{Ext}_{\mathrm{alg}}(G,A').
\]
For this purpose, consider the extension
\[
  1\longrightarrow A'\longrightarrow A\mathbin{\cdot}G
    \longrightarrow A''\mathbin{\cdot}G\longrightarrow1,
\]
obtained in the evident way from the exact sequence
\(0\to A'\to A\to A''\to0\).  If \(u\in Z^1(G,A'')\), then
\(u\) defines in the usual way a section homomorphism
\[
  u:G\longrightarrow A''\mathbin{\cdot}G.
\]
One then has
\[
  d(u)=u^*(A\mathbin{\cdot}G).
\]

\SGANextSectionTarget{sga3:xvii:subsection:a-3}\SGASubsection{A.3}{Comparison of
\texorpdfstring{\(H^2(G,A)\) and
\(\operatorname{Ext}_{\mathrm{alg}}(G,A)\)}
{H2(G,A) and Ext alg(G,A)}}
\SGARefTarget{sga3:xvii:appendix:A:subsection:3}

It is well known that, for abstract groups, there is a functorial
isomorphism between the abelian groups \(H^2(G,A)\) and
\(\operatorname{Ext}(G,A)\).  Likewise in the present setting,
if \(A\) is an object of \((\mathrm{Ab})_G\), one can associate
with every \(2\)-cocycle \(u\in Z^2(G,A)\) an algebraic-group
structure on the prescheme \(A\times_kG\), making it an element
of \(\operatorname{Ext}_{\mathrm{alg}}(G,A)\).  Moreover, this
extension is trivial if and only if \(u\in B^2(G,A)\) (see
\SGARefLink{sga3:iii:lemma:1-2-2}{Exposé III, 1.2.2}).  Recall that the composition law on
\(A\times G\) is defined by
\[
  (a,g)(a',g')=a+{}^g a'+u(g,g').
\]

It is clear that the extensions of \(G\) by \(A\) obtained in
this way are not arbitrary, since they possess a section.
Conversely, however, if
\(E\in\operatorname{Ext}_{\mathrm{alg}}(G,A)\) possesses a
section \(s\), then \(E\) is isomorphic to the extension of
\(G\) by \(A\) associated with the \(2\)-cocycle \(u\) such that
\[
  u(g,g')=s(g)s(g')s(gg')^{-1}.
\]
We thus obtain the following proposition.

\SGAStatement{\SGAInlineRefTarget{sga3:xvii:proposition:a-3-1}Proposition A.3.1}{There is a functorial
isomorphism between the bifunctors with values in abelian
groups
\[
  (G,A)\longmapsto H^2(G,A)
  \qquad\text{and}\qquad
  (G,A)\longmapsto\operatorname{Ext}_s(G,A),
\]
where \(\operatorname{Ext}_s(G,A)\) denotes the subgroup of
\(\operatorname{Ext}_{\mathrm{alg}}(G,A)\) formed by the classes
of extensions of \(G\) by \(A\) that possess a section.}
\SGARefTarget{sga3:xvii:proposition:A:3:1}

\SGANextSectionTarget{sga3:xvii:section:b}\SGASection{B}{Appendix II. Review and Supplements on Radicial
Groups}
\SGARefTarget{sga3:xvii:appendix:B}

Let \(p>1\) be a prime number, and let \(S\) be an
\(\mathbb F_p\)-prescheme.

% END inputs/XVII/tex/components/35_expose_XVII_appendixA_sections2_3_and_appendixB_opening_page044.tex

% BEGIN inputs/XVII/tex/components/36_expose_XVII_appendixB_pages045_046_and_appendixC_opening.tex
\SGASubsection{B.1}{The Frobenius morphism}
\SGARefTarget{sga3:xvii:appendix:B:subsection:1}

For every \(S\)-prescheme \(X\) and every integer \(n>0\), let
\[
  P^n:X\longrightarrow X
\]
denote the \(\mathbb F_p\)-endomorphism corresponding to raising
sections of \(\mathcal O_X\) to the \(p^n\)-th power, and let
\(X^{(n)}\) denote the \(S\)-prescheme obtained from \(X\) by base
change along \(P^n:S\to S\).  There is then a unique
\(S\)-morphism
\[
  F^n:X\longrightarrow X^{(n)}
\]
that makes the following diagram commutative:
\[
\begin{tikzcd}[column sep=3.7em,row sep=3.0em]
  {} & {} & X
    \arrow[dll,"P^n"']
    \arrow[dl,"F^n"]
    \arrow[ddl] \\
  X \arrow[d] &
  X^{(n)} \arrow[l] \arrow[d] &
  {} \\
  S &
  S \arrow[l,"P^n"] &
  {}
\end{tikzcd}
\]
% Loop-1 native reconstruction.  This diagram is included in the
% Session-B exhaustive no-overwrite Loop-2 high-zoom successor set.

It is clear that \(F^n\) identifies with the ``\(n\)-th iterate''
of \(F^1=F\), called the Frobenius morphism of \(X/S\).

If \(G\) is an \(S\)-group prescheme, then \(G^{(n)}\) is an
\(S\)-group prescheme and
\[
  F^n:G\longrightarrow G^{(n)}
\]
is an \(S\)-group morphism.  Its kernel
\({}_{F^n}(G)\) is a characteristic subgroup prescheme of \(G\)
(that is, it is stable under the functor
\(\operatorname{Aut}_{S\text{-}\mathrm{gr}}(G)\)), and is radicial
over \(S\).  If \(G\) is a radicial \(S\)-group prescheme, we say
that \(G\) has height at most \(h\) if
\({}_{F^h}(G)=G\).

\SGASubsection{B.2}{Groups and restricted
\texorpdfstring{\(p\)-Lie algebras}{p-Lie algebras}}
\SGARefTarget{sga3:xvii:appendix:B:subsection:2}

If \(G\) is an \(S\)-group prescheme, then
\(\operatorname{Lie}(G)\) (Exposé II) is naturally a restricted
\(p\)-Lie algebra (Exposé VII\textsubscript{A}, §§5 and 6).  In
particular, one has the following result (cf. Exposé
VII\textsubscript{A}).

\SGAStatement{\SGAInlineRefTarget{sga3:xvii:proposition:b-2-1}Proposition B.2.1}{%
\begin{enumerate}
\item[(i)] \(\operatorname{Lie}(\boldsymbol\alpha_p)_S
=\operatorname{Lie}(\mathbb G_a)_S\)\footnote{Editor's note:
  the definition of \(\boldsymbol\alpha_p\), already given at the
  beginning of this Exposé, has been omitted here.}
  is a sheaf of \(\mathcal O_S\)-modules, free of rank \(1\) over
  \(\mathcal O_S\), generated by an element \(X\) such that
  \(X^{(p)}=0\).
\item[(ii)] \(\operatorname{Lie}(\boldsymbol\mu_p)_S
  =\operatorname{Lie}(\mathbb G_m)_S\) is a sheaf of
  \(\mathcal O_S\)-modules, free of rank \(1\) over
  \(\mathcal O_S\), with a canonical basis \(X\) such that
  \(X^{(p)}=X\).
\end{enumerate}}
\SGARefTarget{sga3:xvii:proposition:B:2:1}

We now recall the fundamental result proved in Exposé
VII\textsubscript{A}, §7.

\SGAStatement{\SGAInlineRefTarget{sga3:xvii:theorem:b-2-2}Theorem B.2.2}{Suppose that \(S\) is affine, with
ring \(A\).  Then the functor
\[
  G\longmapsto\operatorname{Lie}G
\]
defines an equivalence between, on the one hand, the category of
\(S\)-group preschemes \(G\) that are of finite presentation and
flat over \(S\), of height \(1\), and whose Lie algebra is locally
free over \(\mathcal O_S\), and, on the other hand, the category
of restricted \(p\)-\(A\)-Lie algebras that are locally free of
finite rank.

Moreover, if \(G\) is as above and \(H\) is any \(S\)-group
prescheme of finite presentation, the canonical morphism
\[
  \operatorname{Hom}_{S\text{-}\mathrm{gr}}(G,H)
  \longrightarrow
  \operatorname{Hom}_{p\text{-}A\text{-}\mathrm{Lie}}
    \bigl(\operatorname{Lie}G(S),\operatorname{Lie}H(S)\bigr)
\]
is an isomorphism.}
\SGARefTarget{sga3:xvii:theorem:B:2:2}

\SGANextSectionTarget{sga3:xvii:subsection:b-3}\SGASubsection{B.3}{Radicial groups and smooth groups}
\SGARefTarget{sga3:xvii:appendix:B:subsection:3}

We now suppose that \(S\) is the spectrum of a field \(k\) of
characteristic \(p\).  Recall that Exposé
VI\textsubscript{A}, §5 proves that, if \(G\) is an algebraic
\(k\)-group and \(H\) an algebraic subgroup of \(G\), then the
sheaf \(G/H\) (for the fpqc topology) is representable.  We also
recall (VII\textsubscript{A}, \SGARefLink{sga3:VIIA:8.3}{8.3}):

\SGAStatement{\SGAInlineRefTarget{sga3:xvii:proposition:b-3-1}Proposition B.3.1}{Let \(G\) be an algebraic
\(k\)-group.  There exists an integer \(m\) such that, for every
\(n\geq m\), the algebraic group
\[
  G/{}_{F^n}(G)
\]
is smooth over \(k\).}
\SGARefTarget{sga3:xvii:proposition:B:3:1}

\SGAStatement{Proposition B.3.2}{Consider an exact sequence of
algebraic \(k\)-groups
\[
  1\longrightarrow G'\longrightarrow G
    \xrightarrow{u}G''\longrightarrow1
\]
and the following assertions:
\begin{enumerate}
\item[(i)] \SGARefTarget{sga3:xvii:proposition:B:3:2:item:i}
  the morphism \(u\) is smooth;
\item[(ii)] \SGARefTarget{sga3:xvii:proposition:B:3:2:item:ii}
  \(G'\) is smooth over \(k\);
\item[(iii)] \SGARefTarget{sga3:xvii:proposition:B:3:2:item:iii}
  for every integer \(n>0\), the sequence
  \[
    1\longrightarrow{}_{F^n}(G')
      \longrightarrow{}_{F^n}(G)
      \longrightarrow{}_{F^n}(G'')
      \longrightarrow1
  \]
  is exact;
\item[(iv)] \SGARefTarget{sga3:xvii:proposition:B:3:2:item:iv}
  the morphism \({}_FG\to{}_FG''\) is an epimorphism;
\item[(v)] \SGARefTarget{sga3:xvii:proposition:B:3:2:item:v}
  the morphism
  \[
    \bigl(\operatorname{Lie}G\bigr)(k)
      \longrightarrow\bigl(\operatorname{Lie}G''\bigr)(k)
  \]
  is surjective.
\end{enumerate}
Then
\[
  \text{\rm(i)}\Longleftrightarrow\text{\rm(ii)}
  \Longrightarrow\text{\rm(iii)}
  \Longrightarrow\text{\rm(iv)}
  \Longleftrightarrow\text{\rm(v)}.
\]
Moreover, if \(G\) is smooth, all five assertions are equivalent.}
\SGARefTarget{sga3:xvii:proposition:B:3:2}

The equivalence \SGARefLink{sga3:xvii:proposition:B:3:2:item:i}{(i)}\(\Longleftrightarrow\)\SGARefLink{sga3:xvii:proposition:B:3:2:item:ii}{(ii)} follows from
Exposé VI\textsubscript{B}, \SGARefLink{concl:sga3-VIB-9.2-vii}{9.2(vii)}.

For \SGARefLink{sga3:xvii:proposition:B:3:2:item:ii}{(ii)}\(\Longrightarrow\)\SGARefLink{sga3:xvii:proposition:B:3:2:item:iii}{(iii)}, if \(G'\) is smooth, then
\[
  F^n:G'\longrightarrow G'^{(n)}
\]
is an epimorphism, and \SGARefLink{sga3:xvii:proposition:B:3:2:item:iii}{(iii)} follows from the snake diagram.
The implication \SGARefLink{sga3:xvii:proposition:B:3:2:item:iii}{(iii)}\(\Longrightarrow\)\SGARefLink{sga3:xvii:proposition:B:3:2:item:iv}{(iv)} is clear, while
\SGARefLink{sga3:xvii:proposition:B:3:2:item:iv}{(iv)}\(\Longleftrightarrow\)\SGARefLink{sga3:xvii:proposition:B:3:2:item:v}{(v)} follows from \SGARefLink{sga3:xvii:theorem:b-2-2}{Theorem B.2.2}.
Finally, suppose that \(G\) is smooth and that \SGARefLink{sga3:xvii:proposition:B:3:2:item:v}{(v)} holds.  Then
\(G''\) is smooth, and \SGARefLink{sga3:xvii:proposition:B:3:2:item:v}{(v)} implies
\[
  \dim G'=\dim_k\bigl(\operatorname{Lie}G'\bigr)(k),
\]
so \(G'\) is smooth.

\SGASection{C}{Appendix III. Remarks and Supplements Concerning
Exposés XV, XVI, and XVII}
\SGARefTarget{sga3:xvii:appendix:C}

\SGANextSectionTarget{sga3:xvii:subsection:c-1}\SGASubsection{C.1}{}
\SGARefTarget{sga3:xvii:appendix:C:subsection:1}

It may be that \SGARefLink{sga3:xvii:proposition:1:2}{Propositions 1.2} and \SGARefLink{sga3:xv:proposition:1-2-bis}{1.2 bis} of Exposé XV remain
valid if the hypothesis that \(H_0\) is smooth over \(S_0\) is
removed.  This is the case, in particular, if \(G\) is finite,
flat, and commutative.

% END inputs/XVII/tex/components/36_expose_XVII_appendixB_pages045_046_and_appendixC_opening.tex

% BEGIN inputs/XVII/tex/components/37_expose_XVII_appendixC_sections2_through5_and_end_pages047_049.tex
\SGASubsection{C.2}{Complement to \SGARefLink{sga3:xv:remark:4:8}{Exposé XV, 4.8}}
\SGARefTarget{sga3:xvii:appendix:C:subsection:2}

The following proposition, as well as \SGARefLink{sga3:xvii:theorem:c-3-1}{Theorems C.3.1} and \SGARefLink{sga3:xvii:theorem:c-4-1}{C.4.1}
below, will appear in a forthcoming article by M. Raynaud on
group schemes over a discrete valuation ring.\footnote{Editor's
note: comments are to be added here, including references to
Exposé VI\textsubscript{B}, \ldots}

\SGAStatement{\SGAInlineRefTarget{sga3:xvii:proposition:c-2-1}Proposition C.2.1}{Let \(S\) be the spectrum of a
discrete valuation ring, let \(t\) be its generic point, let \(G\)
be a flat \(S\)-group prescheme of finite type, and put
\[
  \widetilde G=\operatorname{Spec}\Gamma(G,\mathcal O_G).
\]
Let \(u:G\to\widetilde G\) be the canonical morphism.  Then:
\begin{enumerate}
\item[(1)] \SGARefTarget{sga3:xvii:proposition:C:2:1:item:1}
  \(\widetilde G\) is naturally an \(S\)-group scheme, and \(u\)
  is a homomorphism.
\item[(2)] \SGARefTarget{sga3:xvii:proposition:C:2:1:item:2}
  \(\ker(u)\) is flat over \(S\), and
  \((\widetilde G,u)\) is the fpqc quotient of \(G\) by
  \(\ker(u)\).  Thus \(\widetilde G\) is the largest affine
  quotient of \(G\).
\item[(3)] \SGARefTarget{sga3:xvii:proposition:C:2:1:item:3}
  If \(G_t\) is affine, then \(\ker(u)\) is an étale group over
  \(S\), and it is the trivial group if and only if \(G\) is
  separated over \(S\).  In particular, a flat, separated
  \(S\)-group scheme of finite type whose generic fiber is affine
  is affine.
\end{enumerate}}
\SGARefTarget{sga3:xvii:proposition:C:2:1}

\SGANextSectionTarget{sga3:xvii:subsection:c-3}\SGASubsection{C.3}{}
\SGARefTarget{sga3:xvii:appendix:C:subsection:3}

In the statement of \SGARefLink{sga3:xv:lemma:6:6}{Exposé XV, 6.6}, the hypothesis that \(H\)
be equal to its connected normalizer is unnecessary in order that the
set of points \(s\) of \(S\) for which \(H_s\) is a parabolic
subgroup of \(G_s\) be open.  Indeed, repeat the proof given in
paragraph~c) preceding \SGARefLink{sga3:xv:lemma:6:9}{Lemma 6.9}, and let \(N\) denote the
scheme-theoretic closure of \(N_t\) in \(G\).  Then \(N\) is a
closed subgroup scheme of \(G\), flat over \(S\), which contains
\(H\) and is contained in \(\mathcal N\).  The theorem below
shows that \(G/N\) is representable.  Since \(G_s/N_s\) is proper
and connected, the proof concludes as in \emph{loc.\ cit.}

\SGAStatement{\SGAInlineRefTarget{sga3:xvii:theorem:c-3-1}Theorem C.3.1\footnote{Editor's note: see \SGARefLink{sga3:xvii:appendix:C:subsection:4}{Theorem
4}C in S. Anantharaman, ``Schémas en groupes, espaces homogènes et
espaces algébriques sur une base de dimension 1,'' \emph{Bull.
Soc. Math. France, Mém.} 33 (1976), 5--79.}}{Let \(S\) be the
spectrum of a henselian discrete valuation ring, let \(G\) be an
\(S\)-group prescheme locally of finite type, and let \(H\) be a
closed subgroup prescheme of \(G\), flat over \(S\).  Then \(G/H\)
is representable.}
\SGARefTarget{sga3:xvii:theorem:C:3:1}

\SGANextSectionTarget{sga3:xvii:subsection:c-4}\SGASubsection{C.4}{Complement to \SGARefLink{sga3:xvi:subsection:1:1}{Exposé XVI, 1.1}}
\SGARefTarget{sga3:xvii:appendix:C:subsection:4}

The following theorem refines \SGARefLink{sga3:viii:proposition:7:9}{Exposé VIII, 7.9}.

\SGAStatement{\SGAInlineRefTarget{sga3:xvii:theorem:c-4-1}Theorem C.4.1}{Let \(S\) be the spectrum of a
discrete valuation ring with residue characteristic \(p>0\), and
let \(G\) be a commutative, smooth, separated \(S\)-group scheme
of finite type.  Then the following conditions are equivalent:
\begin{enumerate}
\item[(i)] \SGARefTarget{sga3:xvii:theorem:C:4:1:item:i}
  \({}_pG\) is finite over \(S\).
\item[(ii)] \SGARefTarget{sga3:xvii:theorem:C:4:1:item:ii}
  For every \(S\)-prescheme \(S'\), every separated
  \(S'\)-group prescheme \(H'\) of finite presentation over
  \(S'\), and every \(S'\)-monomorphism
  \[
    u:G_{S'}\longrightarrow H',
  \]
  the morphism \(u\) is an immersion.
\end{enumerate}}
\SGARefTarget{sga3:xvii:theorem:C:4:1}

\SGANextSectionTarget{sga3:xvii:subsection:c-5}\SGASubsection{C.5}{}
\SGARefTarget{sga3:xvii:appendix:C:subsection:5}

Example XVII, \SGARefLink{sga3:xvii:subsection:6:4:item:a}{6.4(a)}, gives a smooth group over a field \(k\)
whose unipotent radical is not defined over \(k\).  The following
proposition gives a general method for obtaining such groups.

\SGAStatement{\SGAInlineRefTarget{sga3:xvii:proposition:c-5-1}Proposition C.5.1}{Let \(k\) be a field, let \(K\)
be a finite radicial extension of \(k\) of degree greater than
\(1\), and let \(H\) be a smooth, connected \(K\)-algebraic group
of dimension \(r\).  Put
\[
  G=\prod_{K/k}H/K,
\]
which is a smooth, connected \(k\)-algebraic group of dimension
\(r[K:k]\).  Let
\[
  u:G_K\longrightarrow H
\]
be the canonical homomorphism, and put \(R=\ker(u)\).  Then:
\begin{enumerate}
\item[(i)] \SGARefTarget{sga3:xvii:proposition:C:5:1:item:i}
  \(u\) is an epimorphism, and \(R\) is a smooth, connected,
  unipotent algebraic group.
\item[(ii)] \SGARefTarget{sga3:xvii:proposition:C:5:1:item:ii}
  If \(U\) is a smooth algebraic subgroup of \(G\) such that
  \(U_K\) contains \(R\), then \(U=G\).
\end{enumerate}}
\SGARefTarget{sga3:xvii:proposition:C:5:1}

\SGAStatement{Corollary}{Retain the notation of the preceding
proposition.
\begin{enumerate}
\item[(a)] \SGARefTarget{sga3:xvii:corollary:C:5:1:item:a}
  If \(H\) is not unipotent, then the unipotent radical of
  \(G_{\bar k}\) is not defined over \(k\).
\item[(b)] \SGARefTarget{sga3:xvii:corollary:C:5:1:item:b}
  If \(H\) is a nonzero abelian variety, then \(G\) is not an
  extension of an abelian variety by a smooth linear group.
\item[(c)] \SGARefTarget{sga3:xvii:corollary:C:5:1:item:c}
  If \(H\) is not solvable, then the solvable radical of
  \(G_{\bar k}\) is not defined over \(k\), and \(G\) has no
  Borel subgroup defined over \(k\).
\end{enumerate}}
\SGARefTarget{sga3:xvii:corollary:C:5:1}

Let us first show how the corollary follows from the proposition.

\smallskip
\noindent\textit{\SGARefLink{sga3:xvii:appendix:A}{(a)}} Let \(U\) be a unipotent radical of \(G\).
Then \(U_K\) is the unipotent radical of \(G_K\), and hence
contains \(R\), since \(R\) is smooth, connected, and unipotent
by \SGARefLink{sga3:xvii:corollary:C:5:1:item:i}{(i)}.  Thus \(U=G\) by \SGARefLink{sga3:xvii:corollary:C:5:1:item:ii}{(ii)}.  But \(G_K\) has \(H\) as a
quotient, and \(H\) is not unipotent by hypothesis.  Therefore
\(G\) is not unipotent, a contradiction.

\smallskip
\noindent\textit{\SGARefLink{sga3:xvii:section:b}{(b)}} Suppose that \(G\) is an extension of an
abelian variety by a linear group \(L\) that is smooth over
\(k\).  Then \(L_K\) necessarily contains \(R\), so \(L=G\).
But \(G\) cannot be a linear group, because \(G_K\) has a nonzero
abelian variety as a quotient.

\smallskip
\noindent\textit{(c)} Let \(S\) be a solvable radical of \(G\),
and let \(B\) be a Borel subgroup of \(G\).  Then \(S_K\) and
\(B_K\) contain \(R\), so \(S=G\) and \(B=G\).  But then \(G\)
is solvable, contradicting the fact that \(G_K\) has a quotient
that is not solvable.

\SGAProof{Proof of \SGARefLink{sga3:xvii:proposition:c-5-1}{Proposition C.5.1}}
\textit{\SGAInlineRefTarget{sga3:xvii:corollary:C:5:1:item:i}(i)} We begin by describing the morphism \(u\).  A
\(K\)-scheme
\[
  f:S\longrightarrow\operatorname{Spec}(K)
\]
gives rise to the diagram
\[
\begin{tikzcd}[column sep=6.0em,row sep=4.0em]
  \operatorname{Spec}(K)
    \arrow[d,"j"'] &
  T
    \arrow[l,"h"']
    \arrow[d,bend right=28,"j_T"'] \\
  \operatorname{Spec}(k) &
  S
    \arrow[l,"g"]
    \arrow[ul,"f"]
    \arrow[u,bend right=28,"s"']
\end{tikzcd}
\]
where
\[
  g=j\circ f,\qquad
  T=\operatorname{Spec}(K)\times_{\operatorname{Spec}(k)}S,
\]
\(h\) and \(j_T\) are the two projections, and \(s\) is the
section of \(T\) over \(S\) such that \(h\circ s=f\).  The map
\[
  u(S):G_K(S)\longrightarrow H(S)
\]
is simply the composite
\[
  G_K(S)\xrightarrow{\sim}H(T)\longrightarrow H(S),
\]
where the last arrow is defined by the section \(s\).

Take, in particular, \(S\) to be the spectrum of an algebraic
closure \(\bar k\) of \(k\), and let \(f\) be the unique
\(k\)-morphism induced by \(K\to\bar k\), so that \(T\) is a
local Artinian scheme.  To prove \SGARefLink{sga3:xvii:corollary:C:5:1:item:i}{(i)}, it is enough to do so after
extending the base field from \(K\) to \(\bar k\).  It is clear
that
\[
  G_S=\prod_{T/S}H_T/T
\]
represents the Greenberg functor of \(H_T\) relative to \(S\)
(M. J. Greenberg, ``Schemata over local rings,'' \emph{Ann. of
Math.} 73 (1961), 624--648).  Under this identification, the
description above shows that \(u_S\) is the canonical transition
morphism
\[
  \operatorname{Green}(H_T)\longrightarrow
  H_S=H_T\times_T S.
\]
Assertion \SGARefLink{sga3:xvii:corollary:C:5:1:item:i}{(i)} now follows from the smoothness of \(H\) over \(K\)
and M. J. Greenberg, ``Schemata over local rings II,''
\emph{Ann. of Math.} 78 (1963), 256--266.

\smallskip
\noindent\textit{\SGAInlineRefTarget{sga3:xvii:corollary:C:5:1:item:ii}(ii)} To prove \SGARefLink{sga3:xvii:corollary:C:5:1:item:ii}{(ii)}, we may suppose that \(k\)
is separably closed.  Let \(U\) be a smooth algebraic subgroup of
\(G\) such that \(U_K\supset R\), and let us prove that \(U=G\).
Since \(G\) is connected, we may suppose that \(U\) is connected.
Let \(V=u(U_K)\), a smooth, connected algebraic subgroup of
\(H\), and put
\[
  V'=\prod_{K/k}V/K,
\]
a smooth, connected algebraic subgroup of \(G\).  The group
\(V'_K\) is an algebraic subgroup of \(G_K\), and the canonical
morphism \(V'_K\to V\) is the restriction of \(u\) to \(V'_K\).
By hypothesis, \(U_K\) contains \(R\), hence \emph{a fortiori}
contains \(\ker(V'_K\to V)\); by construction, moreover, the
image of \(U_K\) in \(H\) is \(V\).  It follows that \(U_K\)
contains \(V'_K\), and hence that \(U\) contains \(V'\).

On the other hand, the canonical isomorphism
\[
  G(k)\lhook\joinrel\longrightarrow G(K)
    \xrightarrow{u(K)}H(K)
\]
plainly sends \(U(k)\) into \(V(K)=V'(k)\).  Thus
\(U(k)\subset V'(k)\).  Since \(U\) is smooth and \(k\) is
separably closed, it follows that \(U\subset V'\), whence
\(U=V'\).  We then have
\[
  ([K:k]-1)\dim V
  =\dim\ker(V'_K\to V)
  =\dim R
  =([K:k]-1)\dim H.
\]
Since \(K\ne k\), we conclude that \(\dim V=\dim H\), hence that
\(V=H\), and finally that \(U=V'=G\).

% END inputs/XVII/tex/components/37_expose_XVII_appendixC_sections2_through5_and_end_pages047_049.tex


